% L'œil réduit : fonctionnement, défauts et correction — Physique 1ère C — Cours signé OnBuch+
% Programme officiel MINESEC. Compiler avec : tectonic N08-oeil-reduit-defauts-correction.tex
\newcommand{\DOCMATIERE}{Physique}
\newcommand{\DOCNIVEAU}{1ère C}
\newcommand{\DOCMODULE}{Module 3 : Optique}
\newcommand{\DOCLECON}{N08}
\newcommand{\DOCTITRE}{L'œil réduit : fonctionnement, défauts et correction}
% OnBuch+ — préambule « cours signés » ENRICHI (compilé avec tectonic / XeLaTeX)
% Système visuel « Pop » d'OnBuch+ : fond crème (jamais blanc), bordures encre
% épaisses, ombres « dures » décalées, titres Archivo Black en capitales.
% Version MATHÉMATIQUES : graphes (pgfplots/tikz), boîtes pédagogiques
% (définition, propriété, méthode, attention, le savais-tu, exercice résolu,
% exercice, TP, bilan), page de garde et signature OnBuch+ ancrée en bas.
%
% Macros attendues dans main.tex AVANT \input{preamble} :
%   \DOCMATIERE  \DOCNIVEAU  \DOCMODULE  \DOCLECON (numéro)  \DOCTITRE
\documentclass[11pt]{article}

\usepackage{fontspec}
\usepackage[french]{babel}
\usepackage[a4paper,margin=2cm,top=3.4cm,bottom=2.8cm,headheight=1.4cm,footskip=1.3cm]{geometry}
\usepackage{amsmath,amssymb,amsfonts}
\usepackage{mathtools} % \xmapsto, \coloneqq... souvent utilisés par les modèles
\usepackage{cancel} % simplifications barrées (bilans de forces)
% \abs{z} (module, valeur absolue) et \norm{u} (norme), utilisés par les modèles
\providecommand{\abs}{}\let\abs\relax\DeclarePairedDelimiter{\abs}{\lvert}{\rvert}
\providecommand{\norm}{}\let\norm\relax\DeclarePairedDelimiter{\norm}{\lVert}{\rVert}
\usepackage[table]{xcolor} % [table] : fournit aussi \rowcolors (alternance de lignes)
\usepackage{pagecolor}
\usepackage{tikz}
\usetikzlibrary{arrows.meta,positioning,calc,shapes.geometric,shapes.misc,shapes.arrows,decorations.pathmorphing,decorations.pathreplacing,decorations.markings,patterns,shadows,mindmap,trees,fit,backgrounds,matrix,intersections,angles,quotes}
\usepackage{pgfplots}
\usepackage[version=4]{mhchem} % équations nucléaires \ce{^{235}_{92}U}
\usepackage[european,siunitx]{circuitikz} % schémas électriques
\pgfplotsset{compat=1.18}
\usepgfplotslibrary{fillbetween,groupplots} % aires entre courbes (calcul intégral)
\usepackage{enumitem}
\usepackage{fancyhdr}
% [most] charge toutes les bibliothèques tcolorbox (listings, minted, raster,
% documentation...) et gonfle inutilement la mémoire TeX sur un document long
% (~300 boîtes) : on ne charge que ce qui est réellement utilisé.
\usepackage[skins,breakable]{tcolorbox}
\usepackage{graphicx}
\usepackage{colortbl}
\usepackage{booktabs}
\usepackage{tabularx}
\usepackage{diagbox} % case d'en-tête diagonale des tableaux à double entrée
% Colonnes extensibles centrée / gauche / droite (C, L, R), souvent utilisées par les modèles
\newcolumntype{C}{>{\centering\arraybackslash}X}
\newcolumntype{L}{>{\raggedright\arraybackslash}X}
\newcolumntype{R}{>{\raggedleft\arraybackslash}X}
\usepackage{array}
\usepackage{multicol}
\usepackage{siunitx}
% parse-numbers=false : accepte aussi \SI{2,5 \times 10^{-3}}{...}
\sisetup{locale=FR,output-decimal-marker={,},group-minimum-digits=4,parse-numbers=false}
\DeclareSIUnit\ohm{\ensuremath{\symup{\Omega}}} % U+2126 absent de la police
\DeclareSIUnit\division{div} % réglages d'oscilloscope (ms/div, V/div)
\DeclareSIUnit\tr{tr} % tours (vitesse de rotation, ex. \SI{1500}{\tr\per\minute})
\DeclareSIUnit\ch{ch} % cheval-vapeur (puissance moteur, unité usuelle hors SI)
\DeclareSIUnit\franccfa{FCFA} % franc CFA (coûts, contexte camerounais)
\DeclareSIUnit\dioptre{\ensuremath{\delta}} % vergence d'une lentille (dioptrie)
\usepackage{qrcode}
% \degree : commande fréquemment utilisée par les modèles pour un angle ou une
% température en dehors de \SI{}{} (non fournie par les paquets chargés).
\providecommand{\degree}{^{\circ}}
% \qed (fin de démonstration), utilisé par les modèles sans amsthm
\providecommand{\qed}{\hfill$\square$}
% \barre : double barre des valeurs interdites dans un tableau de variations (tkz-tab)
\providecommand{\barre}{\ensuremath{\|}}
% environnement proof (amsthm non chargé) utilisé par les modèles
\ifdefined\proof\else\newenvironment{proof}[1][Démonstration]{\par\noindent\textit{#1.}\ }{\qed\par}\fi
\usepackage{lastpage}
\usepackage{hyperref}
\hypersetup{hidelinks}

% ============================================================
% Palette « Pop » OnBuch+ — mêmes hex que la classe P (plus_ui.dart)
% ============================================================
\definecolor{poporange}{HTML}{F59321}
\definecolor{popdark}{HTML}{E07A0C}
\definecolor{popcream}{HTML}{FFF6E8}
\definecolor{popink}{HTML}{1C1714}
\definecolor{poppurple}{HTML}{7A5AE0}
\definecolor{popgreen}{HTML}{2E9E6A}
\definecolor{poppink}{HTML}{E0457B}
\definecolor{popblue}{HTML}{2F7DE1}
\definecolor{popgold}{HTML}{C98A12}
\definecolor{popmuted}{HTML}{8A7F76}
\definecolor{popline}{HTML}{EADFD2}
\definecolor{popgray}{HTML}{8A7F76}
\colorlet{popgrey}{popgray}
% Alias tolérants (le modèle nomme parfois les couleurs différemment)
\colorlet{popwhite}{white}
\colorlet{popblack}{popink}
\colorlet{popred}{poppink}
\colorlet{popyellow}{popgold}
% Teintes claires pour fonds de tableaux / zones
\colorlet{popblueL}{popblue!10!white}
\colorlet{poporangeL}{poporange!14!white}
\colorlet{popgreenL}{popgreen!12!white}
\colorlet{poppurpleL}{poppurple!12!white}
\colorlet{poppinkL}{poppink!12!white}
\colorlet{popgoldL}{popgold!14!white}

\pagecolor{popcream}

% ============================================================
% Typographie
% ============================================================
\defaultfontfeatures{Ligatures=TeX}
\setmainfont{PlusJakartaSans-Medium.ttf}[Path=../../../fonts/,BoldFont=PlusJakartaSans-Bold.ttf]
% Maths en sans-serif (Fira Math) avec chiffres et lettres droites de Plus
% Jakarta, pour que les formules se fondent dans le texte.
\usepackage[math-style=ISO,bold-style=ISO]{unicode-math}
\setmathfont{FiraMath-Regular.otf}[Path=../../../fonts/]
\setmathfont{PlusJakartaSans-Medium.ttf}[Path=../../../fonts/,range={up/num,up/{Latin,latin},\mathup}]
\usepackage{icomma} % virgule décimale sans espace en mode math
% Caractères mathématiques Unicode absents des polices de texte (Plus Jakarta,
% Archivo Black) : rendus par leur équivalent mathématique, en texte comme en
% formule — sinon ils s'affichent en carré vide (ex. « Arithmétique dans ℤ »).
\usepackage{newunicodechar}
\newunicodechar{ℝ}{\ensuremath{\mathbb{R}}}
\newunicodechar{ℤ}{\ensuremath{\mathbb{Z}}}
\newunicodechar{ℂ}{\ensuremath{\mathbb{C}}}
\newunicodechar{ℕ}{\ensuremath{\mathbb{N}}}
\newunicodechar{ℚ}{\ensuremath{\mathbb{Q}}}
\newunicodechar{𝔻}{\ensuremath{\mathbb{D}}}
\newunicodechar{α}{\ensuremath{\alpha}}
\newunicodechar{β}{\ensuremath{\beta}}
\newunicodechar{γ}{\ensuremath{\gamma}}
\newunicodechar{δ}{\ensuremath{\delta}}
\newunicodechar{ε}{\ensuremath{\varepsilon}}
\newunicodechar{θ}{\ensuremath{\theta}}
\newunicodechar{λ}{\ensuremath{\lambda}}
\newunicodechar{μ}{\ensuremath{\mu}}
\newunicodechar{σ}{\ensuremath{\sigma}}
\newunicodechar{τ}{\ensuremath{\tau}}
\newunicodechar{φ}{\ensuremath{\varphi}}
\newunicodechar{ψ}{\ensuremath{\psi}}
\newunicodechar{ω}{\ensuremath{\omega}}
\newunicodechar{Σ}{\ensuremath{\Sigma}}
% majuscules produites par \MakeUppercase dans les titres (α-aminés -> Α)
\newunicodechar{Α}{\ensuremath{\alpha}}
\newunicodechar{Β}{\ensuremath{\beta}}
\newunicodechar{Γ}{\ensuremath{\Gamma}}
\newunicodechar{Δ}{\ensuremath{\Delta}}
\newunicodechar{Ε}{\ensuremath{\varepsilon}}
\newunicodechar{Θ}{\ensuremath{\Theta}}
\newunicodechar{Λ}{\ensuremath{\Lambda}}
\newunicodechar{Μ}{\ensuremath{\mu}}
\newunicodechar{Τ}{\ensuremath{\tau}}
\newunicodechar{Φ}{\ensuremath{\Phi}}
\newunicodechar{Ψ}{\ensuremath{\Psi}}
\newunicodechar{Ω}{\ensuremath{\Omega}}
\newunicodechar{↦}{\ensuremath{\mapsto}}
\newunicodechar{∈}{\ensuremath{\in}}
\newunicodechar{∉}{\ensuremath{\notin}}
\newunicodechar{∧}{\ensuremath{\wedge}}
\newunicodechar{⇔}{\ensuremath{\Leftrightarrow}}
\newunicodechar{⟺}{\ensuremath{\Longleftrightarrow}}
\newunicodechar{⇒}{\ensuremath{\Rightarrow}}
\newunicodechar{⟹}{\ensuremath{\Longrightarrow}}
\newunicodechar{∘}{\ensuremath{\circ}}
\newunicodechar{∪}{\ensuremath{\cup}}
\newunicodechar{∩}{\ensuremath{\cap}}
\newunicodechar{≡}{\ensuremath{\equiv}}
\newunicodechar{⊂}{\ensuremath{\subset}}
\newunicodechar{⊥}{\ensuremath{\perp}}
\newunicodechar{∃}{\ensuremath{\exists}}
\newunicodechar{∀}{\ensuremath{\forall}}
\newunicodechar{∛}{\ensuremath{\sqrt[3]{\,}}}
\newunicodechar{∜}{\ensuremath{\sqrt[4]{\,}}}
\newunicodechar{⃗}{\ensuremath{{}^{\to}}}
\newunicodechar{ⁿ}{\ifmmode^{n}\else\textsuperscript{\ensuremath{n}}\fi}
\newunicodechar{ˣ}{\ifmmode^{x}\else\textsuperscript{\ensuremath{x}}\fi}
\newunicodechar{ᵇ}{\ifmmode^{b}\else\textsuperscript{\ensuremath{b}}\fi}
\newunicodechar{ʳ}{\ifmmode^{r}\else\textsuperscript{\ensuremath{r}}\fi}
\newunicodechar{ᵘ}{\ifmmode^{u}\else\textsuperscript{\ensuremath{u}}\fi}
\newunicodechar{ʸ}{\ifmmode^{y}\else\textsuperscript{\ensuremath{y}}\fi}
\newunicodechar{ᵃ}{\ifmmode^{a}\else\textsuperscript{\ensuremath{a}}\fi}
\newunicodechar{ᵉ}{\ifmmode^{e}\else\textsuperscript{\ensuremath{e}}\fi}
\newunicodechar{ᵏ}{\ifmmode^{k}\else\textsuperscript{\ensuremath{k}}\fi}
\newunicodechar{ᵖ}{\ifmmode^{p}\else\textsuperscript{\ensuremath{p}}\fi}
\newunicodechar{ᴺ}{\ifmmode^{N}\else\textsuperscript{\ensuremath{N}}\fi}
\newunicodechar{ᶜ}{\ifmmode^{c}\else\textsuperscript{\ensuremath{c}}\fi}
\newunicodechar{ʰ}{\ifmmode^{h}\else\textsuperscript{\ensuremath{h}}\fi}
\newunicodechar{ᵠ}{\ifmmode^{q}\else\textsuperscript{\ensuremath{q}}\fi}
\newunicodechar{ₙ}{\ifmmode_{n}\else\textsubscript{\ensuremath{n}}\fi}
\newunicodechar{ₐ}{\ifmmode_{a}\else\textsubscript{\ensuremath{a}}\fi}
\newunicodechar{ᵢ}{\ifmmode_{i}\else\textsubscript{\ensuremath{i}}\fi}
\newunicodechar{ᵣ}{\ifmmode_{r}\else\textsubscript{\ensuremath{r}}\fi}
\newunicodechar{ₖ}{\ifmmode_{k}\else\textsubscript{\ensuremath{k}}\fi}
\newunicodechar{ᵦ}{\ifmmode_{\beta}\else\textsubscript{\ensuremath{\beta}}\fi}
\newunicodechar{ₓ}{\ifmmode_{x}\else\textsubscript{\ensuremath{x}}\fi}
\newfontfamily\popdisplay{ArchivoBlack-Regular.ttf}[Path=../../../fonts/]
\newfontfamily\poplogo{SpaceGrotesk-Bold.ttf}[Path=../../../fonts/]
\newfontfamily\popsemi{PlusJakartaSans-SemiBold.ttf}[Path=../../../fonts/]
\newfontfamily\popxbold{PlusJakartaSans-ExtraBold.ttf}[Path=../../../fonts/]
\newfontfamily\popxboldls{PlusJakartaSans-ExtraBold.ttf}[Path=../../../fonts/,LetterSpace=3.0]

\setlist[enumerate]{leftmargin=1.7em,itemsep=5pt,topsep=4pt}
\setlist[itemize]{leftmargin=1.7em,itemsep=4pt,topsep=4pt,label={\color{poporange}\textbullet}}
\setlength{\parindent}{0pt}
\setlength{\parskip}{5pt}
\renewcommand{\arraystretch}{1.25}

% pgfplots : style Pop par défaut
\pgfplotsset{
  every axis/.append style={axis line style={popink,line width=1pt},tick style={popink},
    grid style={popline,line width=0.5pt},label style={font=\small},tick label style={font=\footnotesize}},
  popcurve/.style={poporange,line width=1.8pt,no markers,smooth},
}
% Même style disponible dans un \draw TikZ simple (hors pgfplots)
\tikzset{popcurve/.style={poporange,line width=1.8pt}}
\tikzset{popgrid/.style={popline,very thin}}
\colorlet{popcurve}{poporange} % le nom du style est parfois utilisé comme couleur

% ============================================================
% Pill / squiggle
% ============================================================
\newcommand{\poppill}[3]{%
  \tikz[baseline=-0.55ex]\node[draw=popink,line width=1.2pt,rounded corners=2.6pt,%
    fill=#2,inner xsep=6.5pt,inner ysep=3pt]{%
    \popxboldls\fontsize{7}{7}\selectfont\color{#3}\MakeUppercase{#1}};%
}
\newcommand{\popsquiggle}[1]{%
  \tikz[baseline=-0.4ex]{
    \draw[#1,line width=2.6pt,line cap=round]
      plot [smooth] coordinates {(0,0.09) (0.34,0.01) (0.68,0.21) (1.02,0.01) (1.36,0.10)};
  }%
}
% Mot-clé surligné (marqueur orange sous le texte)
\newcommand{\cle}[1]{{\popxbold\color{popdark}#1}}

% ============================================================
% En-tête / pied
% ============================================================
\pagestyle{fancy}
\fancyhf{}
\renewcommand{\headrulewidth}{0pt}
\renewcommand{\footrulewidth}{0pt}
\lhead{\raisebox{-0.15ex}{\poplogo\fontsize{13}{13}\selectfont\color{popink}OnBuch\color{poporange}+}}
\rhead{\poppill{\DOCMATIERE\ \textbullet\ \DOCNIVEAU\ \textbullet\ Leçon \DOCLECON}{white}{popink}}
\lfoot{\popxbold\fontsize{7.6}{7.6}\selectfont\color{popmuted}\MakeUppercase{Cours signé OnBuch+}\ \textbullet\ {\popsemi\DOCTITRE}}
\rfoot{\tikz[baseline=-0.6ex]\node[rounded corners=8.5pt,fill=popink,minimum height=17pt,minimum width=17pt,inner xsep=5pt,inner sep=0]{\popxbold\fontsize{8}{8}\selectfont\color{popcream}\thepage\,/\,\pageref*{LastPage}};}

% ============================================================
% Titres
% ============================================================
\newcommand{\chaptitre}[2]{%
  \begin{tcolorbox}[enhanced,colback=poporange,colframe=popink,boxrule=2.6pt,arc=11pt,
      drop shadow=popink,left=15pt,right=15pt,top=12pt,bottom=11pt,before skip=3pt,after skip=18pt]
    \poppill{#2}{popink}{popcream}\par\vspace{9pt}
    {\raggedright\hyphenpenalty=10000\exhyphenpenalty=10000\popdisplay\fontsize{22}{24}\selectfont\color{popcream}\MakeUppercase{#1}\par}
  \end{tcolorbox}
}
% Section numérotée : pastille orange avec le numéro + titre Archivo
\newcounter{popsec}
\newcommand{\coursec}[1]{%
  \stepcounter{popsec}\setcounter{popsub}{0}%
  \par\vspace{16pt}\noindent
  \begin{minipage}{\linewidth}
  \tikz[baseline=-0.7ex]\node[draw=popink,line width=1.6pt,fill=poporange,rounded corners=4pt,minimum size=22pt,inner sep=0]{\popdisplay\fontsize{12}{12}\selectfont\color{popcream}\thepopsec};%
  \hspace{8pt}{\popdisplay\fontsize{15}{17}\selectfont\color{popink}\MakeUppercase{#1}}\par
  \vspace{4pt}\noindent{\color{poporange}\rule{36pt}{3.6pt}}
  \end{minipage}\par\nopagebreak\vspace{8pt}
}
\newcounter{popsub}
\newcommand{\courssub}[1]{%
  \stepcounter{popsub}%
  \par\vspace{9pt}\noindent{\popxbold\fontsize{11.5}{13}\selectfont\color{popink}{\color{poporange}\thepopsec.\thepopsub}\hspace{6pt}#1}\par\nopagebreak\vspace{4pt}
}

% ============================================================
% Boîtes pédagogiques — même recette : fond blanc, bordure épaisse colorée,
% ombre dure encre, badge plein. Titre optionnel [#1].
% ============================================================
\tcbset{popbase/.style={breakable,enhanced,colback=white,boxrule=2.3pt,arc=9pt,
  shadow={3pt}{-3pt}{0pt}{popink},left=14pt,right=14pt,top=11pt,bottom=11pt,before skip=11pt,after skip=15pt,
  fontupper=\popsemi\fontsize{10.6}{14.5}\selectfont\color{popink},
  fontlower=\popsemi\fontsize{10.6}{14.5}\selectfont\color{popink}}}
% Coche dessinée (le glyphe ✓ n'existe pas dans les polices du document)
\newcommand{\popcheck}[1]{\tikz[baseline=-0.5ex]\draw[#1,line width=1.6pt,line cap=round,line join=round](-0.13,0.0)--(-0.04,-0.09)--(0.13,0.1);}
\providecommand{\coche}{\popcheck{popgreen}}
\providecommand{\resultat}[1]{\par\smallskip\noindent\fcolorbox{popgreen}{popgreenL}{\parbox{\dimexpr\linewidth-2\fboxsep-2\fboxrule\relax}{\textbf{Résultat.} #1}}\par\smallskip}
\newcommand{\popboxhead}[3]{\poppill{#1}{#2}{white}\ifx\relax#3\relax\else\hspace{6pt}{\popxbold\fontsize{10.6}{12}\selectfont\color{popink}#3}\fi\par\vspace{7pt}}

\newtcolorbox{aretenir}[1][]{popbase,colframe=popgreen,before upper={\popboxhead{À retenir}{popgreen}{#1}}}
\newtcolorbox{exemplebox}[1][]{popbase,colframe=poporange,before upper={\popboxhead{Exemple}{poporange}{#1}}}
\newtcolorbox{definition}[1][]{popbase,colframe=popblue,before upper={\popboxhead{Définition}{popblue}{#1}}}
\newtcolorbox{propriete}[1][]{popbase,colframe=poppurple,before upper={\popboxhead{Propriété}{poppurple}{#1}}}
\newtcolorbox{methode}[1][]{popbase,colframe=popgold,before upper={\popboxhead{Méthode}{popgold}{#1}}}
\newtcolorbox{attention}[1][]{popbase,colframe=poppink,before upper={\popboxhead{Attention}{poppink}{#1}}}
\newtcolorbox{experience}[1][]{popbase,colframe=popblue,colback=popblueL,before upper={\popboxhead{Expérience}{popblue}{#1}}}
% « Le savais-tu ? » : carte violette pleine (contexte, culture, Cameroun)
\newtcolorbox{savaistu}[1][]{popbase,colback=poppurple,colframe=popink,
  fontupper=\popsemi\fontsize{10.4}{14.2}\selectfont\color{white},
  before upper={\poppill{Le savais-tu\,?}{white}{poppurple}\ifx\relax#1\relax\else\hspace{6pt}{\popxbold\color{white}#1}\fi\par\vspace{7pt}}}
% Exercice résolu : énoncé en haut, \tcblower puis la solution sur fond crème
\newcounter{popexo}\newcounter{popexr}
\newtcolorbox{exoresolu}[1][]{popbase,colframe=poporange,colbacklower=popcream,
  segmentation style={popink,line width=1.2pt,dashed},
  before upper={\stepcounter{popexr}\popboxhead{Exercice résolu \thepopexr}{poporange}{#1}},
  before lower={\poppill{Solution}{popink}{popcream}\par\vspace{6pt}}}
% Exercice (non résolu en place) — la correction est dans la partie Corrigés
\newtcolorbox{exercice}[1][]{popbase,colframe=popink,
  before upper={\stepcounter{popexo}\popboxhead{Exercice \thepopexo}{popink}{#1}}}
% Corrigé
\newtcolorbox{corrige}[1][]{popbase,colframe=popgreen,colback=popgreenL,
  before upper={\popboxhead{Corrigé}{popgreen}{#1}}}
% Objectifs / prérequis / bilan
\newtcolorbox{objectifs}{popbase,colframe=popink,colback=poporangeL,
  before upper={\popboxhead{Objectifs de la leçon}{popink}{}}}
\newtcolorbox{prerequis}{popbase,colframe=popink,colback=popblueL,
  before upper={\popboxhead{Prérequis}{popblue}{}}}
\newtcolorbox{bilan}{popbase,colframe=popink,colback=white,boxrule=2.8pt,
  before upper={\popboxhead{Fiche bilan}{poporange}{L'essentiel à connaître pour l'examen}}}
% Figure encadrée avec légende
\newtcolorbox{popfigure}[1][]{enhanced,breakable=false,colback=white,colframe=popline,boxrule=1.4pt,arc=8pt,
  left=10pt,right=10pt,top=10pt,bottom=8pt,before skip=10pt,after skip=12pt,halign=center,
  lower separated=false,halign lower=center,
  fontlower=\popsemi\fontsize{9}{11.5}\selectfont\color{popmuted},
  #1}
\newcommand{\legende}[1]{\par\vspace{6pt}{\popsemi\fontsize{9}{11.5}\selectfont\color{popmuted}#1}}

% ============================================================
% Signature OnBuch+
% ============================================================
\newcommand{\poplogoinline}[1]{{\poplogo\fontsize{#1}{#1}\selectfont\color{popink}OnBuch\color{poporange}+}}
\newcommand{\signatureonbuch}{%
  \begin{tcolorbox}[enhanced,colback=popink,colframe=popink,boxrule=2.2pt,arc=10pt,
      drop shadow=poporange,left=14pt,right=14pt,top=10pt,bottom=10pt,before skip=0pt,after skip=0pt]
    \tikz[baseline=-0.7ex]\node[circle,fill=popgreen,draw=popcream,line width=1.3pt,minimum size=22pt,inner sep=0]{\popcheck{white}};%
    \hspace{9pt}\begin{minipage}[c]{0.62\linewidth}
      {\popxbold\fontsize{12}{13}\selectfont\color{popcream}Signé\ \textbullet\ {\poplogo OnBuch\color{poporange}+}}\par\vspace{2pt}
      {\raggedright\popsemi\fontsize{8}{10.5}\selectfont\color{popline}\MakeUppercase{Équipe pédagogique OnBuch+}\\\MakeUppercase{Conforme au programme officiel MINESEC}\par}
    \end{minipage}\hfill
    {\poplogo\fontsize{22}{22}\selectfont\color{popcream}OnBuch\color{poporange}+}
  \end{tcolorbox}%
}

% ============================================================
% Page de garde
% #1 titre · #2 matière · #3 niveau · #4 module · #5 n° leçon · #6 durée ·
% #7 description / objectifs (texte ou itemize)
% ============================================================
\newcommand{\pagegarde}[7]{%
  \thispagestyle{empty}
  \newgeometry{a4paper,margin=1.7cm,top=1.6cm,bottom=1.4cm}%
  \begin{tikzpicture}[remember picture,overlay]
    \fill[poporange!18] ([xshift=-3.2cm,yshift=-3.6cm]current page.north east) circle (4.6cm);
    \fill[poppurple!14] ([xshift=2.4cm,yshift=7.6cm]current page.south west) circle (3.8cm);
  \end{tikzpicture}%
  {\poplogo\fontsize{30}{30}\selectfont\color{popink}OnBuch\color{poporange}+}\hfill
  \raisebox{0.6ex}{\poppill{Cours signé \textbullet\ Premium}{popink}{popcream}}\par
  \vspace{1pt}\popsquiggle{poppurple}\par
  \vspace{8pt}
  \begin{tcolorbox}[enhanced,colback=poporange,colframe=popink,boxrule=2.8pt,arc=13pt,
      drop shadow=popink,left=18pt,right=18pt,top=12pt,bottom=13pt,after skip=10pt]
    \poppill{#2 \textbullet\ #3}{popink}{popcream}\hspace{5pt}\poppill{#4}{white}{popink}\par\vspace{8pt}
    {\popdisplay\fontsize{14}{14}\selectfont\color{popink}LEÇON #5}\par\vspace{4pt}
    {\raggedright\hyphenpenalty=10000\exhyphenpenalty=10000\popdisplay\fontsize{25}{27}\selectfont\color{popcream}\MakeUppercase{#1}\par}
    \vspace{8pt}
    {\popxbold\fontsize{9.5}{9.5}\selectfont\color{popink}\MakeUppercase{Durée conseillée : #6}}
  \end{tcolorbox}
  \begin{tcolorbox}[enhanced,colback=white,colframe=popink,boxrule=2.2pt,arc=10pt,
      drop shadow=popink,left=16pt,right=16pt,top=10pt,bottom=10pt,after skip=8pt]
    \poppill{Dans cette leçon}{poppurple}{white}\par\vspace{5pt}
    \popsemi\fontsize{9.3}{12.6}\selectfont\color{popink}#7
  \end{tcolorbox}
  \noindent\begin{minipage}[t]{0.487\linewidth}
    \begin{tcolorbox}[enhanced,colback=popgreen,colframe=popink,boxrule=2.2pt,arc=10pt,
        drop shadow=popink,left=11pt,right=11pt,top=10pt,bottom=10pt,height=3.1cm,valign=center]
      \tcbox[colback=white,colframe=popink,boxrule=1.3pt,arc=5pt,boxsep=0pt,nobeforeafter,
        left=3pt,right=3pt,top=3pt,bottom=3pt,valign=center]{\qrcode[height=1.9cm]{https://play.google.com/store/apps/details?id=com.luvvix.onbuch&hl=fr}}%
      \hspace{8pt}\begin{minipage}[c]{0.5\linewidth}
        {\popdisplay\fontsize{11}{12.5}\selectfont\color{white}\MakeUppercase{Télécharge OnBuch}}\par\vspace{4pt}
        {\raggedright\popsemi\fontsize{8.4}{11}\selectfont\color{white}Gratuit \textbullet\ Google Play\\Tuteur IA, annales, concours, résultats\par}
      \end{minipage}
    \end{tcolorbox}
  \end{minipage}\hfill
  \begin{minipage}[t]{0.487\linewidth}
    \begin{tcolorbox}[enhanced,colback=poppurple,colframe=popink,boxrule=2.2pt,arc=10pt,
        drop shadow=popink,left=11pt,right=11pt,top=10pt,bottom=10pt,height=3.1cm,valign=center]
      \tikz[baseline=-0.7ex]\node[circle,fill=white,draw=popink,line width=1.3pt,minimum size=1.6cm,inner sep=0]{\popdisplay\fontsize{24}{24}\selectfont\color{poppurple}+};%
      \hspace{8pt}\begin{minipage}[c]{0.6\linewidth}
        {\popdisplay\fontsize{11}{12.5}\selectfont\color{white}\MakeUppercase{Passe à OnBuch+}}\par\vspace{4pt}
        {\raggedright\popsemi\fontsize{8.4}{11}\selectfont\color{white}Tous les cours signés, Tuteur IA illimité, fiches PDF — onglet OnBuch+ de l'app\par}
      \end{minipage}
    \end{tcolorbox}
  \end{minipage}
  \vspace{8pt}
  \popcontenu
  \vfill
  \signatureonbuch
  \restoregeometry
  \newpage
}

% Rangée « Ce cours contient » (page de garde)
\newcommand{\poptile}[3]{% #1 couleur #2 gros texte #3 légende
  \begin{tcolorbox}[enhanced,colback=#1,colframe=popink,boxrule=2pt,arc=9pt,shadow={2.5pt}{-2.5pt}{0pt}{popink},
      left=6pt,right=6pt,top=9pt,bottom=9pt,halign=center,valign=center,height=1.75cm,width=\linewidth,nobeforeafter]
    {\popdisplay\fontsize{12.5}{13}\selectfont\color{white}\MakeUppercase{#2}}\par\vspace{3pt}
    {\centering\popsemi\fontsize{7.8}{9.5}\selectfont\color{white}#3\par}
  \end{tcolorbox}}
\newcommand{\popcontenu}{%
  \par\noindent{\popxboldls\fontsize{7.5}{7.5}\selectfont\color{popmuted}CE COURS SIGNÉ CONTIENT}\par\vspace{7pt}
  \noindent\begin{minipage}[t]{0.235\linewidth}\poptile{popblue}{Cours}{complet, illustré, pas à pas}\end{minipage}\hfill
  \begin{minipage}[t]{0.235\linewidth}\poptile{poporange}{Méthodes}{et exercices résolus}\end{minipage}\hfill
  \begin{minipage}[t]{0.235\linewidth}\poptile{poppink}{Exercices}{type examen + corrigés}\end{minipage}\hfill
  \begin{minipage}[t]{0.235\linewidth}\poptile{popgreen}{Fiche bilan}{l'essentiel à retenir}\end{minipage}\par}

% Sommaire
\newtcolorbox{sommaire}{enhanced,colback=white,colframe=popink,boxrule=2.2pt,arc=10pt,
  drop shadow=popink,left=18pt,right=18pt,top=13pt,bottom=13pt,before skip=6pt,after skip=20pt,
  before upper={\poppill{Sommaire}{poppurple}{white}\par\vspace{9pt}\popsemi\fontsize{10.6}{14.5}\selectfont\color{popink}}}

% Signature de fin de cours — poussée tout en bas de la dernière page
\newcommand{\findecours}{%
  \par\vfill\nopagebreak
  \begin{center}\popsquiggle{poporange}\end{center}\vspace{4pt}
  \signatureonbuch
}
\AtBeginDocument{\def\llbracket{\mathopen{[\mkern-3.5mu[}}\def\rrbracket{\mathclose{]\mkern-3.5mu]}}} % intervalles d'entiers (glyphe ⟦ absent de la police)
% Glyphes absents de Fira Math : redéfinis à partir de glyphes disponibles
\AtBeginDocument{%
  \def\setminus{\mathbin{\backslash}}\let\smallsetminus\setminus
  \def\vdots{\mathord{\vbox{\baselineskip3.2pt\lineskiplimit0pt\kern4pt\hbox{.}\hbox{.}\hbox{.}}}}%
  \def\ddots{\mathinner{\mkern1mu\raise6.5pt\hbox{.}\mkern2mu\raise3.6pt\hbox{.}\mkern2mu\raise0.7pt\hbox{.}\mkern1mu}}%
  \def\triangle{\mathord{\scalebox{0.9}{$\symup{\Delta}$}}}%
  \def\nabla{\mathord{\raisebox{\depth}{\scalebox{1}[-1]{$\symup{\Delta}$}}}}%
  \def\checkmark{\popcheck{popdark}}%
  \def\star{\mathbin{\ast}}\def\bigstar{\mathord{\ast}}%
  \def\diamond{\mathbin{\scalebox{0.6}{$\Diamond$}}}%
  \def\bigcup{\mathop{\mathchoice{\raisebox{-0.3em}{\scalebox{1.7}{$\cup$}}}{\scalebox{1.25}{$\cup$}}{\cup}{\cup}}\displaylimits}%
  \def\bigcap{\mathop{\mathchoice{\raisebox{-0.3em}{\scalebox{1.7}{$\cap$}}}{\scalebox{1.25}{$\cap$}}{\cap}{\cap}}\displaylimits}%
  \def\lesseqqgtr{\mathrel{\lessgtr}}\def\bigoplus{\mathop{\oplus}\displaylimits}%
}
\providecommand{\ssi}{si et seulement si} % abréviation fréquente
\AtBeginDocument{\providecommand{\wideparen}{\overparen}} % arc de cercle

%% FIGFIX-BEGIN v4 — correctifs globaux des figures (docs/onbuch-generation/tools/figfix)
\usepackage{adjustbox}\usepackage{etoolbox}
% toute figure TikZ/pgfplots plus large que la ligne est réduite à la largeur disponible
\BeforeBeginEnvironment{tikzpicture}{\begin{adjustbox}{max width=\linewidth}}
\AfterEndEnvironment{tikzpicture}{\end{adjustbox}}
% légendes pgfplots lisibles par-dessus les courbes
\pgfplotsset{every axis/.append style={legend style={fill=white,fill opacity=0.92,text opacity=1,draw=black!15,inner sep=3pt}}}
% étiquettes d'arêtes (syntaxe edge["..."]) avec fond blanc
\tikzset{every edge quotes/.append style={fill=white,inner sep=1.2pt}}
%% FIGFIX-END
\begin{document}

\pagegarde{L'œil réduit : fonctionnement, défauts et correction}{Physique}{1ère C}{Module 3 : Optique}{N08}{8 h}{Cette leçon modélise l'œil humain par un système optique simple, l'œil réduit, afin d'expliquer la vision nette et le phénomène d'accommodation. Elle permet ensuite d'identifier les défauts de l'œil (myopie, hypermétropie, presbytie) et de proposer les corrections par lentilles adaptées.
\vspace{4pt}
\begin{multicols}{2}\small
\begin{itemize}
\item À la fin de cette leçon, je sais décrire les éléments essentiels de l'œil réel et les associer à ceux de l'œil réduit
\item À la fin de cette leçon, je sais schématiser l'œil réduit (centre optique, foyer image, rétine)
\item À la fin de cette leçon, je sais définir le punctum remotum (PR) et le punctum proximum (PP) et les situer sur un schéma
\item À la fin de cette leçon, je sais expliquer le phénomène d'accommodation comme une variation de la vergence du cristallin
\item À la fin de cette leçon, je sais identifier les défauts d'accommodation (myopie, hypermétropie, presbytie) à partir des positions du PR et du PP
\item À la fin de cette leçon, je sais choisir la nature de la lentille correctrice (convergente ou divergente) adaptée à chaque défaut
\end{itemize}\end{multicols}}

\begin{sommaire}
\begin{multicols}{2}
\begin{enumerate}
\item De l'œil réel à l'œil réduit
\item Accommodation, punctum remotum et punctum proximum
\item La myopie et sa correction
\item L'hypermétropie et sa correction
\item La presbytie et bilan des corrections
\item Méthode générale de résolution et synthèse
\item Activité d'intégration
\item Méthodes et exercices résolus
\item Exercices
\item Corrigés des exercices
\item Fiche bilan
\end{enumerate}
\end{multicols}
\end{sommaire}

\begin{objectifs}
\begin{itemize}
\item À la fin de cette leçon, je sais décrire les éléments essentiels de l'œil réel et les associer à ceux de l'œil réduit
\item À la fin de cette leçon, je sais schématiser l'œil réduit (centre optique, foyer image, rétine)
\item À la fin de cette leçon, je sais définir le punctum remotum (PR) et le punctum proximum (PP) et les situer sur un schéma
\item À la fin de cette leçon, je sais expliquer le phénomène d'accommodation comme une variation de la vergence du cristallin
\item À la fin de cette leçon, je sais identifier les défauts d'accommodation (myopie, hypermétropie, presbytie) à partir des positions du PR et du PP
\item À la fin de cette leçon, je sais choisir la nature de la lentille correctrice (convergente ou divergente) adaptée à chaque défaut
\item À la fin de cette leçon, je sais calculer la vergence de la lentille correctrice pour ramener le PR d'un œil myope à l'infini
\item À la fin de cette leçon, je sais construire graphiquement la marche des rayons lumineux dans un œil corrigé ou non corrigé
\end{itemize}
\end{objectifs}

\begin{prerequis}
\begin{itemize}
\item Lentilles minces convergentes et divergentes : foyers, vergence C = 1/f' (en dioptries)
\item Formule de conjugaison des lentilles minces et construction géométrique des images
\item Notion d'image réelle et d'image virtuelle
\item Modèle du rayon lumineux et propagation rectiligne de la lumière (classes antérieures)
\end{itemize}
\end{prerequis}

\newpage

\coursec{De l'œil réel à l'œil réduit}

\textbf{Situation d'accroche.} À Yaoundé, Maman Ngo Bell, 48 ans, remarque qu'elle doit éloigner son téléphone portable pour lire les messages de ses enfants. Son fils Étienne, élève en Première C, voit parfaitement le tableau au fond de la salle mais peine à lire son cahier posé sur la table. Quant à leur voisin, le chauffeur de taxi M. Fotso, il ne distingue les panneaux routiers que lorsqu'il s'en approche à moins de dix mètres. Comment un même organe, l'œil, peut-il présenter des comportements si différents ? Et comment l'opticien de Bonanjo choisit-il les verres correcteurs adaptés à chacun ?

Ces trois situations quotidiennes, que tu as certainement observées autour de toi, relèvent toutes de la même physique : celle d'un système optique qui doit former une image nette sur un écran. Pour comprendre ces défauts et prévoir leur correction, il faut d'abord connaître l'œil réel, puis le remplacer par un modèle simple sur lequel les lois des lentilles minces, étudiées dans la leçon précédente, pourront s'appliquer. C'est l'objet de cette première section : construire rigoureusement le modèle de \cle{l'œil réduit}.

\courssub{Description sommaire de l'œil réel}

L'œil humain est une sphère d'environ \SI{2,4}{cm} de diamètre, remplie de milieux transparents et bordée à l'arrière par une membrane sensible à la lumière. De l'avant vers l'arrière, la lumière traverse successivement plusieurs organes, chacun jouant un rôle précis.

\begin{description}
\item[La cornée.] Membrane transparente et bombée qui constitue la première surface réfringente de l'œil. C'est elle qui assure l'essentiel de la convergence (environ les deux tiers de la vergence totale), car le passage air--cornée présente le plus grand écart d'indice de tout le trajet.
\item[L'humeur aqueuse.] Liquide transparent situé entre la cornée et le cristallin. Elle maintient la pression interne de l'œil et participe légèrement à la réfraction.
\item[L'iris et la pupille.] L'iris est le diaphragme coloré de l'œil ; son ouverture centrale, la \cle{pupille}, laisse passer la lumière. L'iris se contracte ou se dilate selon l'éclairement : la pupille joue exactement le rôle du diaphragme d'un appareil photographique, en réglant la quantité de lumière qui pénètre dans l'œil (diamètre variable d'environ \SI{2}{mm} en pleine lumière à \SI{8}{mm} dans l'obscurité).
\item[Le cristallin.] Lentille biconvexe transparente, souple et \emph{déformable}, maintenue par des fibres (zonule) et déformée par les muscles ciliaires. C'est l'élément optique essentiel : en modifiant sa courbure, il modifie sa vergence et permet la mise au point.
\item[L'humeur vitrée.] Gel transparent qui remplit la cavité située entre le cristallin et la rétine. Elle maintient la forme sphérique de l'œil.
\item[La rétine.] Membrane tapissant le fond de l'œil, couverte de cellules photoréceptrices (cônes et bâtonnets). C'est \cle{l'écran récepteur} : c'est sur elle que doit se former l'image. La région de vision la plus précise, la fovéa, se situe sur l'axe optique.
\item[Le nerf optique.] Il transmet au cerveau, sous forme de signaux électriques, l'information lumineuse captée par la rétine.
\end{description}

\begin{popfigure}
\begin{center}
\begin{tikzpicture}[scale=1.05,>=Stealth]
% Globe oculaire
\draw[thick,popdark] (0,0) circle (2.2);
% Cornée (bombée à gauche)
\draw[thick,popblue] (-2.2,0) .. controls (-3.0,0.9) and (-3.0,-0.9) .. (-2.2,0);
% Iris / pupille
\draw[very thick,poporange] (-1.55,0.55) -- (-1.55,1.15);
\draw[very thick,poporange] (-1.55,-0.55) -- (-1.55,-1.15);
% Cristallin (lentille biconvexe)
\draw[thick,popgreen,fill=popgreenL] (-0.9,0) ellipse (0.35 and 0.95);
% Rétine (fond de l'œil)
\draw[very thick,popink] (1.1,1.9) arc (60:-60:2.2);
% Nerf optique
\draw[very thick,popgold] (2.2,0.35) -- (3.0,0.5);
\draw[very thick,popgold] (2.2,-0.35) -- (3.0,-0.5);
\draw[very thick,popgold] (3.0,0.5) -- (3.0,-0.5);
% Axe optique
\draw[dashed,popmuted] (-3.4,0) -- (2.4,0);
% Annotations
\draw[->,popblue] (-2.65,0.55) -- (-2.65,1.7) node[above,popblue]{\small cornée};
\draw[->,poporange] (-1.55,1.15) -- (-0.6,1.9) node[above right,poporange,xshift=-4mm]{\small iris et pupille};
\node[popgreen] at (-0.9,-1.6) {\small cristallin};
\draw[->,popgreen] (-0.9,-0.95) -- (-0.9,-1.45);
\draw[->,popink] (1.85,-1.15) -- (2.9,-1.6) node[right,popink]{\small rétine};
\node[popgold] at (3.55,0) {\small nerf};
\node[popgold] at (3.55,-0.35) {\small optique};
\node[popmuted] at (0.6,0.35) {\small humeur vitrée};
\node[popmuted] at (-1.9,-0.35) {\tiny humeur};
\node[popmuted] at (-1.9,-0.6) {\tiny aqueuse};
\end{tikzpicture}
\end{center}
\legende{Coupe schématique de l'œil réel : la lumière entre par la cornée (à gauche), traverse la pupille puis le cristallin, et vient former une image sur la rétine, au fond de l'œil.}
\end{popfigure}

\begin{attention}[Rôles à ne pas confondre]
Retiens bien dès maintenant la correspondance avec l'appareil photographique : le \cle{cristallin} joue le rôle de \textbf{lentille convergente} (objectif), la \cle{rétine} celui d'\textbf{écran} (capteur), et la \cle{pupille} celui de \textbf{diaphragme}. Confondre la rétine (écran) avec le cristallin (lentille) est l'erreur la plus fréquente dans cette leçon.
\end{attention}

\courssub{Pourquoi un modèle simplifié : la nécessité de l'œil réduit}

L'œil réel est un système optique complexe : la lumière subit des réfractions successives à la surface de la cornée, puis aux deux faces du cristallin, dans des milieux d'indices différents (air, humeur aqueuse, cristallin, humeur vitrée). Une étude exacte exigerait d'appliquer les lois de Snell--Descartes à chaque dioptre : c'est long, et surtout inutile pour résoudre les problèmes du Probatoire.

Or nous disposons déjà d'un outil puissant et simple : la théorie des \cle{lentilles minces convergentes}, avec sa relation de conjugaison
\[ \frac{1}{\overline{OA'}} - \frac{1}{\overline{OA}} = \frac{1}{\overline{OF'}} = C, \]
où $C$ est la vergence exprimée en dioptries ($\delta$). Toute la stratégie consiste donc à \textbf{remplacer l'ensemble cornée + humeurs + cristallin par une seule lentille mince convergente}, dont la vergence variable rend compte de la déformabilité du cristallin, et à placer un écran fixe à la place de la rétine. De tels modèles « d'œil réduit » ont été construits au XIX\textsuperscript{e} siècle, notamment par le physicien allemand Hermann von Helmholtz, puis perfectionnés par l'ophtalmologiste suédois Allvar Gullstrand (prix Nobel de physiologie ou médecine en 1911).

\begin{definition}[L'œil réduit]
L'\cle{œil réduit} est un modèle simplifié de l'œil constitué de deux éléments :
\begin{itemize}
\item une \textbf{lentille mince convergente unique}, de centre optique $O$, qui joue le rôle de l'ensemble cornée--cristallin ; sa vergence $C$ est \emph{variable} (le cristallin peut se déformer) ;
\item un \textbf{écran fixe}, la rétine, placé perpendiculairement à l'axe optique principal, à une distance \emph{constante} de la lentille.
\end{itemize}
La distance lentille--rétine vaut environ $15$ à \SI{17}{mm} ; on retiendra la valeur indicative
\[ \overline{OA'}_{\text{rétine}} \approx \SI{15}{mm} = \SI{1,5e-2}{m} \quad \text{(œil réduit au repos)}. \]
\end{definition}

\begin{popfigure}
\begin{center}
\begin{tikzpicture}[scale=1.0,>=Stealth]
% Axe optique
\draw[->,popmuted] (-4.5,0) -- (5.2,0) node[below left]{\small axe optique};
% Lentille mince (cristallin)
\draw[very thick,popgreen] (0,-1.9) -- (0,1.9);
\draw[very thick,popgreen,->] (0,1.9) -- (-0.35,1.55);
\draw[very thick,popgreen,->] (0,1.9) -- (0.35,1.55);
\draw[very thick,popgreen,->] (0,-1.9) -- (-0.35,-1.55);
\draw[very thick,popgreen,->] (0,-1.9) -- (0.35,-1.55);
\node[popgreen] at (0,2.35) {\small lentille (cristallin)};
% Centre optique
\fill[popdark] (0,0) circle (1.5pt) node[below left=1pt]{\small $O$};
% Foyers
\fill[popblue] (1.5,0) circle (1.5pt) node[below]{\small $F'$};
\fill[popblue] (-1.5,0) circle (1.5pt) node[below]{\small $F$};
% Rétine (écran)
\draw[very thick,popink] (1.5,-1.9) -- (1.5,1.9);
\node[popink,rotate=90] at (1.85,0) {\small rétine (écran fixe)};
% Distance O - rétine
\draw[<->,poporange,thick] (0,-2.25) -- (1.5,-2.25) node[midway,below]{\small $\approx \SI{15}{mm}$ (constante)};
% Objet et rayons
\fill[popdark] (-3.75,0) circle (1.2pt);
\draw[very thick,popdark,->] (-3.75,0) -- (-3.75,1.3) node[midway,left]{\small $AB$};
% Rayon parallèle
\draw[thick,popblue] (-3.75,1.3) -- (0,1.3);
\draw[thick,popblue,->] (0,1.3) -- (1.5,0);
% Rayon par O
\draw[thick,poppurple,->] (-3.75,1.3) -- (1.5,-0.433);
% Image
\fill[popink] (1.5,-0.433) circle (1.5pt);
\draw[very thick,popink,->] (1.5,0) -- (1.5,-0.433) node[midway,right]{\small $A'B'$};
\node[popmuted] at (-2.2,1.9) {\small rayon parallèle à l'axe};
\node[popmuted] at (-2.6,-1.3) {\small rayon passant par $O$};
\end{tikzpicture}
\end{center}
\legende{Modèle de l'œil réduit : une lentille mince convergente de centre $O$ forme l'image $A'B'$ de l'objet $AB$ sur l'écran fixe (la rétine), situé à environ \SI{15}{mm} de $O$. L'image est renversée.}
\end{popfigure}

\courssub{Vergence de l'œil au repos et condition de vision nette}

Considérons un œil \textbf{normal} (on dit \emph{emmétrope}) \textbf{au repos}, c'est-à-dire dont les muscles ciliaires ne travaillent pas : le cristallin est alors le moins bombé, sa vergence est minimale. Dans cet état, l'œil voit nettement les objets très éloignés (à l'infini). Pour un objet à l'infini, l'image se forme dans le plan focal image : $\overline{OF'} = f'$. Comme cette image doit se former sur la rétine, on a $f' \approx \SI{15}{mm}$, d'où la vergence au repos :
\[ C_{\text{repos}} = \frac{1}{f'} = \frac{1}{\num{1,5e-2}} \approx \SI{67}{\delta}. \]
Selon les modèles d'œil réduit (distance lentille--rétine de $15$ à \SI{17}{mm}), on trouve une vergence au repos comprise entre $60$ et $67$ dioptries.

\begin{aretenir}[Ordres de grandeur à connaître]
\begin{itemize}
\item Distance cristallin--rétine (constante) : environ \SI{15}{mm} (valeur indicative de l'œil réduit au repos).
\item Vergence de l'œil au repos : $C \approx 60$ à $\SI{67}{\delta}$ (ordre de grandeur à retenir : $\approx \SI{60}{\delta}$).
\end{itemize}
Vérification dimensionnelle : la vergence est l'inverse d'une longueur, $[C] = \mathrm{m}^{-1}$ ; l'unité SI est la dioptrie ($\delta$), avec $1\ \delta = 1\ \mathrm{m}^{-1}$.
\end{aretenir}

\begin{exemplebox}[Calcul de la vergence au repos]
Un œil réduit au repos a une distance lentille--rétine de $\SI{16,7}{mm}$. Sa vergence au repos vaut :
\[ C = \frac{1}{f'} = \frac{1}{\num{1,67e-2}\ \mathrm{m}} \approx \SI{60}{\delta}. \]
C'est l'ordre de grandeur annoncé. Un œil est donc un système optique extrêmement convergent : à titre de comparaison, les lunettes correctrices ont des vergences de quelques dioptries seulement.
\end{exemplebox}

Venons-en maintenant à la condition fondamentale qui gouvernera toute la leçon. La rétine est un écran \emph{fixe} : sa distance au cristallin ne change jamais. Pour qu'un objet soit vu nettement, il faut et il suffit que la lentille forme l'image de cet objet \emph{exactement} sur cet écran.

\begin{propriete}[Condition de vision nette (admise)]
La vision est nette si et seulement si l'image $A'B'$ de l'objet $AB$ observé, donnée par le cristallin, se forme \textbf{exactement sur la rétine}. Si l'image se forme en avant ou en arrière de la rétine, la vision est floue : c'est un défaut de l'œil.
\end{propriete}

Cette propriété, jointe à la relation de conjugaison des lentilles minces, est la clé de tous les exercices de la leçon. Remarque la conséquence immédiate : puisque la distance image $\overline{OA'}$ est imposée (constante, égale à la distance cristallin--rétine) alors que la distance objet $\overline{OA}$ varie selon la position de l'objet regardé, la relation de conjugaison
\[ \frac{1}{\overline{OA'}} - \frac{1}{\overline{OA}} = C \]
montre que la vergence $C$ doit nécessairement \textbf{varier} avec la position de l'objet. C'est précisément ce que réalise le cristallin en se déformant : c'est le phénomène d'\cle{accommodation}, que nous étudierons en détail dans la section suivante.

\begin{attention}[Ce qui varie et ce qui ne varie pas]
Piège classique du Probatoire : la distance cristallin--rétine ne varie \textbf{jamais} (l'œil ne change pas de taille pour faire la mise au point). Quand l'œil passe de la vision de loin à la vision de près, c'est la \textbf{vergence} du cristallin qui augmente (le cristallin se bombe, ses rayons de courbure diminuent), et non la position de la rétine. Écrire dans une copie que « la rétine recule » est une erreur éliminatoire.
\end{attention}

\begin{exoresolu}[Vergence nécessaire pour voir un objet proche]
\textbf{Énoncé.} Un œil réduit est modélisé par une lentille mince de centre $O$ placée à $\SI{15}{mm}$ devant la rétine. Cet œil observe un objet $AB$ situé à $\overline{OA} = -\SI{25}{cm}$ (distance de lecture confortable). Déterminer la vergence $C$ que doit prendre le cristallin pour que la vision soit nette. Comparer à la vergence au repos ($\SI{60}{\delta}$) et conclure.
\tcblower
\textbf{Solution.} La vision est nette si l'image se forme sur la rétine, donc $\overline{OA'} = +\SI{15}{mm} = +\num{1,5e-2}\ \mathrm{m}$. La relation de conjugaison s'écrit :
\begin{align*}
C &= \frac{1}{\overline{OA'}} - \frac{1}{\overline{OA}} \\
&= \frac{1}{\num{1,5e-2}} - \frac{1}{-\num{0,25}} \\
&= 66,7 + 4,0 \\
&\approx \SI{71}{\delta}.
\end{align*}
Au repos, la vergence vaut environ $\SI{60}{\delta}$ (objet à l'infini) ; pour lire à \SI{25}{cm}, elle doit atteindre environ $\SI{71}{\delta}$. Le cristallin doit donc \textbf{augmenter sa vergence d'environ 11 dioptries} en se bombant : c'est l'accommodation. Plus l'objet est proche, plus l'effort d'accommodation est grand — d'où la fatigue visuelle lors d'une lecture prolongée.
\end{exoresolu}

\begin{savaistu}[Une image à l'envers, redressée par le cerveau]
Comme le montre la construction géométrique de la figure de l'œil réduit, l'image formée par une lentille convergente sur la rétine est \textbf{renversée} : le haut d'un objet est projeté en bas de la rétine, et la droite à gauche. Pourtant, nous voyons le monde à l'endroit ! C'est le cerveau qui, dès la naissance, apprend à réinterpréter les signaux du nerf optique et « remet l'image à l'endroit ». Des expériences célèbres avec des lunettes inversant les images ont montré qu'au bout de quelques jours, le cerveau d'un adulte s'adapte et rétablit une vision correcte. L'optique s'arrête à la rétine ; la vision, elle, se termine dans le cerveau.
\end{savaistu}

\begin{methode}[Ce qu'il faut retenir de cette section]
Pour aborder la suite de la leçon (défauts et corrections), garde en mémoire la démarche en quatre points :
\begin{enumerate}
\item L'œil réel est remplacé par l'\textbf{œil réduit} : une lentille convergente de centre $O$ (le cristallin) et un écran fixe (la rétine) à $\approx \SI{15}{mm}$.
\item La distance lentille--rétine est \textbf{constante} ; seule la \textbf{vergence} varie, grâce à la déformation du cristallin.
\item Vergence au repos : $C \approx 60$ à $\SI{67}{\delta}$ (objet à l'infini vu net).
\item Condition de vision nette : l'image doit se former \textbf{exactement sur la rétine} ; on applique alors la relation de conjugaison avec $\overline{OA'} = +\SI{15}{mm}$.
\end{enumerate}
Dans les sections suivantes, nous verrons comment l'accommodation permet à l'œil normal de voir net de l'infini au punctum proximum, puis nous analyserons les défauts de Maman Ngo Bell, d'Étienne et de M. Fotso — myopie, hypermétropie, presbytie — et les verres correcteurs que l'opticien de Bonanjo leur prescrira.
\end{methode}

\coursec{Accommodation, punctum remotum et punctum proximum}

Dans la section précédente, nous avons remplacé l'œil réel, organe complexe (cornée, humeur aqueuse, cristallin, humeur vitrée), par un modèle simple et puissant : \cle{l'œil réduit}, constitué d'une lentille convergente (le cristallin) et d'un écran (la rétine) placé à une distance fixe $d \approx \SI{15}{\milli\meter}$ derrière la lentille. Mais ce modèle soulève immédiatement une question fondamentale : si la distance lentille--écran est \emph{fixe}, comment l'œil peut-il voir net aussi bien un avion dans le ciel que le texte de ton cahier à \SI{25}{\centi\meter} ? C'est toute la magie de \cle{l'accommodation}, que nous allons maintenant étudier en détail.

\courssub{Le problème posé : une distance image fixe, des objets à toutes les distances}

\begin{experience}[Un constat qui dérange]
Prenons une lentille convergente de vergence $C = \SI{15}{\dioptre}$ et un écran placé à \SI{10}{\centi\meter} derrière elle. On place une bougie allumée à différentes distances de la lentille et on observe l'écran.
\begin{itemize}
\item Bougie à \SI{20}{\centi\meter} : l'image sur l'écran est \emph{nette}.
\item Bougie à \SI{50}{\centi\meter} : l'image sur l'écran devient \emph{floue} ; pour la rendre nette, il faut \emph{rapprocher} l'écran.
\item Bougie à \SI{10}{\centi\meter} : l'image est de nouveau floue ; il faut \emph{reculer} l'écran.
\end{itemize}
\textbf{Interprétation.} D'après la formule de conjugaison des lentilles minces,
\[ \frac{1}{\overline{OA'}} - \frac{1}{\overline{OA}} = C, \]
lorsque l'objet se rapproche ($\overline{OA}$ varie), la position de l'image $\overline{OA'}$ varie aussi. Or, dans l'œil, la rétine est \emph{solidement fixée} au fond du globe oculaire : elle ne peut ni avancer ni reculer. L'œil ne peut donc pas jouer sur $\overline{OA'}$.
\end{experience}

Ce constat est capital. Dans un appareil photographique ancien, on fait la mise au point en \emph{déplaçant} l'objectif par rapport à la pellicule. Dans l'œil, c'est impossible : la distance cristallin--rétine vaut environ \SI{15}{\milli\meter} et ne change jamais. Il ne reste alors qu'une seule variable d'ajustement : la \cle{vergence} $C$ de la lentille, c'est-à-dire la courbure du cristallin.

\begin{attention}[Le point de départ de tout le raisonnement]
La distance cristallin--rétine est \textbf{fixe} ($\approx \SI{15}{\milli\meter}$). Pour que l'image d'un objet reste sur la rétine quelle que soit la distance de l'objet, l'œil doit obligatoirement modifier sa \textbf{vergence}. Toute la physiologie de la vision nette découle de cette contrainte géométrique.
\end{attention}

\courssub{L'accommodation : le cristallin, une lentille déformable}

\begin{definition}[Accommodation]
L'\cle{accommodation} est le phénomène par lequel l'œil modifie la \textbf{courbure du cristallin} (et donc sa vergence) grâce à l'action des \textbf{muscles ciliaires}, afin que l'image d'un objet se forme exactement sur la rétine, quelle que soit la distance de l'objet.
\end{definition}

Le cristallin est une lentille biologique souple et élastique, suspendue par des fibres (les zonules) à un anneau de muscles : les \cle{muscles ciliaires}. Le mécanisme est le suivant :

\begin{itemize}
\item \textbf{Muscles ciliaires relâchés} : les zonules sont tendues, elles \emph{tirent} sur le cristallin qui s'\emph{aplatit}. Ses rayons de courbure augmentent, sa vergence \emph{diminue} : c'est la vergence \textbf{minimale} $C_{\min}$. L'œil est dit \emph{au repos}.
\item \textbf{Muscles ciliaires contractés} : les zonules se détendent, le cristallin, par son élasticité propre, se \emph{bombe} (il devient plus épais, plus courbé). Sa vergence \emph{augmente} : l'œil \emph{accommode}. Au maximum de contraction, la vergence atteint sa valeur \textbf{maximale} $C_{\max}$.
\end{itemize}

\begin{aretenir}[Le réflexe d'accommodation en une phrase]
Accommoder, c'est \textbf{bomber le cristallin} pour \textbf{augmenter la vergence} de l'œil afin de voir net un objet \textbf{proche}. Au repos, le cristallin est aplati, la vergence est minimale, et l'œil regarde au loin.
\end{aretenir}

Vérifions la cohérence physique de ce mécanisme par la formule de conjugaison. Pour un objet proche, $\overline{OA}$ est petit en valeur absolue, donc $-\dfrac{1}{\overline{OA}} = \dfrac{1}{|\overline{OA}|}$ est grand : il faut une vergence $C$ plus grande pour conserver le même $\overline{OA'} = \SI{15}{\milli\meter}$. Bomber le cristallin augmente bien la vergence : le mécanisme physiologique répond exactement à l'exigence optique.

\begin{popfigure}
\begin{center}
\begin{tikzpicture}[scale=0.9, every node/.style={scale=0.85}]
% ===== Panneau 1 : oeil au repos =====
% Retine
\draw[line width=1.6pt, popink] (3.2,-1.6) arc[start angle=-60, end angle=60, radius=1.85];
\node[popink, align=center] at (4.1,0) {rétine};
% Cristallin aplati (lentille mince)
\draw[line width=1.4pt, popblue] (0,-1.1) .. controls (0.18,0) .. (0,1.1);
\draw[line width=1.4pt, popblue] (0,-1.1) .. controls (-0.18,0) .. (0,1.1);
\node[popblue, align=center] at (0,-1.5) {cristallin\\aplati};
% Axe optique
\draw[popmuted, dashed] (-5.6,0) -- (3.4,0);
% Rayons paralleles
\foreach \y in {-0.7,0,0.7}{
  \draw[poporange, -{Stealth[length=2mm]}] (-5.4,\y) -- (0,\y);
  \draw[poporange] (0,\y) -- (2.95,0);
}
\node[poporange, align=center] at (-3.6,1.15) {rayons parallèles\\(objet à l'infini)};
% Foyer image sur retine
\fill[popdark] (2.95,0) circle (1.6pt);
\node[popdark, below] at (2.95,-0.15) {$F'$ sur la rétine};
\node[align=center, font=\bfseries] at (-1,-2.6) {Œil au repos : $C = C_{\min}$};
% ===== Panneau 2 : oeil accommodant =====
\begin{scope}[xshift=9.5cm]
\draw[line width=1.6pt, popink] (3.2,-1.6) arc[start angle=-60, end angle=60, radius=1.85];
\node[popink] at (4.1,0) {rétine};
% Cristallin bombe
\draw[line width=1.4pt, popblue, fill=popblueL] (0,-1.15) .. controls (0.42,0) .. (0,1.15);
\draw[line width=1.4pt, popblue, fill=popblueL] (0,-1.15) .. controls (-0.42,0) .. (0,1.15);
\node[popblue, align=center] at (0,-1.6) {cristallin\\bombé};
\draw[popmuted, dashed] (-5.6,0) -- (3.4,0);
% Objet proche
\draw[popgreen, line width=1.4pt, -{Stealth[length=2.2mm]}] (-3.4,0) -- (-3.4,1.0);
\node[popgreen, below] at (-3.4,-0.1) {objet proche};
% Rayons divergents
\foreach \y in {-0.75,0,0.75}{
  \draw[poporange] (-3.4,1.0) -- (0,\y);
  \draw[poporange] (0,\y) -- (2.95,0);
}
\draw[poporange, -{Stealth[length=2mm]}] (-2.6,0.76) -- (-2.2,0.65);
\fill[popdark] (2.95,0) circle (1.6pt);
\node[popdark, below] at (2.95,-0.15) {image sur la rétine};
\node[align=center, font=\bfseries] at (-1,-2.6) {Œil accommodant : $C = C_{\max}$};
\end{scope}
\end{tikzpicture}
\end{center}
\legende{Œil réduit au repos (à gauche) : le cristallin aplati ($C_{\min}$) focalise les rayons parallèles venus de l'infini sur la rétine. Œil accommodant (à droite) : le cristallin bombé ($C_{\max}$) fait converger sur la rétine les rayons divergents venus d'un objet proche. La distance cristallin--rétine est la même dans les deux cas.}
\end{popfigure}

\courssub{Le punctum remotum (PR) : la limite de la vision de loin}

\begin{definition}[Punctum remotum]
Le \cle{punctum remotum} (PR) est le point de l'axe optique \textbf{le plus éloigné} dont l'œil voit une image nette, \textbf{sans accommoder} (muscles ciliaires relâchés, vergence $C_{\min}$).
\end{definition}

Lorsque l'œil est au repos, sa vergence vaut $C_{\min}$. La formule de conjugaison appliquée à l'œil réduit, avec l'image sur la rétine ($\overline{OA'} = +\SI{15}{\milli\meter}$), donne la position de l'objet conjugué :
\[ \frac{1}{\overline{OA'}} - \frac{1}{\overline{OA}_{\text{PR}}} = C_{\min}
\quad\Longrightarrow\quad
\overline{OA}_{\text{PR}} = \frac{1}{\dfrac{1}{\overline{OA'}} - C_{\min}}. \]

\begin{definition}[Œil emmétrope]
Un œil est dit \cle{emmétrope} (normal) lorsque, au repos, il voit nettement les objets situés à \textbf{l'infini} : son punctum remotum est \emph{rejeté à l'infini} ($\overline{OA}_{\text{PR}} \to -\infty$).
\end{definition}

Pour l'œil emmétrope, $\dfrac{1}{\overline{OA}_{\text{PR}}} = 0$, donc la formule de conjugaison se réduit à :
\[ C_{\min} = \frac{1}{\overline{OA'}} = \frac{1}{\SI{15e-3}{\meter}} \approx \num{66,7}\ \delta. \]
C'est un résultat remarquable : la vergence au repos de l'œil normal est entièrement déterminée par la profondeur du globe oculaire. Un œil emmétrope regarde les étoiles, le sommet du mont Cameroun ou un avion \emph{sans aucun effort musculaire} : c'est pour cela que la vision de loin ne fatigue pas.

\courssub{Le punctum proximum (PP) : la limite de la vision de près}

\begin{definition}[Punctum proximum]
Le \cle{punctum proximum} (PP) est le point de l'axe optique \textbf{le plus proche} dont l'œil voit une image nette, avec une \textbf{accommodation maximale} (muscles ciliaires contractés au maximum, vergence $C_{\max}$).
\end{definition}

Lorsque tu rapproches un texte de ton œil, les muscles ciliaires se contractent de plus en plus pour bomber le cristallin. Mais un muscle ne peut pas se contracter indéfiniment : en dessous d'une certaine distance, la vergence maximale $C_{\max}$ ne suffit plus à ramener l'image sur la rétine, et le texte devient flou. Cette distance limite est la distance du PP.

\begin{propriete}[Limites de l'œil emmétrope]
Pour un œil \textbf{emmétrope jeune} :
\begin{itemize}
\item le punctum remotum est à l'\textbf{infini} : $\text{PR} \to \infty$ ;
\item le punctum proximum est situé à environ \SI{25}{\centi\meter} de l'œil : $d_{\text{PP}} \approx \SI{25}{\centi\meter}$.
\end{itemize}
La distance $d_{\text{PP}} = \SI{25}{\centi\meter}$ est appelée \cle{distance minimale de vision distincte} ; on la note souvent $d_m$.
\end{propriete}

\begin{attention}[Accommoder fatigue]
L'accommodation est un \textbf{travail musculaire} : les muscles ciliaires se contractent en permanence tant que tu regardes de près. Une lecture prolongée à moins de \SI{25}{\centi\meter} (téléphone collé au visage, cahier trop proche) impose une accommodation forte et durable, source de \textbf{fatigue visuelle} : picotements, maux de tête, vision trouble passagère. Règle d'hygiène : garde tes documents à au moins \SI{30}{\centi\meter} et fais des pauses en regardant au loin (l'œil se repose... sans rien faire !).
\end{attention}

\courssub{Domaine de vision nette et amplitude d'accommodation}

\begin{definition}[Domaine de vision nette]
Le \cle{domaine de vision nette} (ou champ de netteté) est la région de l'espace, sur l'axe optique, comprise entre le punctum proximum (PP) et le punctum remotum (PR). Tout objet placé dans ce domaine est vu net, moyennant une accommodation adaptée ; tout objet situé en dehors est vu flou.
\end{definition}

Pour l'œil emmétrope, le domaine de vision nette s'étend de \SI{25}{\centi\meter} à l'infini : c'est un domaine immense, qui explique l'extraordinaire souplesse de notre vision.

\begin{popfigure}
\begin{center}
\begin{tikzpicture}[scale=1, every node/.style={scale=0.85}]
% Axe optique
\draw[-{Stealth[length=2.5mm]}, line width=1pt, popdark] (-1.2,0) -- (11.5,0) node[below left=1pt and -18pt] {axe optique};
% Oeil
\draw[line width=1.3pt, popblue, fill=popblueL] (0,-0.55) .. controls (0.25,0) .. (0,0.55);
\draw[line width=1.3pt, popblue, fill=popblueL] (0,-0.55) .. controls (-0.25,0) .. (0,0.55);
\draw[line width=1.3pt, popink] (0.85,-0.7) arc[start angle=-55, end angle=55, radius=0.85];
\node[popblue, below] at (0,-0.7) {œil};
% PP
\draw[line width=1.2pt, popgreen] (2.5,-0.35) -- (2.5,0.35);
\node[popgreen, above, align=center] at (2.5,0.4) {PP\\\SI{25}{\centi\meter}};
% Graduations
\foreach \x/\t in {5/{\SI{50}{\centi\meter}}, 7.5/{\SI{1}{\meter}}}{
  \draw[popmuted] (\x,-0.15) -- (\x,0.15);
  \node[popmuted, below] at (\x,-0.2) {\t};
}
% PR a l'infini
\node[poporange, above] at (10.6,0.15) {PR : $\infty$};
\draw[poporange, line width=1.2pt] (10.4,-0.3) -- (10.4,0.3);
% Zone de vision nette
\draw[line width=4pt, popgreen!45, opacity=0.7] (2.5,0.75) -- (10.4,0.75);
\draw[-{Stealth[length=2mm]}, popgreen, line width=1pt] (10.4,0.75) -- (11.2,0.75);
\node[popgreen!60!black, above, align=center] at (6.4,1.0) {domaine de vision nette};
% Zone floue proche
\draw[line width=4pt, poppink!55, opacity=0.7] (0.15,0.75) -- (2.5,0.75);
\node[poppink!70!black, above, align=center] at (1.3,1.0) {flou};
\end{tikzpicture}
\end{center}
\legende{Domaine de vision nette de l'œil emmétrope : il s'étend du PP (\SI{25}{\centi\meter}) au PR (rejeté à l'infini). Entre l'œil et le PP, la vision est floue car l'accommodation maximale ne suffit plus.}
\end{popfigure}

Pour quantifier la capacité d'accommodation de l'œil, on introduit une grandeur essentielle :

\begin{definition}[Amplitude d'accommodation]
L'\cle{amplitude d'accommodation} $A$ est la différence entre la vergence maximale et la vergence minimale de l'œil :
\[ A = C_{\max} - C_{\min} \quad \text{(en dioptries, } \delta\text{)}. \]
Elle mesure la \emph{réserve de puissance} dont dispose l'œil pour voir de près.
\end{definition}

\begin{propriete}[Vieillissement de l'amplitude d'accommodation]
L'amplitude d'accommodation \textbf{diminue avec l'âge} : le cristallin perd progressivement son élasticité et ne se bombe plus aussi facilement. Un jeune de \num{15} ans possède typiquement $A \approx \num{8}\ \text{à}\ \num{12}\ \delta$ ; à \num{45} ans, il ne reste souvent que $\num{3}\ \text{à}\ \num{4}\ \delta$, et le PP recule au-delà de \SI{25}{\centi\meter} : c'est le début de la \emph{presbytie}, que nous étudierons plus loin.
\end{propriete}

\courssub{Exemple chiffré : vérification par la formule de conjugaison}

\begin{exemplebox}[De $C_{\min}$ à $C_{\max}$ pour un œil emmétrope]
Considérons un œil emmétrope dont la distance cristallin--rétine vaut $\overline{OA'} = \SI{15}{\milli\meter}$, le punctum proximum $d_{\text{PP}} = \SI{25}{\centi\meter}$ et l'amplitude d'accommodation $A = C_{\max} - C_{\min} = \num{4}\ \delta$. Vérifions la cohérence de ces données.

\textbf{1. Vergence au repos.} Au repos, l'œil emmétrope voit l'infini ($\overline{OA}_{\text{PR}} \to -\infty$) :
\begin{align*}
C_{\min} &= \frac{1}{\overline{OA'}} - \frac{1}{\overline{OA}_{\text{PR}}}
= \frac{1}{\SI{15e-3}{\meter}} - 0 \\
&\approx \num{66,7}\ \delta.
\end{align*}

\textbf{2. Vergence maximale.} Avec $A = \num{4}\ \delta$ :
\[ C_{\max} = C_{\min} + A = \num{66,7} + \num{4} = \num{70,7}\ \delta. \]

\textbf{3. Position du PP prédite.} À accommodation maximale, l'image doit rester sur la rétine :
\begin{align*}
\frac{1}{\overline{OA}_{\text{PP}}} &= \frac{1}{\overline{OA'}} - C_{\max}
= \num{66,7} - \num{70,7} = -\num{4}\ \delta,\\
\overline{OA}_{\text{PP}} &= \frac{1}{-\num{4}} = -\SI{0,25}{\meter} = -\SI{25}{\centi\meter}.
\end{align*}
Le signe négatif confirme que l'objet est réel, placé devant l'œil, à \SI{25}{\centi\meter} : on retrouve exactement le PP de l'œil normal. Le modèle de l'œil réduit est parfaitement cohérent.
\end{exemplebox}

\begin{attention}[Pièges fréquents dans les calculs]
\begin{itemize}
\item \textbf{Signe de $\overline{OA}$} : l'objet est toujours \emph{devant} l'œil, donc $\overline{OA} < 0$ dans la convention des lentilles. Écrire $\overline{OA} = +\SI{25}{\centi\meter}$ est une erreur classique qui fausse tout le calcul.
\item \textbf{Unités} : la vergence s'exprime en dioptries ($\delta = \text{m}^{-1}$) ; toutes les distances doivent être converties en \emph{mètres} avant le calcul (\SI{15}{\milli\meter} = \SI{0,015}{\meter}).
\item \textbf{Ne pas confondre} $A = C_{\max} - C_{\min}$ (amplitude, en $\delta$) avec la distance PP--PR (en mètres) : ce sont deux grandeurs de nature différente.
\end{itemize}
\end{attention}

\courssub{Méthode : un objet est-il vu net ?}

\begin{methode}[Utiliser le domaine de vision nette]
Pour déterminer si un œil voit nettement un objet, et dans quelles conditions :
\begin{enumerate}
\item \textbf{Repérer les limites} : noter la distance du PP et celle du PR (données de l'énoncé ou résultats d'un calcul par la formule de conjugaison).
\item \textbf{Localiser l'objet} : mesurer sa distance $d$ à l'œil.
\item \textbf{Comparer} :
\begin{itemize}
\item si $d_{\text{PP}} \leq d \leq d_{\text{PR}}$ : l'objet est \textbf{vu net}, avec une accommodation d'autant plus forte que $d$ est proche du PP ;
\item si $d < d_{\text{PP}}$ : l'objet est \textbf{trop proche}, vision floue (accommodation insuffisante) ;
\item si $d > d_{\text{PR}}$ : l'objet est \textbf{trop loin}, vision floue (l'œil ne peut pas \emph{dés}accommoder en dessous de $C_{\min}$).
\end{itemize}
\item \textbf{Conclure} en précisant, si besoin, la vergence nécessaire : $C = \dfrac{1}{\overline{OA'}} - \dfrac{1}{\overline{OA}}$.
\end{enumerate}
\end{methode}

\begin{exoresolu}[Un élève lit-il dans de bonnes conditions ?]
\textbf{Énoncé.} Un élève emmétrope ($C_{\min} = \num{66,7}\ \delta$, $A = \num{8}\ \delta$, distance cristallin--rétine \SI{15}{\milli\meter}) lit son manuel de physique placé à \SI{20}{\centi\meter} de ses yeux.
\begin{enumerate}
\item Déterminer la position de son punctum proximum.
\item Voit-il le manuel nettement ? Quelle vergence son œil adopte-t-il ?
\end{enumerate}
\tcblower
\textbf{Solution.}
\textbf{1. Position du PP.} La vergence maximale vaut :
\[ C_{\max} = C_{\min} + A = \num{66,7} + \num{8} = \num{74,7}\ \delta. \]
La formule de conjugaison, avec l'image sur la rétine ($\overline{OA'} = \SI{0,015}{\meter}$), donne :
\begin{align*}
\frac{1}{\overline{OA}_{\text{PP}}} &= \frac{1}{\overline{OA'}} - C_{\max} = \num{66,7} - \num{74,7} = -\num{8}\ \delta,\\
\overline{OA}_{\text{PP}} &= -\frac{1}{\num{8}} = -\SI{0,125}{\meter} = -\SI{12,5}{\centi\meter}.
\end{align*}
Le PP de cet élève est à \SI{12,5}{\centi\meter} devant l'œil : grâce à sa grande amplitude d'accommodation (il est jeune !), il voit net beaucoup plus près que la norme de \SI{25}{\centi\meter}.

\textbf{2. Vision du manuel à \SI{20}{\centi\meter}.} On a $d = \SI{20}{\centi\meter} > d_{\text{PP}} = \SI{12,5}{\centi\meter}$, et le PR est à l'infini (œil emmétrope) : l'objet est \textbf{dans le domaine de vision nette}, donc vu net. La vergence nécessaire est :
\begin{align*}
C &= \frac{1}{\overline{OA'}} - \frac{1}{\overline{OA}}
= \frac{1}{\SI{0,015}{\meter}} - \frac{1}{-\SI{0,20}{\meter}} \\
&= \num{66,7} + \num{5} = \num{71,7}\ \delta.
\end{align*}
L'œil accommode de $\num{71,7} - \num{66,7} = \num{5}\ \delta$, soit \SI{62}{\%} de son amplitude maximale : la lecture est possible mais assez fatigante si elle se prolonge. Il serait plus confortable de tenir le manuel vers \SI{30}{\centi\meter}.
\end{exoresolu}

\begin{savaistu}[Pourquoi les enfants rapprochent-ils tant leurs bandes dessinées ?]
Un enfant de \num{10} ans possède une amplitude d'accommodation pouvant dépasser \num{12}\ \delta$ : son PP peut descendre jusqu'à \SI{7}{\centi\meter} ! Il peut donc lire le nez collé sur sa bande dessinée sans effort apparent. Ce n'est \emph{pas} le signe d'un défaut de vision. En revanche, vers \num{45}--\num{50} ans, l'amplitude tombe sous \num{2}\ \delta$ et le PP recule au-delà d'un mètre : les parents « allongent le bras » pour lire les messages sur leur téléphone... C'est la presbytie, conséquence naturelle du durcissement du cristallin, et non une maladie.
\end{savaistu}

\begin{aretenir}[L'essentiel de la section]
\begin{itemize}
\item La distance cristallin--rétine étant \textbf{fixe}, l'œil voit net à toutes les distances en modifiant sa \textbf{vergence} : c'est l'\textbf{accommodation}, réalisée par la déformation du cristallin sous l'action des muscles ciliaires.
\item \textbf{PR} : point le plus lointain vu net \emph{sans} accommoder (vergence $C_{\min}$). Pour l'œil emmétrope, le PR est à l'\textbf{infini}.
\item \textbf{PP} : point le plus proche vu net avec accommodation \emph{maximale} (vergence $C_{\max}$). Pour l'œil emmétrope jeune, $d_{\text{PP}} \approx \SI{25}{\centi\meter}$.
\item Le \textbf{domaine de vision nette} s'étend du PP au PR.
\item L'\textbf{amplitude d'accommodation} $A = C_{\max} - C_{\min}$ (en dioptries) diminue avec l'âge.
\item Outil de calcul : la formule de conjugaison $\dfrac{1}{\overline{OA'}} - \dfrac{1}{\overline{OA}} = C$, avec $\overline{OA'} = +\SI{15}{\milli\meter}$ et $\overline{OA} < 0$.
\end{itemize}
\end{aretenir}

Nous savons désormais décrire le fonctionnement de l'œil \emph{normal}. Mais que se passe-t-il lorsque le globe oculaire est trop long, ou la vergence au repos trop forte ? L'image des objets lointains se forme alors \emph{devant} la rétine : c'est la \textbf{myopie}, premier défaut d'accommodation que nous allons étudier et corriger dans la section suivante.

\coursec{La myopie et sa correction}

\courssub{Transition : quand l'accommodation ne suffit plus}

Dans la section précédente, nous avons vu qu'un œil \emph{emmélope} (normal) possède un \cle{punctum remotum} (PR) à l'infini et un \cle{punctum proximum} (PP) à une distance finie (environ \SI{25}{cm} pour un jeune adulte). L'accommodation, variation de la vergence du cristallin, permet de faire coïncider l'image nette d'un objet situé entre le PP et l'infini avec la rétine.

Mais que se passe-t-il si, même au \emph{repos complet} (accommodation minimale, cristallin le moins convergent possible), l'image d'un objet très lointain se forme \emph{devant} la rétine ? L'œil est alors \textbf{trop convergent} par rapport à sa longueur axiale, ou son globe est \textbf{trop long} par rapport à sa puissance optique. C'est le défaut le plus fréquent au Cameroun comme ailleurs : la \textbf{myopie}. Étudions-le en détail, de sa description physique à sa correction mathématique rigoureuse.

\begin{popfigure}
\centering
\begin{tikzpicture}[scale=1.1, >=Stealth]
% Couleurs
\colorlet{retina}{poppink!80!popdark}
\colorlet{lens}{popblue!60!popdark}
\colorlet{ray}{poporange!80!popdark}
\colorlet{blur}{popmuted}
% Œil : globe
\draw[popdark, thick] (0,0) circle (2.5);
% Cornée + cristallin (lentille mince équivalente)
\draw[lens, thick, fill=lens!10] (-0.3,-1.2) -- (-0.3,1.2) node[midway, right=2pt] {\small Lentille oculaire $L_{\text{œil}}$};
% Rétine
\draw[retina, very thick] (2.0,-1.5) -- (2.0,1.5) node[midway, right=3pt] {\small Rétine};
% Nerf optique
\draw[popdark, thick] (-2.5,0) -- (-3.5,0) node[left] {\small Nerf optique};
% Objet à l'infini : rayons parallèles incidents
\foreach \y in {-1.0,-0.5,0,0.5,1.0} {
    \draw[ray, thick, ->] (-5,\y) -- (-2.5,\y);
}
% Réfraction par la lentille oculaire : convergence trop forte
\foreach \y in {-1.0,-0.5,0,0.5,1.0} {
    \draw[ray, thick] (-0.3,\y) -- (1.0,0); % Convergence avant la rétine
    \draw[ray, dashed, blur] (1.0,0) -- (2.0,\y/2.5); % Divergence après le foyer -> tache sur rétine
}
% Foyer image de l'œil seul (F')
\fill[ray] (1.0,0) circle (1.5pt) node[below=2pt] {$F'_{\text{œil}}$};
% Tache floue sur la rétine
\draw[blur, thick, fill=blur!20] (2.0,-0.4) rectangle (2.0,0.4) node[midway, right=3pt] {\small Tache floue};
% Légendes
\node[align=center, popdark] at (0,-3.0) {\textbf{Œil myope non corrigé} : image d'un objet à l'infini\\formée en avant de la rétine (tache floue)};
\end{tikzpicture}
\legende{Œil myope : la vergence de l'œil au repos est trop grande, le foyer image $F'_{\text{œil}}$ se situe devant la rétine. Les rayons convergent puis divergent, formant une tache floue sur la rétine.}
\end{popfigure}

\courssub{Définition et conséquences optiques de la myopie}

\begin{definition}[Myopie]
Un œil est dit \textbf{myope} lorsque son \cle{punctum remotum} (PR) se trouve à une \textbf{distance finie} devant l'œil. L'œil myope voit net les objets proches (entre son PP et son PR) mais \textbf{flou les objets lointains} (au-delà de son PR).
\end{definition}

Physiquement, deux causes anatomiques principales (souvent combinées) expliquent ce défaut :
\begin{itemize}
    \item \textbf{Myopie axile :} le globe oculaire est trop long (distance cornée-rétine excessive) pour la puissance optique de l'œil.
    \item \textbf{Myopie d'indice (ou de courbure) :} la cornée et/ou le cristallin sont trop convergents (courbures trop prononcées ou indices trop élevés) pour la longueur axiale normale.
\end{itemize}
Dans les deux cas, la \textbf{vergence de l'œil au repos} $C_{\text{œil, repos}}$ est \textbf{trop forte}. L'image d'un objet à l'infini (rayons incidents parallèles) se forme en un point $F'_{\text{œil}}$ situé \emph{entre} la lentille oculaire et la rétine.

\courssub{Le punctum remotum (PR) et le punctum proximum (PP) de l'œil myope}

Puisque l'œil au repos forme l'image nette d'un objet à l'infini \emph{en avant} de la rétine, il existe une distance finie $d_{\text{PR}}$ pour laquelle l'objet placé à cette distance donne une image \emph{exactement sur la rétine} sans accommodation. C'est la définition même du \cle{punctum remotum} (PR).

\begin{propriete}[Position du PR et du PP chez le myope]
\begin{itemize}
    \item Le \textbf{PR} est à une distance finie $d_{\text{PR}}$ (ex. \SI{50}{cm}, \SI{1}{m}, \SI{2}{m}) \textbf{devant} l'œil. C'est la \textbf{vision la plus lointaine nette sans correction}.
    \item Le \textbf{PP} (vision la plus proche nette à accommodation maximale) est \textbf{plus proche} de l'œil que chez l'emmétrope. Le myope peut lire de très près (ex. \SI{10}{cm}) sans effort, ce qui est souvent perçu comme un « avantage » de près.
\end{itemize}
\end{propriete}

\begin{savaistu}[Myopie et scolarité au Cameroun]
Dans nos lycées de Yaoundé, Douala ou Bafoussam, on observe souvent des élèves qui plissent les yeux pour lire le tableau (vision de loin floue) mais qui lisent leur cahier de très près sans lunettes. C'est le signe typique de la myopie. Le dépistage précoce (dès le primaire) est crucial : une myopie non corrigée fatigue l'élève, provoque des maux de tête et peut nuire à la réussite scolaire. Les campagnes de dépistage organisées par le MINSANTE et des ONG dans les zones rurales (ex. région de l'Extrême-Nord) permettent de distribuer des lunettes à bas coût.
\end{savaistu}

\courssub{Principe de la correction : rendre l'œil moins convergent}

L'objectif est de permettre au myope de voir net un objet à l'infini (étoiles, tableau, panneau de signalisation sur l'autoroute Yaoundé-Douala) \textbf{sans accommodation} (œil au repos). Pour cela, il faut que l'image de cet objet à l'infini, après passage à travers le système correcteur + œil, se forme \textbf{sur la rétine}.

Puisque l'œil seul forme cette image \emph{devant} la rétine, il faut \textbf{diminuer la convergence totale} du système (correcteur + œil). On place donc devant l'œil une \textbf{lentille divergente} (vergence négative, bords plus minces que le centre).

\begin{popfigure}
\centering
\begin{tikzpicture}[scale=1.1, >=Stealth]
\colorlet{retina}{poppink!80!popdark}
\colorlet{lens}{popblue!60!popdark}
\colorlet{corr}{popgreen!70!popdark}
\colorlet{ray}{poporange!80!popdark}
% Œil
\draw[popdark, thick] (0,0) circle (2.5);
\draw[lens, thick, fill=lens!10] (-0.3,-1.2) -- (-0.3,1.2) node[midway, right=2pt] {\small $L_{\text{œil}}$};
\draw[retina, very thick] (2.0,-1.5) -- (2.0,1.5) node[midway, right=3pt] {\small Rétine};
\draw[popdark, thick] (-2.5,0) -- (-3.5,0) node[left] {\small Nerf optique};
% Lentille correctrice (divergente)
\draw[corr, thick, fill=corr!10] (-3.0,-1.0) -- (-3.0,1.0) node[midway, left=2pt] {\small Verre correcteur $L_c$ (div.)};
% Rayons incidents parallèles (objet à l'infini)
\foreach \y in {-0.8,-0.4,0,0.4,0.8} {
    \draw[ray, thick, ->] (-5.5,\y) -- (-3.0,\y);
}
% Réfraction par lentille divergente : rayons divergent comme s'ils venaient de F'_c (foyer image de la lentille correctrice)
\foreach \y in {-0.8,-0.4,0,0.4,0.8} {
    % Prolongement arrière pour montrer image virtuelle
    \draw[ray, dashed, thin] (-3.0,\y) -- (-4.0,0); 
    % Rayons réels après lentille correctrice
    \draw[ray, thick] (-3.0,\y) -- (-0.3,\y*0.6); % Légère divergence
}
% Foyer image de la lentille correctrice (virtuel, côté objet)
\fill[corr] (-4.0,0) circle (1.5pt) node[below=2pt] {$F'_c$};
% Réfraction par l'œil : convergence sur la rétine
\foreach \y in {-0.8,-0.4,0,0.4,0.8} {
    \draw[ray, thick] (-0.3,\y*0.6) -- (2.0,\y*0.1); % Ajusté pour converger sur rétine
}
% Point sur rétine
\fill[ray] (2.0,0) circle (1.5pt);
\node[align=center, popdark] at (0,-3.0) {\textbf{Œil myope corrigé} : lentille divergente\\écarte les rayons, image nette sur la rétine};
\end{tikzpicture}
\legende{Correction de la myopie : la lentille divergente $L_c$ crée une image virtuelle de l'objet à l'infini au niveau du PR de l'œil. Cette image virtuelle sert d'objet réel pour l'œil qui la met au point sur la rétine.}
\end{popfigure}

\courssub{Condition de correction rigoureuse : la formule de conjugaison}

Plaçons la lentille correctrice \textbf{accolée} à l'œil (hypothèse standard en Première : distance verre-œil négligeable, \SI{\approx 0}{cm}). Notons :
\begin{itemize}
    \item $C_c$ : vergence de la lentille correctrice (en dioptries $\delta$ ou $\text{m}^{-1}$). On cherche $C_c < 0$.
    \item $f'_c = 1/C_c$ : foyer image de la lentille correctrice.
    \item $d_{\text{PR}}$ : distance du punctum remotum de l'œil nu (en mètres, $d_{\text{PR}} > 0$).
\end{itemize}

\textbf{Raisonnement physique (guidé) :}
\begin{enumerate}
    \item L'objet est à l'infini : $\overline{OA} = -\infty$ (distance algébrique objet/lentille).
    \item La lentille correctrice doit donner une \textbf{image virtuelle} située \textbf{au PR de l'œil} pour que l'œil, sans accommodation, la mette au point sur la rétine.
    \item Cette image virtuelle est du \emph{même côté} que l'objet par rapport à la lentille divergente : sa distance algébrique $\overline{A'B'} = -d_{\text{PR}}$ (négative car image virtuelle, côté objet).
    \item On applique la \textbf{formule de conjugaison de Newton} (ou de Descartes) pour la lentille correctrice seule.
\end{enumerate}

\begin{propriete}[Vergence de la lentille correctrice de la myopie — Démonstration]
\textbf{Énoncé :} Pour corriger une myopie dont le punctum remotum est à la distance $d_{\text{PR}}$ (en mètres) de l'œil, la lentille correctrice (accolée à l'œil) doit être \textbf{divergente} de vergence :
\[
\boxed{C_c = -\dfrac{1}{d_{\text{PR}}}}
\]
\textbf{Démonstration (exigible au Probatoire) :}
On utilise la relation de Newton pour la lentille correctrice : $\overline{FA} \cdot \overline{F'A'} = -f'^2$.
\begin{itemize}
    \item Objet à l'infini $\Rightarrow$ $\overline{FA} = -\infty$ (ou on utilise Descartes : $\frac{1}{\overline{A'B'}} - \frac{1}{\overline{AB}} = C_c$).
    \item Image virtuelle au PR : $\overline{A'B'} = -d_{\text{PR}}$.
    \item Distance objet $\overline{AB} = -\infty \Rightarrow \frac{1}{\overline{AB}} = 0$.
\end{itemize}
Formule de Descartes (conjugaison) :
\[
\frac{1}{\overline{A'B'}} - \frac{1}{\overline{AB}} = C_c
\]
\[
\frac{1}{-d_{\text{PR}}} - 0 = C_c \quad \Rightarrow \quad C_c = -\frac{1}{d_{\text{PR}}}
\]
\emph{Analyse dimensionnelle :} $d_{\text{PR}}$ en $\text{m} \Rightarrow 1/d_{\text{PR}}$ en $\text{m}^{-1} = \delta$. Le signe moins confirme une lentille \textbf{divergente}. \qed
\end{propriete}

\begin{attention}[Pièges classiques au Probatoire]
\begin{enumerate}
    \item \textbf{Signe de la vergence :} Une lentille divergente a \textbf{obligatoirement une vergence négative} ($C < 0$). Oublier le signe « $-$ » fait perdre la moitié des points (ou plus).
    \item \textbf{Conversion des unités :} $d_{\text{PR}}$ est souvent donné en \textbf{centimètres} (ex. \SI{40}{cm}). Il \textbf{faut impérativement convertir en mètres} (\SI{0,40}{m}) avant de calculer $C = -1/d_{\text{PR}}$. Erreur fréquente : $C = -1/40 = -0,025\,\delta$ (faux) au lieu de $-1/0,40 = -2,5\,\delta$ (vrai).
    \item \textbf{Distance verre-œil :} En Première, on suppose la lentille \textbf{accolée} à l'œil ($d=0$). En Terminale (ou si précisé), on tient compte de la distance vertex (ex. \SI{1,5}{cm}) ce qui modifie légèrement la vergence prescrite.
\end{enumerate}
\end{attention}

\courssub{Méthode de résolution type}

\begin{methode}[Corriger une myopie — 3 étapes]
\begin{enumerate}
    \item \textbf{Repérer le PR :} Lire l'énoncé : « Le punctum remotum est à $x$ cm » ou « Le sujet voit net jusqu'à $x$ cm sans lunettes ». Convertir en mètres : $d_{\text{PR}} = x/100$.
    \item \textbf{Imposer la condition de correction :} L'image de l'objet à l'infini par le verre correcteur doit se former au PR. $\Rightarrow \overline{A'B'} = -d_{\text{PR}}$.
    \item \textbf{Calculer la vergence :} $C_c = -1/d_{\text{PR}}$ (en $\delta$). Vérifier : $C_c < 0$ (divergente). La puissance commerciale s'arrondit souvent au quart de dioptrie près (ex. $-2,50\,\delta$, $-2,75\,\delta$).
\end{enumerate}
\end{methode}

\courssub{Exemple chiffré détaillé}

\begin{exoresolu}[Correction de la myopie de M. Fotso]
\textbf{Énoncé :} M. Fotso, enseignant à Bertoua, consulte un ophtalmologiste. Celui-ci détermine que son punctum remotum est situé à \textbf{\SI{40}{cm}} de son œil. On suppose la lentille correctrice accolée à l'œil. Quelle est la vergence des verres à prescrire ? Quel est leur foyer image ?
\tcblower
\textbf{Solution détaillée :}
\begin{enumerate}
    \item \textbf{Données :} $d_{\text{PR}} = \SI{40}{cm} = \SI{0,40}{m}$.
    \item \textbf{Inconnues :} $C_c$ (en $\delta$), $f'_c$ (en m).
    \item \textbf{Formule :} $C_c = -\dfrac{1}{d_{\text{PR}}}$.
    \item \textbf{Calcul :}
    \[
    C_c = -\frac{1}{0,40} = -2,5\,\delta
    \]
    \item \textbf{Foyer image :} $f'_c = \dfrac{1}{C_c} = \dfrac{1}{-2,5} = -\SI{0,40}{m} = -\SI{40}{cm}$.
    \item \textbf{Interprétation :} Il faut prescrire des verres \textbf{divergents} de \textbf{-2,50 dioptries}. Leur foyer image $F'_c$ est virtuel, situé à \SI{40}{cm} du verre côté objet. C'est exactement la position du PR de M. Fotso.
    \item \textbf{Vérification dimensionnelle :} $[C] = \text{m}^{-1} = \delta$. Cohérent.
\end{enumerate}
\textbf{Réponse :} Verres divergents de vergence $C = \mathbf{-2,5\,\delta}$ (foyer image $f' = \mathbf{-40\,cm}$).
\end{exoresolu}

\courssub{Effet de la correction sur le punctum proximum (PP)}

La correction éloigne l'image virtuelle des objets proches. Le PP \textbf{reculé} (s'éloigne de l'œil).
\begin{itemize}
    \item Sans correction : PP très proche (ex. \SI{10}{cm}). Lecture très aisée de près.
    \item Avec correction : Le PP apparent (à travers les lunettes) se trouve plus loin. Pour un myope fort (ex. $-6\,\delta$), le PP peut reculer jusqu'à \SI{30}{cm} ou \SI{40}{cm}, ce qui peut gêner la lecture de très près (petits caractères, téléphone portable).
\end{itemize}
C'est un compromis acceptable pour la vision de loin. Chez le \textbf{presbyte myope} (sujet âgé), on utilise souvent des verres progressifs ou on retire les lunettes pour lire de près.

\begin{popfigure}
\centering
\begin{tikzpicture}[scale=1.2, >=Stealth]
% Construction géométrique : Objet à l'infini -> Lentille divergente -> Image virtuelle au PR
\colorlet{corr}{popgreen!70!popdark}
\colorlet{ray}{poporange!80!popdark}
\colorlet{obj}{popblue!70!popdark}
\colorlet{img}{poppink!70!popdark}
% Axes
\draw[popline, thick, ->] (-5,0) -- (3,0) node[right] {Axe optique};
% Lentille divergente
\draw[corr, very thick] (0,-1.2) -- (0,1.2) node[above=2pt] {$L_c$ (div.)};
% Foyers de la lentille correctrice
\fill[corr] (-1.5,0) circle (2pt) node[below=3pt] {$F_c$ (objet)};
\fill[corr] (1.5,0) circle (2pt) node[below=3pt] {$F'_c$ (image)};
% Objet à l'infini : rayon incident parallèle (hauteur h)
\def\h{1.0}
\draw[ray, thick, ->] (-4.5,\h) -- (0,\h) node[midway, above] {Objet à l'infini};
% Rayon 1 : parallèle -> passe par F'_c (virtuel, prolongation arrière)
\draw[ray, thick] (0,\h) -- (1.5,0);
\draw[ray, dashed, thin] (0,\h) -- (-1.5,0); % Prolongement pour construction
% Rayon 2 : passe par centre optique O (non dévié)
\draw[ray, thick, ->] (-4.5,0) -- (0,0) -- (2.5,0);
% Image virtuelle A'B' : intersection des prolongements arrière
\draw[img, very thick, dashed] (-1.5,-0.5) -- (-1.5,0.5) node[above=2pt] {$A'B'$ (Image virtuelle au PR)};
% Flèches objet/image
\draw[obj, <->] (-4.5,0.2) -- (-4.5,\h) node[midway, left] {$AB$ (à l'infini)};
\draw[img, <->] (-1.5,0) -- (-1.5,0.5) node[midway, left] {$A'B'$};
% Cotes
\draw[popmuted, <->] (0,-1.5) -- (-1.5,-1.5) node[midway, below] {$f'_c = -d_{\text{PR}}$};
\draw[popmuted, <->] (-1.5,-1.5) -- (-4.5,-1.5) node[midway, below] {Infini};
% Légende
\node[align=center, popdark] at (0,-2.5) {\textbf{Construction géométrique} : Objet à l'infini $\to$ Image virtuelle au PR par lentille divergente};
\end{tikzpicture}
\legende{Construction de l'image virtuelle d'un objet à l'infini par la lentille divergente correctrice. L'image $A'B'$ se forme au foyer image $F'_c$ de la lentille, qui doit coïncider avec le PR de l'œil myope.}
\end{popfigure}

\courssub{Synthèse de la section : Myopie}

\begin{aretenir}[Bilan — Myopie et correction]
\begin{itemize}
    \item \textbf{Défaut :} Œil trop convergent ou globe trop long $\Rightarrow$ Image de l'infini devant la rétine.
    \item \textbf{PR :} Distance finie $d_{\text{PR}}$ (ex. \SI{50}{cm}). Vision nette de loin \textbf{impossible} sans correction.
    \item \textbf{PP :} Plus proche que la normale. Vision de près \textbf{excellente} sans lunettes.
    \item \textbf{Correction :} Lentille \textbf{divergente} (vergence $C < 0$) accolée à l'œil.
    \item \textbf{Formule clé :} $\boxed{C = -\dfrac{1}{d_{\text{PR}}}}$ ($d_{\text{PR}}$ en mètres, $C$ en dioptries).
    \item \textbf{Effet sur PP :} La correction recule le PP (vision de près légèrement moins performante, mais confortable pour la plupart).
\end{itemize}
\end{aretenir}

\begin{exercice}[Application directe]
\begin{enumerate}
    \item Un élève de Première C à Garoua a son punctum remotum à \SI{60}{cm}. Calculer la vergence des verres correcteurs (accolés à l'œil) qu'il doit porter pour voir net le tableau.
    \item Une patiente porte des verres de vergence $-3,00\,\delta$. Quelle est la position de son punctum remotum sans lunettes ?
    \item \textbf{Erreur à corriger :} Un étudiant écrit : « $d_{\text{PR}} = 80\,\text{cm}$, donc $C = -1/80 = -0,0125\,\delta$ ». Expliquez l'erreur et donnez le bon résultat.
\end{enumerate}
\end{exercice}

\begin{corrige}[Exercice n]
\begin{enumerate}
    \item $d_{\text{PR}} = \SI{60}{cm} = \SI{0,60}{m}$. $C = -1/0,60 = \mathbf{-1,67\,\delta}$ (verres divergents).
    \item $C = -3,00\,\delta \Rightarrow d_{\text{PR}} = -1/C = 1/3,00 = \mathbf{0,33\,m = 33\,cm}$.
    \item Erreur : non-conversion des cm en m. $80\,\text{cm} = 0,80\,\text{m}$. $C = -1/0,80 = \mathbf{-1,25\,\delta}$.
\end{enumerate}
\end{corrige}

Cette section établit les bases rigoureuses de la correction de la myopie. La section suivante traitera de l'hypermétropie, défaut symétrique où l'œil est \emph{pas assez} convergent, nécessitant une lentille \emph{convergente} pour rapprocher le PR virtuel de l'infini.

\coursec{L'hypermétropie et sa correction}

Dans la section précédente, nous avons étudié la myopie : un œil \emph{trop convergent}, dont le globe est trop long, et qui voit net de près mais flou de loin. Nous allons maintenant rencontrer le défaut exactement symétrique : l'\cle{hypermétropie}. Tu verras que la démarche est identique — comprendre le défaut géométriquement, repérer les points particuliers PP et PR, puis choisir la lentille qui ramène l'image là où l'œil peut la voir nette — mais que tous les signes sont inversés. C'est une excellente occasion de consolider ta maîtrise de la formule de conjugaison des lentilles minces.

\courssub{Observation et définition de l'hypermétropie}

\textbf{Une situation concrète.}\par

Merveille, 16 ans, est une excellente élève. Pourtant, depuis quelques mois, elle se plaint de maux de tête en fin de journée, surtout après les séances de lecture. Curieusement, elle voit très bien le tableau au fond de la salle, mais elle doit tenir son cahier à bout de bras pour le lire, et même ainsi, les lettres finissent par devenir floues. Sa grand-mère fait la même chose... mais pour une raison différente, comme nous le verrons à la section suivante. Merveille, elle, est probablement \cle{hypermétrope}.

\textbf{Analyse géométrique du défaut.}\par

Rappelons le fonctionnement de l'œil normal au repos : les rayons issus d'un objet très éloigné (à l'infini) arrivent parallèles à l'axe optique, traversent le système cornée--cristallin, et convergent exactement \emph{sur la rétine}. L'hypermétropie apparaît dans deux situations équivalentes :

\begin{itemize}
\item le système optique de l'œil (cornée + cristallin) est \textbf{pas assez convergent} : sa vergence au repos est trop faible ;
\item ou bien le globe oculaire est \textbf{trop court} : la rétine est placée trop près du cristallin.
\end{itemize}

Dans les deux cas, la conséquence est la même : pour un œil hypermétrope \emph{au repos}, les rayons parallèles venant de l'infini convergent en un point situé \textbf{derrière la rétine}. Sur la rétine, on ne récolte donc qu'une tache floue.

\begin{popfigure}
\begin{center}
\begin{tikzpicture}[scale=1.0,>=Stealth]
% axe optique
\draw[popmuted] (-0.5,0) -- (6.5,0);
% cristallin (double flèche)
\draw[popblue,thick] (0.8,-1.3) -- (0.8,1.3);
\draw[popblue,thick,->] (0.8,1.3) -- (0.8,1.55);
\draw[popblue,thick,->] (0.8,-1.3) -- (0.8,-1.55);
\node[popblue,below] at (0.8,-1.6) {\small cristallin};
% rétine
\draw[popdark,very thick] (3.6,-1.3) -- (3.6,1.3);
\node[popdark,above,align=center] at (3.6,1.35) {\small rétine\\ \small (trop proche)};
% foyer image derrière la rétine
\fill[poporange] (5.2,0) circle (0.07);
\node[poporange,below,align=center] at (5.2,-0.12) {\small foyer image $F'$\\ \small (derrière la rétine)};
% rayons parallèles
\draw[popgreen,thick,->] (-0.4,0.9) -- (2.4,0.27);
\draw[popgreen,thick] (2.4,0.27) -- (5.2,0);
\draw[popgreen,thick,->] (-0.4,-0.9) -- (2.4,-0.27);
\draw[popgreen,thick] (2.4,-0.27) -- (5.2,0);
\draw[popgreen,thick,->] (-0.4,0) -- (2.6,0);
\draw[popgreen,thick] (2.6,0) -- (5.2,0);
\node[popgreen,left,align=right] at (-0.4,0.9) {\small rayons venant\\ \small de l'infini};
% tache floue sur la rétine
\fill[poporange!40] (3.6,0) ellipse (0.06 and 0.45);
\node[poporange,right,align=left] at (3.75,0.85) {\small tache floue};
\end{tikzpicture}
\end{center}
\legende{Œil hypermétrope au repos, non corrigé : l'image d'un objet à l'infini se formerait \emph{derrière} la rétine. La rétine ne reçoit qu'une tache floue.}
\end{popfigure}

\begin{definition}[L'hypermétropie]
Un œil \cle{hypermétrope} est un œil \textbf{insuffisamment convergent} (ou dont le globe oculaire est trop court). Au repos, l'image d'un objet éloigné se formerait \textbf{derrière la rétine}. Pour voir net, même de loin, l'œil hypermétrope doit \textbf{accommoder en permanence}, et sa vision de près est très limitée.
\end{definition}

\begin{attention}[Hypermétropie ou presbytie ?]
Le mot exact est \textbf{hypermétropie} (du grec \emph{hyper} : au-delà, et \emph{metron} : mesure). Évite aussi de confondre avec la \textbf{presbytie} : l'hypermétropie est un défaut de \emph{convergence} (forme du globe ou puissance du cristallin), présent \textbf{dès la jeunesse} ; la presbytie est une perte d'\emph{amplitude d'accommodation} liée au \textbf{vieillissement} du cristallin, qui apparaît vers 45 ans même chez un œil autrefois normal. Les deux se corrigent avec des lentilles convergentes, mais leurs causes sont différentes.
\end{attention}

\courssub{Les points particuliers de l'œil hypermétrope : PR et PP}

\textbf{Le punctum remotum : un point virtuel.}\par

Pour un œil normal, le \cle{punctum remotum} (PR) — le point le plus lointain vu net \emph{sans accommodation} — est rejeté à l'infini. Pour l'œil hypermétrope, la situation est subtile : même au repos, son œil est trop peu convergent pour focaliser sur la rétine des rayons parallèles. Pour que l'image tombe sur la rétine sans accommoder, il faudrait que l'objet envoie des rayons \emph{déjà convergents} vers l'œil, c'est-à-dire que l'objet soit... \textbf{virtuel, situé derrière l'œil}. On dit que le PR de l'œil hypermétrope est \textbf{virtuel}.

Retiens simplement ceci : l'œil hypermétrope \textbf{doit accommoder pour voir net à l'infini}. Cet effort permanent du muscle ciliaire explique les symptômes de Merveille : \textbf{fatigue visuelle, maux de tête, picotements} en fin de journée. L'hypermétropie modérée passe souvent inaperçue chez le jeune enfant précisément parce que son cristallin, très souple, compense sans cesse — au prix d'une fatigue chronique.

\textbf{Le punctum proximum : repoussé au loin.}\par

Le \cle{punctum proximum} (PP) est le point le plus proche vu net, avec accommodation maximale. Or l'œil hypermétrope « dépense » déjà une partie de son pouvoir d'accommodation pour la vision de loin : il ne lui en reste plus assez pour la vision de près. Son PP est donc \textbf{plus éloigné que la normale} : là où un œil normal lit confortablement à $d_m = \SI{25}{cm}$, un œil hypermétrope peut avoir son PP à $\SI{50}{cm}$, $\SI{1}{m}$ ou davantage. C'est pourquoi Merveille écarte son cahier à bout de bras.

\begin{aretenir}[Les deux signatures de l'hypermétropie]
\begin{itemize}
\item \textbf{PR virtuel} (derrière l'œil) : l'œil accommode même pour regarder à l'infini $\Rightarrow$ fatigue visuelle permanente ;
\item \textbf{PP trop éloigné} ($d_m > \SI{25}{cm}$) : la lecture de près est impossible ou pénible.
\end{itemize}
\end{aretenir}

\courssub{Principe de la correction : une lentille convergente}

Le défaut est un \emph{déficit de convergence}. La correction est donc logique : on \textbf{ajoute de la convergence} en plaçant devant l'œil une \textbf{lentille convergente} (vergence positive). C'est l'inverse de la myopie, que l'on corrigeait par une lentille divergente.

\begin{propriete}[Correction de l'hypermétropie]
L'hypermétropie se corrige par une \textbf{lentille convergente} (verre correcteur de \textbf{vergence positive}, $C > 0$). La lentille pré-converge les rayons lumineux, ce qui soulage l'accommodation de l'œil et ramène l'image sur la rétine.
\end{propriete}

\begin{popfigure}
\begin{center}
\begin{tikzpicture}[scale=1.0,>=Stealth]
% axe optique
\draw[popmuted] (-0.8,0) -- (5.6,0);
% lentille correctrice convergente
\draw[poppurple,thick] (-0.3,-1.3) -- (-0.3,1.3);
\draw[poppurple,thick,->] (-0.3,1.3) -- (-0.3,1.55);
\draw[poppurple,thick,->] (-0.3,-1.3) -- (-0.3,-1.55);
\node[poppurple,above] at (-0.3,1.6) {\small lentille convergente};
% cristallin
\draw[popblue,thick] (1.6,-1.3) -- (1.6,1.3);
\draw[popblue,thick,->] (1.6,1.3) -- (1.6,1.55);
\draw[popblue,thick,->] (1.6,-1.3) -- (1.6,-1.55);
\node[popblue,below] at (1.6,-1.6) {\small cristallin};
% rétine
\draw[popdark,very thick] (4.4,-1.3) -- (4.4,1.3);
\node[popdark,above] at (4.4,1.35) {\small rétine};
% rayons
\draw[popgreen,thick,->] (-0.7,0.95) -- (0.6,0.55);
\draw[popgreen,thick] (0.6,0.55) -- (4.4,0);
\draw[popgreen,thick,->] (-0.7,-0.95) -- (0.6,-0.55);
\draw[popgreen,thick] (0.6,-0.55) -- (4.4,0);
\draw[popgreen,thick,->] (-0.7,0) -- (2.2,0);
\draw[popgreen,thick] (2.2,0) -- (4.4,0);
% point image net
\fill[poporange] (4.4,0) circle (0.08);
\node[poporange,below] at (4.4,-0.15) {\small image nette sur la rétine};
\end{tikzpicture}
\end{center}
\legende{Œil hypermétrope corrigé : la lentille convergente fait démarrer la convergence plus tôt ; l'image se forme exactement sur la rétine.}
\end{popfigure}

\textbf{Condition de correction usuelle.}\par

En pratique, l'ophtalmologiste cherche à rendre à l'œil hypermétrope une \textbf{vision de près confortable} : il impose que l'œil corrigé puisse lire à la distance normale $d_m = \SI{25}{cm}$ \emph{avec son accommodation maximale habituelle}. Traduite en langage optique, la condition s'écrit :

\begin{aretenir}[Condition de correction de l'hypermétropie]
L'objet réel placé au PP normal ($\overline{OA} = -\SI{25}{cm}$) doit, à travers la lentille correctrice, donner une \textbf{image virtuelle située au PP réel de l'œil} ($\overline{OA'} = -d_m^{\text{œil}}$). L'œil voit alors cette image virtuelle nette en accommodant au maximum, comme s'il regardait son propre PP.
\end{aretenir}

Remarque bien les signes : l'image est \emph{virtuelle}, située \emph{avant} la lentille (du côté de l'objet), donc $\overline{OA'} < 0$. C'est le point le plus délicat du calcul.

\begin{methode}[Calculer la vergence correctrice d'un œil hypermétrope]
\begin{enumerate}
\item \textbf{Choisir la position d'objet souhaitée} : en général la lecture au PP normal, $\overline{OA} = -\SI{0,25}{m}$.
\item \textbf{Imposer la position de l'image} : l'image virtuelle doit se former au PP réel de l'œil, $\overline{OA'} = -d_m^{\text{œil}}$ (signe négatif car image virtuelle, en amont de la lentille).
\item \textbf{Appliquer la formule de conjugaison} des lentilles minces :
\[ \frac{1}{\overline{OA'}} - \frac{1}{\overline{OA}} = \frac{1}{f'} = C \]
en exprimant toutes les distances en \textbf{mètres} et en respectant scrupuleusement les signes.
\item \textbf{Conclure} : on doit trouver $C > 0$ (lentille convergente). Si tu trouves une vergence négative, tu as fait une erreur de signe !
\end{enumerate}
\end{methode}

\courssub{Exemple chiffré complet}

\begin{exoresolu}[Correction d'un œil hypermétrope dont le PP est à 1 m]
Un œil hypermétrope a son punctum proximum à la distance $d_m = \SI{1,0}{m}$. On souhaite lui permettre de lire un texte placé à $\SI{25}{cm}$ de l'œil, en accommodant au maximum. La lentille correctrice est placée au voisinage immédiat de l'œil (on confond le centre optique $O$ de la lentille avec celui de l'œil réduit).
\begin{enumerate}
\item Déterminer la nature de la lentille correctrice.
\item Calculer sa vergence $C$ et sa distance focale $f'$.
\end{enumerate}
\tcblower
\textbf{1. Nature de la lentille.} L'œil hypermétrope n'est pas assez convergent : il faut renforcer sa convergence. La lentille correctrice est donc \textbf{convergente}, de vergence $C > 0$.

\textbf{2. Vergence correctrice.} Appliquons la méthode :

\emph{Étape 1 — Objet.} Le texte est placé à $\SI{25}{cm}$ devant la lentille : c'est un objet réel, donc
\[ \overline{OA} = -\SI{0,25}{m}. \]

\emph{Étape 2 — Image.} L'œil ne voit net de près qu'à partir de $\SI{1,0}{m}$. La lentille doit donc donner du texte une image \textbf{virtuelle} située à $\SI{1,0}{m}$ devant la lentille (c'est cette image que l'œil regardera) :
\[ \overline{OA'} = -\SI{1,0}{m}. \]

\emph{Étape 3 — Formule de conjugaison.}
\begin{align*}
C &= \frac{1}{\overline{OA'}} - \frac{1}{\overline{OA}} \\
&= \frac{1}{-1{,}0} - \frac{1}{-0{,}25} \\
&= -1 + 4 = +3\ \delta.
\end{align*}

\emph{Étape 4 — Distance focale.}
\[ f' = \frac{1}{C} = \frac{1}{3} \approx \SI{0,33}{m} = \SI{33}{cm}. \]

\textbf{Conclusion.} Il faut prescrire une lentille \textbf{convergente de vergence $C = +3\ \delta$} (distance focale $f' \approx \SI{33}{cm}$). Le signe positif confirme la nature convergente : le résultat est cohérent. Avec ces verres, le texte placé à $\SI{25}{cm}$ donne une image virtuelle à $\SI{1}{m}$, exactement là où l'œil de Merveille sait voir net de près.
\end{exoresolu}

\begin{attention}[Les trois pièges classiques du calcul]
\begin{itemize}
\item \textbf{Oublier le signe de $\overline{OA'}$.} L'image est \emph{virtuelle} : elle se trouve du même côté que l'objet, donc $\overline{OA'} = -1{,}0\ \text{m}$ et non $+1{,}0\ \text{m}$. Avec le mauvais signe, tu trouverais $C = \dfrac{1}{1{,}0} - \dfrac{1}{-0{,}25} = 1 + 4 = +5\ \delta$ : une vergence fausse, trop forte. Écris toujours les deux fractions explicitement avant de calculer.
\item \textbf{Oublier de convertir en mètres.} La vergence s'exprime en dioptries ($\delta = \text{m}^{-1}$) : $\SI{25}{cm} = \SI{0,25}{m}$ obligatoirement.
\item \textbf{Confondre PP et PR.} Pour l'hypermétropie, la correction usuelle vise la \emph{vision de près} : c'est le \textbf{PP} qui intervient dans le calcul, alors que pour la myopie c'était le \textbf{PR}. Lis bien l'énoncé pour savoir quelle distance utiliser.
\end{itemize}
\end{attention}

\begin{attention}[Analyse dimensionnelle de la formule de conjugaison]
Chaque terme de $\dfrac{1}{\overline{OA'}} - \dfrac{1}{\overline{OA}} = C$ est l'inverse d'une longueur, donc s'exprime en $\text{m}^{-1}$. La dioptrie est bien homogène : $1\ \delta = 1\ \text{m}^{-1}$. Si ton calcul fait apparaître des mètres au numérateur, c'est qu'une inversion a été oubliée.
\end{attention}

\courssub{Comparaison avec la myopie : tableau récapitulatif}

Les deux défauts d'accommodation étudiés jusqu'ici sont rigoureusement symétriques. Ce tableau, à savoir reproduire le jour du Probatoire, en fait la synthèse.

\begin{center}
\begin{tabularx}{\linewidth}{|l|X|X|X|}
\hline
\rowcolor{popblueL} \textbf{Caractéristique} & \textbf{Œil normal (emmétrope)} & \textbf{Œil myope} & \textbf{Œil hypermétrope} \\
\hline
Convergence & normale & trop convergent (globe trop long) & pas assez convergent (globe trop court) \\
\hline
Image d'un objet à l'infini (œil au repos) & sur la rétine & en avant de la rétine & derrière la rétine \\
\hline
Punctum remotum (PR) & à l'infini & réel, rapproché (ex. $\SI{1}{m}$ devant l'œil) & virtuel, derrière l'œil (accommode pour voir loin) \\
\hline
Punctum proximum (PP) & à $\SI{25}{cm}$ & très rapproché ($< \SI{25}{cm}$) & trop éloigné ($> \SI{25}{cm}$) \\
\hline
Vision & nette de près et de loin & bonne de près, floue de loin & fatigante de loin, mauvaise de près \\
\hline
Verre correcteur & aucun & lentille \textbf{divergente} ($C < 0$) & lentille \textbf{convergente} ($C > 0$) \\
\hline
Rôle de la correction & --- & ramener le PR à l'infini & ramener le PP à $\SI{25}{cm}$ \\
\hline
\end{tabularx}
\end{center}

\begin{savaistu}[L'hypermétropie au quotidien]
Chez le jeune enfant, une hypermétropie modérée est fréquente et souvent compensée sans effort apparent grâce à la souplesse du cristallin ; elle peut cependant provoquer un strabisme accommodatif, que le port précoce de lunettes convergentes corrige remarquablement. Au Cameroun comme ailleurs, les campagnes de dépistage visuel en milieu scolaire permettent de repérer ces défauts avant qu'ils ne pénalisent les apprentissages : un enfant qui « n'aime pas lire » souffre parfois simplement d'une hypermétropie non corrigée. Les verres convergents, épais au centre, donnent aux yeux de leur porteur un aspect légèrement grossi — à l'inverse des verres divergents des myopes, épais au bord, qui rapetissent l'œil vu de face.
\end{savaistu}

\begin{exercice}[À toi de jouer]
Un élève hypermétrope a son punctum proximum à $d_m = \SI{50}{cm}$. On veut lui permettre de lire à $\SI{25}{cm}$ en accommodant au maximum, avec une lentille accolée à l'œil.
\begin{enumerate}
\item Quelle doit être la position de l'image donnée par la lentille ? Préciser sa nature.
\item Calculer la vergence $C$ de la lentille correctrice et en déduire sa distance focale.
\item Sans refaire le calcul complet, dire qualitativement comment évoluerait la vergence correctrice si le PP de l'élève était à $\SI{2}{m}$ au lieu de $\SI{50}{cm}$.
\end{enumerate}
\end{exercice}

\begin{corrige}[Exercice : correction d'un œil hypermétrope (PP à 50 cm)]
\textbf{1.} L'image doit se former au \textbf{PP réel de l'œil}, c'est-à-dire à $\SI{50}{cm}$ devant la lentille. Elle est \textbf{virtuelle} (l'objet réel est plus proche que le foyer objet de la lentille convergente), donc $\overline{OA'} = -\SI{0,50}{m}$.

\textbf{2.} Formule de conjugaison avec $\overline{OA} = -\SI{0,25}{m}$ :
\begin{align*}
C &= \frac{1}{\overline{OA'}} - \frac{1}{\overline{OA}} = \frac{1}{-0{,}50} - \frac{1}{-0{,}25} \\
&= -2 + 4 = +2\ \delta,
\qquad f' = \frac{1}{C} = \SI{0,50}{m} = \SI{50}{cm}.
\end{align*}
La lentille est convergente, de vergence $+2\ \delta$ : résultat cohérent avec le défaut à corriger.

\textbf{3.} Si le PP passe à $\SI{2}{m}$, le terme $\dfrac{1}{\overline{OA'}} = -\dfrac{1}{2} = -0{,}5\ \delta$ devient plus petit en valeur absolue, donc $C = 4 - 0{,}5 = +3{,}5\ \delta$ augmente : plus l'hypermétropie est marquée (PP plus lointain), plus la vergence correctrice doit être \textbf{forte}. C'est physiquement logique : l'œil a besoin de davantage de convergence supplémentaire.
\end{corrige}

\begin{aretenir}[L'essentiel de la section]
\begin{itemize}
\item L'œil \textbf{hypermétrope} n'est \textbf{pas assez convergent} (globe trop court) : au repos, l'image d'un objet lointain se formerait \textbf{derrière la rétine}.
\item Il \textbf{accommode en permanence}, même pour voir loin (PR virtuel) : fatigue visuelle et maux de tête. Son \textbf{PP est trop éloigné} : il voit mal de près.
\item Correction : lentille \textbf{convergente} ($C > 0$) qui ramène le PP à $\SI{25}{cm}$.
\item Calcul type : $\overline{OA} = -\SI{0,25}{m}$, $\overline{OA'} = -d_m^{\text{œil}}$ (image virtuelle, signe négatif !), puis $C = \dfrac{1}{\overline{OA'}} - \dfrac{1}{\overline{OA}}$.
\item Ne pas confondre hypermétropie (défaut de convergence, dès la jeunesse) et presbytie (vieillissement du cristallin) — objet de la prochaine section.
\end{itemize}
\end{aretenir}

\coursec{La presbytie et bilan des corrections}

Après avoir étudié la myopie et l'hypermétropie — deux défauts liés à une inadéquation entre la puissance optique de l'œil et sa longueur axiale — nous abordons maintenant un troisième défaut, fondamentalement différent : la \cle{presbytie}. Contrairement aux deux précédents, elle n'est pas une erreur de « fabrication » de l'œil (forme ou longueur), mais une conséquence inéluctable du \emph{vieillissement} du cristallin. Elle touche \emph{tout le monde} à partir de 40--45 ans, y compris les anciens myopes et hypermétropes. Comprendre la presbytie, c'est comprendre comment l'amplitude d'accommodation s'effondre et comment on peut la compenser optiquement.

\courssub{Origine physiologique : le cristallin durcit}

\emph{Observation.} Un enfant de 10 ans peut lire un livre tenu à 7 cm de son nez sans effort. À 20 ans, la distance de lecture confortable est d'environ 15 cm. À 45 ans, on tend le bras (40--50 cm). À 60 ans, le bras n'est plus assez long : il faut des lunettes.

Cette évolution reflète la perte progressive d'élasticité du cristallin. Rappelons le mécanisme d'accommodation (section 2) : le muscle ciliaire se contracte $\rightarrow$ le ligament de Zinn se détend $\rightarrow$ le cristallin, élastique, \emph{bombe} (ses rayons de courbure diminuent) $\rightarrow$ sa vergence \textbf{augmente} pour faire le point sur un objet proche.

Avec l'âge, les fibres du cristallin se compactent, le noyau durcit (sclérose nucléaire). Le cristallin devient rigide : même si le muscle ciliaire se contracte à fond, le cristallin ne bombe plus suffisamment. La \cle{vergence maximale} $C_{\text{max}}$ de l'œil diminue, tandis que la \cle{vergence minimale} $C_{\text{min}}$ (œil au repos, vision de loin) reste globalement inchangée pour un œil émétrope.

\begin{definition}[Presbytie]
La \textbf{presbytie} est la diminution physiologique de l'amplitude d'accommodation $\Delta C = C_{\text{max}} - C_{\text{min}}$ due au vieillissement du cristallin (perte d'élasticité, sclérose nucléaire). Elle apparaît vers 40--45 ans et progresse jusqu'à 60--65 ans.
\begin{itemize}
    \item Le \cle{punctum proximum} (PP) \textbf{recule} progressivement (vision nette de plus en plus lointaine).
    \item Le \cle{punctum remotum} (PR) \textbf{reste à l'infini} pour un œil initialement émétrope (vision de loin conservée).
\end{itemize}
Ce n'est \textbf{pas} un défaut de forme (longueur axiale, cornée), mais un \textbf{défaut d'accommodation}.
\end{definition}

\begin{savaistu}[L'amplitude d'accommodation au fil de la vie]
\begin{center}
\begin{tabularx}{\linewidth}{|c|X|X|}\hline
\textbf{Âge} & \textbf{PP (distance minimale)} & \textbf{Amplitude $\Delta C$ (dioptries)} \\ \hline
10 ans & 7 cm & $\approx 14\,\delta$ \\
20 ans & 10 cm & $\approx 10\,\delta$ \\
40 ans & 25 cm & $\approx 4\,\delta$ \\
50 ans & 50 cm & $\approx 2\,\delta$ \\
60 ans & 1 m & $\approx 1\,\delta$ \\
65 ans & $\infty$ (accommodation nulle) & $0\,\delta$ \\ \hline
\end{tabularx}
\end{center}
Vers 60 ans, l'amplitude tombe sous 1 dioptrie : l'œil est devenu un \emph{système à focale fixe}. C'est pourquoi la presbytie est universelle : même un myope ou un hypermétrope corrigé pour la vision de loin devient presbyte.
\end{savaistu}

\courssub{Évolution du punctum proximum : modélisation graphique}

Le recul du PP est la signature clinique de la presbytie. On le mesure en demandant au patient d'approcher un texte jusqu'à la limite de netteté.

\begin{popfigure}
\begin{tikzpicture}
\begin{axis}[
    width=\linewidth,
    height=0.55\linewidth,
    xlabel={Âge (ans)},
    ylabel={Distance du PP (cm)},
    xmin=10, xmax=70,
    ymin=0, ymax=110,
    xtick={10,20,30,40,50,60,70},
    ytick={0,10,25,50,100},
    grid=major,
    grid style={popmuted},
    axis line style={popdark},
    tick style={popdark},
    label style={font=\small,popdark},
    tick label style={font=\small,popdark},
    title={Recul du punctum proximum avec l'âge},
    title style={font=\normalsize\bfseries,popdark},
]
\addplot[popcurve, mark=*, mark size=2.5pt, popblue] coordinates {
    (10,7) (20,10) (30,14) (40,25) (45,35) (50,50) (55,70) (60,100) (65,150)
};
\addplot[only marks, mark=*, mark size=3pt, poporange] coordinates {
    (20,10) (40,25) (50,50) (60,100)
};
\node[poporange, anchor=south west, font=\small] at (axis cs:20,10) {20 ans : 10 cm};
\node[poporange, anchor=south west, font=\small] at (axis cs:40,25) {40 ans : 25 cm};
\node[poporange, anchor=south west, font=\small] at (axis cs:50,50) {50 ans : 50 cm};
\node[poporange, anchor=north west, font=\small] at (axis cs:60,100) {60 ans : 1 m};
\end{axis}
\end{tikzpicture}
\legende{Évolution typique de la distance du punctum proximum (PP) en fonction de l'âge. La courbe s'infléchit fortement après 45 ans. Au-delà de 60--65 ans, le PP tend vers l'infini : l'accommodation est nulle.}
\end{popfigure}

\emph{Interprétation.} La relation n'est pas linéaire : l'amplitude $\Delta C = 1/\text{PP}$ (en mètres) décroît approximativement linéairement avec l'âge (loi de Duane). À 20 ans, $\Delta C \approx 10\,\delta$ (PP = 10 cm) ; à 50 ans, $\Delta C \approx 2\,\delta$ (PP = 50 cm).

\courssub{Correction de la presbytie : le verre de lecture}

Puisque l'œil ne peut plus augmenter sa vergence pour voir net de près, on doit \textbf{lui apporter la vergence manquante} sous forme d'une lentille convergente placée devant l'œil (lunettes de lecture).

\begin{propriete}[Correction de la presbytie par une lentille convergente]
Pour un presbyte dont le PP est à $d_{\text{PP}}$ (en mètres) et qui veut lire distinctement à une distance $d_{\text{lecture}}$ (usuellement 25 cm = 0,25 m), la vergence $C_{\text{add}}$ de la lentille correctrice (addition) est :
\[
C_{\text{add}} = \frac{1}{d_{\text{lecture}}} - \frac{1}{d_{\text{PP}}} \quad (\text{en dioptries, } \delta)
\]
Cette lentille est \textbf{convergente} ($C_{\text{add}} > 0$). Elle forme une image virtuelle de l'objet (placé à $d_{\text{lecture}}$) au niveau du PP de l'œil, que l'œil peut alors voir net sans accommoder.
\end{propriete}

\begin{exemplebox}[Maman Ngo Bell : calcul de l'addition pour la lecture]
Maman Ngo Bell, 52 ans, commerçante au Marché Central de Douala, tient ses factures à 50 cm pour les lire distinctement. Son PP est donc à $d_{\text{PP}} = \SI{0,50}{\meter}$. Elle souhaite retrouver une lecture confortable à $d_{\text{lecture}} = \SI{0,25}{\meter}$.

\emph{Données :} $d_{\text{PP}} = \SI{0,50}{\meter}$, $d_{\text{lecture}} = \SI{0,25}{\meter}$.
\emph{Chercher :} Vergence $C_{\text{add}}$ du verre de lecture.

\emph{Solution :}
\begin{align*}
C_{\text{add}} &= \frac{1}{d_{\text{lecture}}} - \frac{1}{d_{\text{PP}}} \\
               &= \frac{1}{0,25} - \frac{1}{0,50} \\
               &= 4,00 - 2,00 \\
               &= +2,00\,\delta
\end{align*}
Maman Ngo Bell a besoin de verres \textbf{convergents de $+2,00\,\delta$} pour lire à 25 cm.
\vspace{0,5em}
\emph{Vérification dimensionnelle :} $[1/\text{m}] = \text{m}^{-1} = \delta$. Correct.
\emph{Ordre de grandeur :} À 50 ans, l'addition usuelle est $+1,50$ à $+2,00\,\delta$. Cohérent.
\end{exemplebox}

\begin{attention}[Piège : confusion PP / distance de lecture]
Ne pas confondre la \textbf{distance du PP} (où l'œil voit net \emph{sans} lunettes, accommodation maximale) et la \textbf{distance de lecture souhaitée} (où l'on place l'objet \emph{avec} lunettes).
\begin{itemize}
    \item Si le patient dit : « Je vois net à 50 cm » $\rightarrow$ $d_{\text{PP}} = \SI{0,50}{\meter}$.
    \item Si le patient dit : « Je veux lire à 30 cm » $\rightarrow$ $d_{\text{lecture}} = \SI{0,30}{\meter}$.
\end{itemize}
La formule utilise \textbf{les deux} distances. Oublier l'une des deux est l'erreur n°1 au Probatoire.
\end{attention}

\courssub{Presbytie associée à un autre défaut : verres bifocaux et progressifs}

La presbytie s'ajoute à l'état réfractif préexistant de l'œil.
\begin{itemize}
    \item \textbf{Émétrope presbyte} : vision de loin parfaite, lunettes convergentes \emph{uniquement pour le près} (demi-lunes ou lunettes qu'on ôte pour regarder loin).
    \item \textbf{Myope presbyte} : il porte une lentille divergente pour la vision de loin. Pour lire, il a besoin d'une puissance \emph{moins divergente} (ou convergente). \textbf{Astuce :} un myope léger (ex. $-2\,\delta$, PP $\approx$ 50 cm) peut \textbf{retirer ses lunettes} pour lire : sa myopie compense sa presbytie !
    \item \textbf{Hypermétrope presbyte} : il porte déjà une lentille convergente pour la vision de loin. Pour le près, il lui faut une convergence \emph{supplémentaire} (addition).
\end{itemize}
Dans les deux derniers cas, changer de lunettes sans cesse est pénible. On utilise :
\begin{itemize}
    \item \textbf{Verres à double foyer (bifocaux) :} partie supérieure pour la vision de loin (correction du défaut principal), partie inférieure (segment) pour la vision de près (addition). Ligne de séparation visible.
    \item \textbf{Verres progressifs :} variation continue de la vergence de la zone de loin vers la zone de près. Pas de ligne visible, esthétique, mais zone de progression latérale (distorsion) et adaptation nécessaire.
\end{itemize}

\begin{popfigure}
\begin{tikzpicture}[
    scale=0.9,
    every node/.style={font=\small},
    eye/.style={draw=popdark, thick, fill=popblueL!30, minimum width=2.5cm, minimum height=1.8cm, rounded corners=0.9cm},
    lens/.style={draw=poppurple, very thick, fill=poppurpleL!20, minimum width=0.4cm, minimum height=1.5cm, rounded corners=0.2cm},
    ray/.style={->, poporange, thick},
    obj/.style={draw=popgreen, thick, fill=popgreenL!30, minimum width=0.3cm, minimum height=1cm},
    img/.style={draw=poppink, thick, fill=poppinkL!30, minimum width=0.3cm, minimum height=1cm, dashed},
]
% --- MYOPIE ---
\begin{scope}[xshift=0cm]
\node[eye] (eye1) at (0,0) {};
\node[lens, left=0.5cm of eye1.west] (lens1) {};
\node[obj, left=1.5cm of lens1.west, align=center] (obj1){Objet\\loin ($\infty$)};
\draw[ray] (obj1.east) -- (lens1.west) node[midway,above,sloped,font=\tiny] {$\parallel$};
\draw[ray] (lens1.east) -- (eye1.west) node[midway,above,sloped,font=\tiny] {divergentes};
\node[img, left=0.8cm of lens1.west, align=center] (img1){Image\\virtuelle au PR};
\draw[ray, dashed] (lens1.west) -- (img1.east);
\node[align=center, popdark] at (0,-1.6) {\textbf{MYOPIE}\\\textit{Verre divergent ($C<0$)}\\\textit{Recule l'image au PR}};
\end{scope}

% --- HYPERMÉTROPIE ---
\begin{scope}[xshift=5.5cm]
\node[eye] (eye2) at (0,0) {};
\node[lens, left=0.5cm of eye2.west] (lens2) {};
\node[obj, left=1.5cm of lens2.west, align=center] (obj2){Objet\\loin ($\infty$)};
\draw[ray] (obj2.east) -- (lens2.west) node[midway,above,sloped,font=\tiny] {$\parallel$};
\draw[ray] (lens2.east) -- (eye2.west) node[midway,above,sloped,font=\tiny] {convergentes};
\node[img, right=1.2cm of eye2.east, align=center] (img2){Image\\réelle sur rétine};
\draw[ray, dashed] (lens2.east) -- (img2.west);
\node[align=center, popdark] at (0,-1.6) {\textbf{HYPERMÉTROPIE}\\\textit{Verre convergent ($C>0$)}\\\textit{Avance le foyer sur rétine}};
\end{scope}

% --- PRESBYTIE ---
\begin{scope}[xshift=11cm]
\node[eye] (eye3) at (0,0) {};
\node[lens, left=0.5cm of eye3.west] (lens3) {};
\node[obj, left=0.8cm of lens3.west, align=center] (obj3){Objet\\près (25 cm)};
\draw[ray] (obj3.east) -- (lens3.west);
\draw[ray] (lens3.east) -- (eye3.west) node[midway,above,sloped,font=\tiny] {convergentes};
\node[img, left=1.5cm of lens3.west, align=center] (img3){Image\\virtuelle au PP};
\draw[ray, dashed] (lens3.west) -- (img3.east);
\node[align=center, popdark] at (0,-1.6) {\textbf{PRESBYTIE}\\\textit{Verre convergent ($C>0$)}\\\textit{Pour la vision de près seulement}};
\end{scope}
\end{tikzpicture}
\legende{Synthèse visuelle des trois défauts et de leurs corrections. L'œil réduit est schématisé en bleu. Le verre correcteur (violet) modifie la divergence/convergence du faisceau incident pour que l'œil forme une image nette sur la rétine (ou au PP/PR).}
\end{popfigure}

\courssub{Complément Probatoire : association lentille correctrice + œil (lentilles accolées)}

Au Probatoire, on peut vous demander de justifier le calcul de la correction en utilisant le modèle des \textbf{lentilles minces accolées}. L'œil réduit (cornée + cristallin au repos) est modélisé comme une lentille unique de vergence $C_{\text{œil}}$. La lentille de correction (lunettes) de vergence $C_{\text{corr}}$ est placée au contact (distance négligeable).

\begin{propriete}[Vergences de lentilles minces accolées (admise, exigible au Probatoire)]
Si deux lentilles minces de vergences $C_1$ et $C_2$ sont accolées (distance nulle entre elles), le système équivaut à une lentille unique de vergence :
\[
C_{\text{eq}} = C_1 + C_2
\]
\emph{Justification (rappel) :} Pour un objet à l'infini, la première lentille donne une image à son foyer $f_1' = 1/C_1$. Cette image devient l'objet de la seconde lentille (distance objet $p_2 = -f_1'$ car objet virtuel pour la seconde si $C_1>0$). En utilisant la relation de conjugaison $C_2 = 1/p_2' - 1/p_2$ et $p_2 = -1/C_1$, on trouve $1/p_2' = C_1 + C_2$. Le foyer image du système est $f' = 1/(C_1+C_2)$.
\end{propriete}

\emph{Application à la correction de l'amétropie.}
L'œil a une vergence $C_{\text{œil}}$. Pour un objet à l'infini, l'œil seul formerait une image en $F_{\text{œil}}'$ tel que $\overline{F_{\text{œil}}'} = 1/C_{\text{œil}}$.
\begin{itemize}
    \item \textbf{Myopie :} $F_{\text{œil}}'$ est \emph{devant} la rétine. Il faut \emph{diminuer} la vergence totale : $C_{\text{corr}} < 0$ (divergente). $C_{\text{œil}} + C_{\text{corr}} = C_{\text{émétrope}}$.
    \item \textbf{Hypermétropie :} $F_{\text{œil}}'$ est \emph{derrière} la rétine. Il faut \emph{augmenter} la vergence totale : $C_{\text{corr}} > 0$ (convergente).
    \item \textbf{Presbytie (vision de près) :} L'objet est à distance finie $p$. L'œil au repos ($C_{\text{min}}$) ne peut pas faire le point. On ajoute $C_{\text{add}} > 0$ pour que $(C_{\text{min}} + C_{\text{add}})$ permette la conjugaison objet $\to$ rétine.
\end{itemize}

\begin{exoresolu}[Correction d'un hypermétrope presbyte : addition et vision de loin]
M. Essomba, 55 ans, est hypermétrope. Son œil réduit a une vergence au repos $C_{\text{œil}} = +58,0\,\delta$. La distance rétine-cristallin (longueur axiale) est $L = \SI{22,0}{\milli\meter} = \SI{0,022}{\meter}$. Son PP est à 80 cm.
\begin{enumerate}
    \item Vérifier qu'il est hypermétrope et calculer la vergence $C_{\text{loin}}$ de son verre de loin.
    \item Calculer l'addition $C_{\text{add}}$ pour lire à 25 cm.
    \item Donner la vergence du segment de près de son verre bifocal.
\end{enumerate}
\tcblower
\emph{Solution.}
\begin{enumerate}
    \item \textbf{Vision de loin (objet à $\infty$).} L'œil seul forme l'image en $F_{\text{œil}}'$ tel que $F_{\text{œil}}' = 1/C_{\text{œil}} = 1/58,0 \approx \SI{0,01724}{\meter} = \SI{17,24}{\milli\meter}$.
    La rétine est à $L = \SI{22,0}{\milli\meter}$. $F_{\text{œil}}'$ est \textbf{devant} la rétine $\rightarrow$ \textbf{myopie} !
    \textbf{Attention :} L'énoncé disait « hypermétrope » pour tester la vigilance. Le calcul montre $f'_{\text{œil}} = \SI{17,2}{\milli\meter} < L = \SI{22,0}{\milli\meter}$ : le foyer tombe devant la rétine, c'est une myopie.
    
    Calcul pour myope : la vergence totale nécessaire pour mettre l'image de l'infini sur la rétine est $C_{\text{total}} = 1/L = 1/0,022 \approx 45,45\,\delta$.
    Avec lentilles accolées : $C_{\text{total}} = C_{\text{œil}} + C_{\text{loin}}$.
    \[
    C_{\text{loin}} = C_{\text{total}} - C_{\text{œil}} = 45,45 - 58,0 = -12,55\,\delta \approx \mathbf{-12,6\,\delta}
    \]
    Verre \textbf{divergent} de $-12,6\,\delta$.
    
    \item \textbf{Vision de près (lecture à 25 cm).} PP = \SI{0,80}{\meter}. Addition pour lecture à \SI{0,25}{\meter} :
    \[
    C_{\text{add}} = \frac{1}{0,25} - \frac{1}{0,80} = 4,00 - 1,25 = \mathbf{+2,75\,\delta}
    \]
    
    \item \textbf{Verre bifocal :}
    \begin{itemize}
        \item Partie supérieure (vision de loin) : $C_{\text{loin}} = -12,6\,\delta$.
        \item Partie inférieure (vision de près) : $C_{\text{près}} = C_{\text{loin}} + C_{\text{add}} = -12,6 + 2,75 = \mathbf{-9,85\,\delta}$.
    \end{itemize}
    Le segment de près est \emph{moins divergent} que la partie de loin (ou plus convergent).
\end{enumerate}
\emph{Moralité :} Vérifiez toujours le type d'amétropie par le calcul ($f'_{\text{œil}}$ vs $L$) avant d'appliquer les formules toutes faites !
\end{exoresolu}

\courssub{Bilan général des corrections et méthode de résolution}

\begin{aretenir}[Tableau de synthèse : les trois défauts et leurs corrections]
\begin{center}
\begin{tabularx}{\linewidth}{|c|X|X|X|}\hline
\textbf{Défaut} & \textbf{Origine} & \textbf{Verre correcteur} & \textbf{Principe} \\ \hline
\textbf{Myopie} & Œil trop long ou trop convergent ($F'$ devant rétine). PR \textbf{fini} devant l'œil. & \textbf{Divergent} ($C < 0$) & Recule l'image virtuelle de l'objet lointain au PR de l'œil. \\ \hline
\textbf{Hypermétropie} & Œil trop court ou pas assez convergent ($F'$ derrière rétine). PR \textbf{virtuel} derrière l'œil. & \textbf{Convergent} ($C > 0$) & Avance l'image réelle de l'objet lointain sur la rétine (ou au PR virtuel). \\ \hline
\textbf{Presbytie} & Cristallin rigide : $\Delta C \downarrow$. PP \textbf{recule}. PR inchangé (si émétrope). & \textbf{Convergent} ($C > 0$) \emph{pour le près seulement} & Forme une image virtuelle de l'objet proche au PP de l'œil. \\ \hline
\end{tabularx}
\end{center}
\end{aretenir}

\begin{methode}[Résolution systématique d'un problème de correction (Probatoire)]
\begin{enumerate}
    \item \textbf{Identifier le défaut} à partir des données : PR (myopie/hypermétropie) ou PP/âge (presbytie). Calculer $f'_{\text{œil}}$ si $C_{\text{œil}}$ et $L$ donnés.
    \item \textbf{Dessiner le schéma} : œil réduit, verre correcteur au contact, objet (loin $\infty$ ou près $d$), image (sur rétine, au PR ou au PP). Noter les vergences.
    \item \textbf{Écrire la relation de conjugaison} pour le système (verre + œil) ou utiliser la formule directe :
    \begin{itemize}
        \item Myopie (vision loin) : $C_{\text{corr}} = \dfrac{1}{d_{\text{PR}}}$ ($d_{\text{PR}} < 0$ si PR réel devant l'œil $\rightarrow C<0$).
        \item Hypermétropie (vision loin) : $C_{\text{corr}} = \dfrac{1}{L} - C_{\text{œil}}$ (souvent $>0$).
        \item Presbytie (vision près) : $C_{\text{add}} = \dfrac{1}{d_{\text{lecture}}} - \dfrac{1}{d_{\text{PP}}}$.
    \end{itemize}
    \item \textbf{Calculer} avec unités SI (mètres $\rightarrow$ dioptries). Vérifier le signe ($+$ convergent, $-$ divergent).
    \item \textbf{Conclure} par une phrase : « Le verre correcteur est un verre [convergent/divergent] de vergence $C = \dots\,\delta$. »
\end{enumerate}
\end{methode}

\begin{attention}[Erreurs fatales au Probatoire]
\begin{itemize}
    \item \textbf{Signe des distances :} PR réel (myopie) : distance \textbf{négative} (objet virtuel pour l'œil) $\rightarrow C_{\text{corr}}$ négatif. PR virtuel (hypermétropie) : distance \textbf{positive} $\rightarrow C_{\text{corr}}$ positif. PP : toujours distance \textbf{positive} (objet réel pour l'œil).
    \item \textbf{Unités :} Distances en \textbf{mètres} pour obtenir des dioptries ($\text{m}^{-1}$). $\SI{50}{\centi\meter} = \SI{0,50}{\meter}$.
    \item \textbf{Confusion loin/près :} Le verre de presbytie ne sert \textbf{que pour le près}. Il floute la vision de loin. Ne pas l'utiliser pour conduire !
    \item \textbf{Myope presbyte :} Ne pas prescrire systématiquement un bifocal. Demander : « Voyez-vous net de près \emph{sans} vos lunettes de loin ? » Si oui (myopie $\approx$ addition), pas de segment de près nécessaire.
\end{itemize}
\end{attention}

\courssub{Exercices d'entraînement}

\begin{exercice}[Exercice 1 : Presbytie simple]
Mme Atangana, 48 ans, institutrice à Yaoundé, constate qu'elle doit tenir ses copies à 40 cm pour les corriger distinctement. Elle souhaite lire à 25 cm.
\begin{enumerate}
    \item Quel est son défaut visuel ? Justifier.
    \item Calculer la vergence des verres de lecture dont elle a besoin.
    \item Si elle met ces verres pour regarder un élève au tableau (loin), que se passe-t-il ?
\end{enumerate}
\end{exercice}

\begin{exercice}[Exercice 2 : Myope presbyte - Le cas « M. Ngono »]
M. Ngono, 55 ans, est myope. Il porte des verres $-3,00\,\delta$ pour la vision de loin. Son PP \emph{sans lunettes} est à 33 cm.
\begin{enumerate}
    \item Quelle est son amplitude d'accommodation résiduelle (en $\delta$) ?
    \item Peut-il lire à 25 cm \emph{sans} ôter ses lunettes de loin ? Justifier par un calcul.
    \item Quelle addition ($C_{\text{add}}$) faut-il ajouter à sa correction de loin pour qu'il lise à 25 cm \emph{avec} ses lunettes ?
\end{enumerate}
\end{exercice}

\begin{corrige}[Exercice 1]
\begin{enumerate}
    \item PP = 40 cm = \SI{0,40}{\meter}. Âge 48 ans $\rightarrow$ \textbf{Presbytie} (PP reculé, PR à l'infini supposé).
    \item $C_{\text{add}} = 1/0,25 - 1/0,40 = 4,00 - 2,50 = \mathbf{+1,50\,\delta}$.
    \item Avec des verres $+1,50\,\delta$, un objet à l'infini donne une image virtuelle au foyer image du verre : $f' = 1/1,50 \approx \SI{0,67}{\meter} = \SI{67}{\centi\meter}$ \emph{devant} le verre. L'œil (émétrope, PR=$\infty$) ne peut pas faire le point sur un objet virtuel à 67 cm (il ne peut accommoder que jusqu'à 40 cm). \textbf{Vision de loin floue.} Elle doit ôter ses lunettes pour voir le tableau.
\end{enumerate}
\end{corrige}

\begin{corrige}[Exercice 2]
\begin{enumerate}
    \item Sans lunettes, PP = \SI{0,33}{\meter}. PR = $\infty$ (myope corrigé $-3\,\delta$ $\rightarrow$ œil nu a PR à $-1/3 \approx -\SI{0,33}{\meter}$). Amplitude $\Delta C = 1/0,33 - 0 \approx \mathbf{3,0\,\delta}$. Faible, typique 55 ans.
    \item \textbf{Sans ôter lunettes :} Il porte $-3,00\,\delta$. Objet à 25 cm ($p = -\SI{0,25}{\meter}$). Vergence incidente sur l'œil : $C_{\text{inc}} = C_{\text{verre}} + 1/p = -3,00 + (-4,00) = -7,00\,\delta$. L'œil nu (myope $-3\,\delta$) a une vergence $C_{\text{œil}} \approx 60+3=63\,\delta$ (exemple). L'œil voit net si l'objet est à son PR (33 cm) ou s'il accommode. Ici l'objet virtuel pour l'œil est à $1/7 \approx \SI{14}{\centi\meter}$ (très près). L'œil devrait accommoder $7\,\delta$ : impossible ($\Delta C = 3\,\delta$). \textbf{Ne peut pas lire avec ses lunettes de loin.}
    \item Il faut que l'ensemble (verre de loin + addition) présente à l'œil nu un objet virtuel à son PP (\SI{0,33}{\meter}).
    Vergence nécessaire incidente sur l'œil nu : $V_{\text{œil}} = -1/0,33 \approx -3,0\,\delta$ (objet virtuel à 33 cm).
    Objet réel à 25 cm : $V_{\text{obj}} = -4,0\,\delta$.
    Verge totale du segment de près (loin + add) : $C_{\text{près}} = V_{\text{œil}} - V_{\text{obj}} = -3,0 - (-4,0) = \mathbf{+1,0\,\delta}$.
    Verre de loin = $-3,0\,\delta$. Donc l'addition $C_{\text{add}} = C_{\text{près}} - C_{\text{loin}} = +1,0 - (-3,0) = \mathbf{+4,0\,\delta}$.
    \textbf{Réponse :} Addition $C_{\text{add}} = \mathbf{+4,0\,\delta}$. Segment de près = $+1,0\,\delta$ (moins divergent que le segment de loin).
\end{enumerate}
\end{corrige}

\begin{savaistu}[Presbytie et chirurgie : au-delà des lunettes]
Au Cameroun, comme partout, la chirurgie réfractive (LASIK, PKR) corrige myopie/hypermétropie/astigmatisme en sculptant la cornée. Mais elle \textbf{ne restaure pas l'accommodation} (le cristallin reste rigide). Un myope opéré à 30 ans deviendra presbyte vers 45 ans comme tout le monde.
Des solutions chirurgicales spécifiques existent pour la presbytie :
\begin{itemize}
    \item \textbf{Implant multifocal / à profondeur de champ étendue (EDOF) :} Remplace le cristallin (chirurgie de la cataracte ou \emph{clear lens exchange}). Donne vision loin + près simultanée. Standard actuel à l'Hôpital Général de Douala ou à la Clinique de l'Espérance.
    \item \textbf{PresbyLASIK / Supracor :} Crée une zone de près sur la cornée (multifocalité cornéenne). Moins utilisé aujourd'hui.
    \item \textbf{Inlays cornéens (ex. Kamra) :} Petit anneau opaque au centre de la cornée (effet sténopé $\rightarrow$ profondeur de champ accrue). En régression.
\end{itemize}
Le choix dépend de l'âge, de l'état du cristallin (cataracte naissante ?), des besoins professionnels (conducteur de nuit ? chirurgien ?) et du budget. L'optométriste et l'ophtalmologiste guident le patient.
\end{savaistu}

Cette section achève l'étude des défauts de l'œil réduit. Vous maîtrisez maintenant la modélisation optique de l'œil, l'analyse de ses défauts (myopie, hypermétropie, presbytie, astigmatisme vu en C) et le calcul rigoureux de leurs corrections. La clé du Probatoire : \textbf{schéma clair $\rightarrow$ identification du défaut $\rightarrow$ formule adaptée $\rightarrow$ calcul en SI $\rightarrow$ conclusion signée}. Bon travail !

\coursec{Méthode générale de résolution et synthèse}

Nous avons étudié séparément la myopie, l'hypermétropie et la presbytie, avec pour chacune sa stratégie de correction. Au Probatoire, l'exercice ne précise jamais le défaut : c'est à toi de le diagnostiquer à partir des données, puis de construire la correction. Cette dernière section te donne une \cle{méthode unique}, applicable à \emph{tous} les exercices de correction de l'œil, puis une synthèse complète de la leçon.

\courssub{La fiche de résolution type en six étapes}

\begin{methode}[Résoudre tout exercice de correction de l'œil]
\begin{enumerate}
\item \textbf{Lire et convertir.} Repérer dans l'énoncé le punctum remotum (PR), le punctum proximum (PP), les distances et les vergences. Convertir \emph{immédiatement} toutes les distances en \textbf{mètres} : $25\ \mathrm{cm} = \num{0,25}\ \mathrm{m}$, $2\ \mathrm{m}$ reste $2\ \mathrm{m}$.
\item \textbf{Diagnostiquer.} Comparer PR et PP à ceux de l'œil normal (PR à l'infini, PP à $d_m = 25\ \mathrm{cm}$) :
\begin{itemize}
\item PR fini (devant l'œil) $\Rightarrow$ \cle{myopie} ;
\item PR à l'infini mais PP reculé ($> 25\ \mathrm{cm}$) $\Rightarrow$ \cle{hypermétropie} si le sujet est jeune, \cle{presbytie} si le sujet est âgé (le PR restant normal) ;
\item PR et PP normaux $\Rightarrow$ œil emmétrope, pas de correction.
\end{itemize}
\item \textbf{Écrire la condition de correction.} La lentille correctrice doit donner, d'un objet placé là où l'œil \emph{devrait} voir, une image située là où l'œil \emph{peut} voir :
\begin{itemize}
\item myope : objet réel à l'infini $\rightarrow$ image \emph{virtuelle} au PR ;
\item hypermétrope ou presbyte : objet réel à $25\ \mathrm{cm}$ $\rightarrow$ image \emph{virtuelle} au PP.
\end{itemize}
\item \textbf{Appliquer la relation de conjugaison} des lentilles minces, axe orienté dans le sens de la lumière :
\[ \frac{1}{\overline{OA'}} - \frac{1}{\overline{OA}} = \frac{1}{f'} = C \]
en remplaçant $\overline{OA}$ et $\overline{OA'}$ par leurs valeurs \emph{algébriques} (signe inclus !).
\item \textbf{Conclure} sur la nature du verre : $C > 0$ $\Rightarrow$ lentille \textbf{convergente} ; $C < 0$ $\Rightarrow$ lentille \textbf{divergente}. Donner $C$ en dioptries ($\delta$) et $f' = 1/C$ en mètres.
\item \textbf{Vérifier la cohérence.} Le signe de $C$ correspond-il au défaut diagnostiqué à l'étape 2 ? L'ordre de grandeur est-il plausible (les verres correcteurs courants vont de $-10\ \delta$ à $+4\ \delta$ environ) ? Si possible, recalculer une grandeur non utilisée (par exemple le PP corrigé) pour valider le résultat.
\end{enumerate}
\end{methode}

\begin{attention}[L'erreur classique du Probatoire : les signes algébriques]
La faute la plus fréquente consiste à écrire $\dfrac{1}{OA'} - \dfrac{1}{OA} = \dfrac{1}{f'}$ avec des \emph{distances positives} partout. C'est faux : la relation de conjugaison est \textbf{algébrique}. Pour un œil myope corrigé, l'image au PR est \emph{virtuelle}, située \emph{avant} la lentille : $\overline{OA'} < 0$. De même, l'objet à $25\ \mathrm{cm}$ pour un hypermétrope est réel, donc $\overline{OA} = -\num{0,25}\ \mathrm{m}$ et non $+\num{0,25}\ \mathrm{m}$. \textbf{Réflexe obligatoire} : orienter l'axe dans le sens de propagation de la lumière, placer l'origine $O$ au centre optique de la lentille, puis affecter un signe à chaque grandeur \emph{avant} tout calcul. Une vérification rapide : un objet réel donne toujours $\overline{OA} < 0$ ; une image virtuelle donne toujours $\overline{OA'} < 0$.
\end{attention}

\courssub{Arbre de décision : du diagnostic au verre correcteur}

L'organigramme ci-dessous résume la démarche des étapes 1, 2 et 5 : deux questions suffisent pour identifier le défaut et la nature du verre.

\begin{popfigure}
\begin{center}
\begin{tikzpicture}[
  font=\small,
  quest/.style={draw=popblue, fill=popblueL, rounded corners=2pt, align=center, minimum width=3.2cm, minimum height=0.9cm},
  diag/.style={draw=poporange, fill=poporangeL, rounded corners=2pt, align=center, minimum width=2.6cm, minimum height=0.9cm},
  verre/.style={draw=popgreen, fill=popgreenL, rounded corners=2pt, align=center, minimum width=2.6cm, minimum height=0.9cm},
  fleche/.style={-{Stealth[length=2.5mm]}, thick, popdark},
  node distance=0.55cm and 0.9cm]
\node[quest] (q1) {PR à l'infini ?};
\node[diag, below left=0.7cm and 0.2cm of q1, align=center] (myope){\textbf{Myopie}\\ PR fini devant l'œil};
\node[quest, below right=0.7cm and 0.2cm of q1] (q2) {PP à $25\ \mathrm{cm}$ ?};
\node[verre, below=0.7cm of myope, align=center] (vdiv){Lentille \textbf{divergente}\\ $C = \dfrac{1}{\overline{\mathrm{PR}}} < 0$};
\node[diag, below left=0.7cm and 0.1cm of q2, align=center] (hyper){\textbf{Hypermétropie}\\ PP $> 25\ \mathrm{cm}$,\\ sujet jeune};
\node[diag, below right=0.7cm and 0.1cm of q2, align=center] (presb){\textbf{Presbytie}\\ PP reculé,\\ sujet âgé, PR normal};
\node[verre, below=0.7cm of hyper, align=center] (vconv1){Lentille \textbf{convergente}\\ $C > 0$};
\node[verre, below=0.7cm of presb, align=center] (vconv2){Lentille \textbf{convergente}\\ pour la vision de près};
\node[diag, right=1.1cm of q2, align=center] (normal){\textbf{Œil normal}\\ (emmétrope)};
\draw[fleche] (q1) -- node[above left, font=\footnotesize]{non} (myope);
\draw[fleche] (q1) -- node[above right, font=\footnotesize]{oui} (q2);
\draw[fleche] (myope) -- (vdiv);
\draw[fleche] (q2) -- node[above left, font=\footnotesize]{non} (hyper);
\draw[fleche] (q2) -- node[above right, font=\footnotesize]{non, âgé} (presb);
\draw[fleche] (q2) -- node[above, font=\footnotesize]{oui} (normal);
\draw[fleche] (hyper) -- (vconv1);
\draw[fleche] (presb) -- (vconv2);
\end{tikzpicture}
\end{center}
\legende{Arbre de décision : deux questions (position du PR, position du PP) conduisent au diagnostic et au type de verre correcteur. Pour le myope, $\overline{\mathrm{PR}}$ désigne la distance algébrique du PR à la lentille, négative car le PR est devant l'œil.}
\end{popfigure}

\courssub{Exemple de synthèse complet}

\begin{exoresolu}[Œil de PR $= 2\ \mathrm{m}$ et PP $= 20\ \mathrm{cm}$ : diagnostic, correction, vérification]
Un élève de Première C consulte l'ophtalmologiste de l'hôpital central de Yaoundé. On mesure : punctum remotum $\mathrm{PR} = 2\ \mathrm{m}$ devant l'œil ; punctum proximum $\mathrm{PP} = 20\ \mathrm{cm}$. (1) Diagnostiquer le défaut. (2) Déterminer la vergence du verre correcteur (lentille accolée à l'œil) permettant de voir net un objet à l'infini. (3) Calculer le PP corrigé, c'est-à-dire la distance minimale de vision nette avec les lunettes.
\tcblower
\textbf{Étape 1 — Données en mètres.} $\overline{\mathrm{PR}} = -2\ \mathrm{m}$ (le PR est devant l'œil, donc en amont de la lentille, dans le sens opposé à la lumière émergente) ; $\mathrm{PP} = \num{0,20}\ \mathrm{m}$.
\par\textbf{Étape 2 — Diagnostic.} Le PR est \emph{fini} (2 m au lieu de l'infini) et le PP est légèrement \emph{rapproché} ($20\ \mathrm{cm} < 25\ \mathrm{cm}$) : c'est une \textbf{myopie légère}. L'œil est trop convergent (ou trop long) ; l'image d'un objet lointain se forme en avant de la rétine.
\par\textbf{Étape 3 — Condition de correction.} Le verre doit donner d'un objet réel à l'infini ($\overline{OA} = -\infty$) une image \emph{virtuelle} située au PR de l'œil nu : $\overline{OA'} = -2\ \mathrm{m}$. L'œil voit alors cette image nette sans effort d'accommodation.
\par\textbf{Étape 4 — Relation de conjugaison.}
\begin{align*}
C &= \frac{1}{\overline{OA'}} - \frac{1}{\overline{OA}} = \frac{1}{-2} - \frac{1}{-\infty} = -\num{0,5}\ \delta.
\end{align*}
\par\textbf{Étape 5 — Conclusion.} $C = -\num{0,5}\ \delta < 0$ : verre \textbf{divergent}, de distance focale image $f' = \dfrac{1}{C} = -2\ \mathrm{m}$. Cohérent avec le diagnostic de myopie.
\par\textbf{Étape 6 — Vérification par le PP corrigé.} Avec les lunettes, l'objet le plus proche visible nettement est celui dont le verre donne une image virtuelle au PP de l'œil nu : $\overline{OA'} = -\num{0,20}\ \mathrm{m}$. Alors :
\begin{align*}
\frac{1}{\overline{OA}} &= \frac{1}{\overline{OA'}} - C = \frac{1}{-\num{0,20}} - (-\num{0,5}) = -5 + \num{0,5} = -\num{4,5}\ \delta,\\
\overline{OA} &= \frac{1}{-\num{4,5}} \approx -\num{0,22}\ \mathrm{m}.
\end{align*}
Le PP corrigé est à environ $22\ \mathrm{cm}$ : la correction éloigne légèrement le point de vision de près, ce qui reste confortable ($22\ \mathrm{cm} < 25\ \mathrm{cm}$). Le résultat est cohérent : la myopie est corrigée sans gêner la lecture.
\end{exoresolu}

\courssub{Constructions graphiques : marche des rayons avant et après correction}

Savoir tracer les rayons est un savoir-faire exigible. Retiens les trois rayons particuliers d'une lentille mince : le rayon passant par le centre optique $O$ n'est \textbf{pas dévié} ; le rayon incident parallèle à l'axe émerge en passant par le foyer image $F'$ (lentille convergente) ou \emph{comme s'il} venait de $F'$ (lentille divergente) ; le rayon passant par le foyer objet $F$ émerge parallèle à l'axe.

\textbf{Avant correction (œil myope).} Les rayons issus d'un point objet à l'infini convergent en un point \emph{en avant} de la rétine, puis divergent : la tache sur la rétine est floue.

\textbf{Après correction.} La lentille divergente rend les rayons légèrement divergents avant l'entrée dans l'œil ; le cristallin les ramène alors exactement sur la rétine. Pour la construction, on procède en deux temps : (1) la lentille correctrice donne de l'objet à l'infini une image virtuelle au PR (les prolongements des rayons émergents passent par le PR) ; (2) cette image virtuelle sert d'objet pour l'œil réduit, qui la projette nettement sur la rétine.

\begin{attention}[Pièges de tracé]
Pour une lentille divergente, les rayons émergents ne se croisent jamais réellement : l'image est virtuelle, obtenue en \textbf{prolongeant en pointillés} les rayons émergents vers l'amont. Ne trace jamais un rayon réel en pointillé ni un prolongement en trait plein : le correcteur y est attentif. Respecte aussi l'échelle annoncée et oriente l'axe optique (flèche dans le sens de la lumière).
\end{attention}

\courssub{Tableau final à retenir}

\begin{propriete}[Bilan des défauts et corrections — à connaître par cœur]
\begin{center}
\begin{tabularx}{\linewidth}{|l|X|X|X|}
\hline
\rowcolor{popblueL} \textbf{Défaut} & \textbf{Signature} & \textbf{Cause} & \textbf{Correction} \\
\hline
Myopie & PR fini (rapproché), PP $< 25\ \mathrm{cm}$ & Œil trop convergent ou trop long : image en avant de la rétine & Lentille \textbf{divergente}, $C = \dfrac{1}{\overline{\mathrm{PR}}} < 0$ \\
\hline
Hypermétropie & PR à l'infini (ou derrière l'œil), PP $> 25\ \mathrm{cm}$ & Œil pas assez convergent : image d'un objet proche derrière la rétine & Lentille \textbf{convergente}, $C > 0$ \\
\hline
Presbytie & PP qui recule avec l'âge, PR inchangé & Cristallin vieillissant, accommodation insuffisante & Lentille \textbf{convergente} pour la vision de près \\
\hline
\end{tabularx}
\end{center}
Dans les trois cas, la stratégie est identique : la lentille correctrice fabrique une image \emph{virtuelle} située dans le domaine de vision nette de l'œil (entre PP et PR), image que l'œil projette ensuite sur la rétine.
\end{propriete}

\courssub{Carte mentale de synthèse de la leçon}

\begin{popfigure}
\begin{center}
\begin{tikzpicture}[
  font=\small,
  racine/.style={draw=popdark, fill=popblue, text=white, rounded corners=3pt, align=center, minimum width=3.4cm, minimum height=1cm, font=\small\bfseries},
  branche/.style={draw=popblue, fill=popblueL, rounded corners=2pt, align=center, minimum width=2.9cm, minimum height=0.85cm},
  feuille/.style={draw=popmuted, fill=popcream, rounded corners=2pt, align=center, minimum width=2.7cm, minimum height=0.8cm, font=\footnotesize},
  lien/.style={-{Stealth[length=2.2mm]}, thick, popmuted}]
\node[racine, align=center] (oeil){L'ŒIL RÉDUIT\\ (lentille $+$ écran : rétine)};
\node[branche, above left=0.9cm and 0.4cm of oeil, align=center] (accom){Accommodation\\ Vergence variable};
\node[branche, below left=0.9cm and 0.4cm of oeil, align=center] (points){Points particuliers\\ PR et PP};
\node[branche, above right=0.9cm and 0.4cm of oeil, align=center] (defauts){Défauts\\ d'accommodation};
\node[branche, below right=0.9cm and 0.4cm of oeil, align=center] (correc){Corrections\\ Relation de conjugaison};
\node[feuille, above=0.6cm of accom, align=center] (a1){Cristallin se bombe\\ pour voir de près};
\node[feuille, left=0.6cm of points, align=center] (p1){Œil normal : PR $= \infty$,\\ PP $= 25\ \mathrm{cm}$};
\node[feuille, above=0.6cm of defauts, align=center] (d1){Myopie : PR fini\\ Hypermétropie : PP reculé\\ Presbytie : PP reculé (âge)};
\node[feuille, below=0.6cm of correc, align=center] (c1){Myope : divergente\\ Hypermétrope : convergente\\ Presbyte : convergente (près)};
\draw[lien] (oeil) -- (accom);
\draw[lien] (oeil) -- (points);
\draw[lien] (oeil) -- (defauts);
\draw[lien] (oeil) -- (correc);
\draw[lien] (accom) -- (a1);
\draw[lien] (points) -- (p1);
\draw[lien] (defauts) -- (d1);
\draw[lien] (correc) -- (c1);
\end{tikzpicture}
\end{center}
\legende{Carte mentale de synthèse : de la modélisation de l'œil réduit aux corrections des défauts d'accommodation.}
\end{popfigure}

\begin{aretenir}[L'essentiel de la section]
\begin{itemize}
\item Toute correction de l'œil se résout par la même méthode en six étapes : convertir, diagnostiquer (PR et PP), écrire la condition de correction, appliquer $\dfrac{1}{\overline{OA'}} - \dfrac{1}{\overline{OA}} = C$ en \textbf{algébrique}, conclure, vérifier.
\item La lentille correctrice donne toujours une image \emph{virtuelle} placée dans le domaine de vision nette de l'œil : au PR pour le myope (objet à l'infini), au PP pour l'hypermétrope ou le presbyte (objet à $25\ \mathrm{cm}$).
\item Myopie $\Rightarrow$ verre divergent ($C<0$) ; hypermétropie et presbytie $\Rightarrow$ verre convergent ($C>0$).
\item Un résultat sans unité (dioptrie), sans signe justifié ou sans vérification de cohérence perd des points : soigne la rédaction jusqu'à la dernière ligne.
\end{itemize}
\end{aretenir}

\begin{exercice}[Pour t'entraîner]
Un patient de 55 ans a un PR à l'infini et un PP à $50\ \mathrm{cm}$. (1) Identifier le défaut. (2) Calculer la vergence du verre lui permettant de lire à $25\ \mathrm{cm}$. (3) Vérifier le signe et l'ordre de grandeur du résultat.
\end{exercice}

\begin{corrige}[Exercice : presbytie]
(1) PR normal, PP reculé chez un sujet âgé : \textbf{presbytie}. (2) Condition : objet réel à $\overline{OA} = -\num{0,25}\ \mathrm{m}$, image virtuelle au PP : $\overline{OA'} = -\num{0,50}\ \mathrm{m}$. Alors $C = \dfrac{1}{-\num{0,50}} - \dfrac{1}{-\num{0,25}} = -2 + 4 = +2\ \delta$. (3) $C > 0$ : lentille convergente, cohérent avec une presbytie ; $2\ \delta$ est une valeur courante de lunettes de lecture.
\end{corrige}

\newpage

\coursec{Activité d'intégration}

\coursec{Leçon N08 : L'œil réduit : fonctionnement, défauts et correction}

\courssub{Objectifs de l'activité}

À l'issue de cette activité d'intégration, tu seras capable de :

\begin{itemize}
\item schématiser l'\cle{œil réduit} et y situer correctement le \cle{punctum remotum} (PR) et le \cle{punctum proximum} (PP) d'un patient sur un axe gradué ;
\item \cle{diagnostiquer} un défaut d'accommodation (myopie, hypermétropie, presbytie) en comparant les positions du PR et du PP à celles de l'œil emmétrope ;
\item calculer la \cle{vergence} du verre correcteur adapté à chaque situation (vision de loin ou vision de près) ;
\item construire la marche de rayons lumineux à travers le système \{verre correcteur + œil réduit\} pour justifier la correction ;
\item rédiger une \og ordonnance \fg{} en dioptries, comme le ferait un opticien professionnel.
\end{itemize}

\courssub{Situation}

Tu es stagiaire dans le cabinet d'optique \og Vision Claire \fg{} situé à Bonanjo, à Douala, non loin de la place du Gouvernement. Ce samedi matin, l'opticien titulaire est retenu à un séminaire à Yaoundé et te confie l'analyse préliminaire de trois fiches patients établies la veille par l'ophtalmologue du cabinet. Chaque fiche indique les positions du punctum remotum (PR, point le plus éloigné vu nettement sans accommodation) et du punctum proximum (PP, point le plus proche vu nettement avec accommodation maximale), mesurées par rapport au centre optique $O$ de l'œil réduit du patient.

\textbf{Fiche patient A — Étienne, 52 ans, comptable.} Il se plaint de devoir éloigner ses documents pour lire. Mesures : PR rejeté à l'infini ; PP à \SI{60}{cm} devant l'œil. Distance cristallin--rétine : $d = \SI{15}{mm}$ (valeur standard). Étienne souhaite lire confortablement un document placé à \SI{25}{cm}.

\textbf{Fiche patient B — M. Fotso, 34 ans, chauffeur de mototaxi.} Il plisse les yeux pour lire les panneaux routiers au loin. Mesures : PR à \SI{50}{cm} devant l'œil ; PP à \SI{12}{cm} devant l'œil. Il souhaite voir nettement les objets très éloignés (vision de loin).

\textbf{Fiche patient C — Maman Ngo Bell, 48 ans, commerçante au marché Mokolo.} Elle ne lit plus les étiquettes de prix de près, mais voit parfaitement au loin. Mesures : PR rejeté à l'infini ; PP à \SI{50}{cm} devant l'œil. Elle souhaite lire à \SI{25}{cm}.

\textbf{Données de référence — œil emmétrope (œil normal) :} PR à l'infini ; PP à environ \SI{25}{cm} (distance minimale de vision distincte $d_m = \SI{25}{cm}$). L'œil réduit est modélisé par une lentille mince convergente unique de centre optique $O$, de distance focale image variable, l'écran (rétine) étant placé à la distance fixe $d = \SI{15}{mm}$ derrière $O$.

\textbf{Ta mission :} pour chaque patient, établir le diagnostic, déterminer la vergence du verre correcteur et rédiger l'ordonnance que l'opticien signera à son retour.

\courssub{Consignes}

Le travail se fait par groupes de trois ou quatre élèves. Chaque groupe traite les trois fiches, puis la mise en commun au tableau permet une vérification croisée entre groupes.

\begin{enumerate}
\item \textbf{Schématisation.} Pour chaque patient, dessiner l'œil réduit (lentille mince de centre $O$, rétine à \SI{15}{mm}) et placer sur un axe gradué en centimètres les positions du PR et du PP. Comparer avec l'axe de l'œil emmétrope.
\item \textbf{Diagnostic.} En comparant le PR et le PP de chaque patient à ceux de l'œil emmétrope, identifier le défaut d'accommodation dont souffre chacun. Justifier en une phrase (défaut de convergence de l'œil, ou fatigue du cristallin liée à l'âge).
\item \textbf{Calcul de la correction.}
\begin{itemize}
\item Patient B : déterminer la vergence $C_B$ du verre qui ramène la vision de loin (objet à l'infini) sur la rétine, c'est-à-dire qui donne d'un objet à l'infini une image au PR du patient.
\item Patients A et C : déterminer la vergence $C_A$ puis $C_C$ du verre qui permet de lire à \SI{25}{cm}, c'est-à-dire qui donne d'un objet placé à \SI{25}{cm} une image au PP du patient.
\end{itemize}
\item \textbf{Construction graphique.} Pour le patient B, construire la marche de deux rayons issus d'un point objet $B$ à l'infini (rayons parallèles entre eux) traversant le verre correcteur puis l'œil réduit, jusqu'à la rétine. Faire de même pour le patient C avec un objet à \SI{25}{cm}.
\item \textbf{Rédaction de l'ordonnance.} Pour chaque patient, rédiger l'ordonnance en dioptries (symbole $\delta$), en précisant la nature des verres (convergents ou divergents) et leur usage (vision de loin ou vision de près).
\item \textbf{Vérification croisée.} Comparer vos résultats avec ceux des autres groupes : les vergences doivent être identiques ; discuter les éventuels écarts de construction graphique.
\end{enumerate}

\courssub{Ressources à mobiliser}

\begin{aretenir}[Formules utiles pour l'activité]
\begin{itemize}
\item \textbf{Relation de conjugaison des lentilles minces} (notions algébriques, origine au centre optique $O$) :
\[ \frac{1}{\overline{OA'}} - \frac{1}{\overline{OA}} = \frac{1}{f'} = C \]
où $C$ est la vergence en dioptries ($\delta$), $f'$ la distance focale image en mètres.
\item \textbf{Principe général de correction :} le verre correcteur doit donner, de tout objet que le patient veut voir, une \emph{image virtuelle} située dans la \cle{zone de vision distincte} du patient (entre son PP et son PR). C'est cette image qui joue le rôle d'objet pour l'œil.
\item \textbf{Correction de la myopie (vision de loin) :} verre \emph{divergent} de foyer image au PR :
\[ C = \frac{1}{\overline{OF_R}} \qquad (\overline{OF_R} < 0) \]
\item \textbf{Correction de l'hypermétropie (vision de loin) :} verre \emph{convergent} de foyer objet au PR (PR virtuel derrière l'œil, $\overline{OF_R} > 0$) :
\[ C = \frac{1}{\overline{OF_R}} \qquad (\overline{OF_R} > 0) \]
\item \textbf{Correction de la presbytie ou de l'hypermétropie (vision de près) :} verre \emph{convergent} donnant d'un objet à $d_m = \SI{25}{cm}$ une image au PP :
\[ C = \frac{1}{\overline{OF_P}} - \frac{1}{\overline{OA}} = \frac{1}{\overline{OF_P}} + \frac{1}{d_m} \]
avec $\overline{OA} = -d_m = -\SI{0,25}{m}$ et $\overline{OF_P} = -\overline{OP}$ (image virtuelle du même côté que l'objet, donc abscisse négative).
\item \textbf{Rayons particuliers pour la construction graphique :}
\begin{itemize}
\item Tout rayon passant par le centre optique $O$ n'est pas dévié.
\item Un rayon incident parallèle à l'axe optique émerge en passant par le foyer image $F'$.
\item Un rayon incident passant par le foyer objet $F$ émerge parallèle à l'axe optique.
\end{itemize}
\end{itemize}
\end{aretenir}

\begin{attention}[Pièges à éviter]
\begin{itemize}
\item Toutes les distances doivent être converties en \textbf{mètres} avant le calcul des vergences : une vergence s'exprime en dioptries ($\delta = \mathrm{m}^{-1}$).
\item Respecte scrupuleusement les \textbf{signes algébriques} : les points situés devant la lentille (côté objet) ont une abscisse négative. Ainsi $\overline{OA} = -\SI{0,25}{m}$ et non $+\SI{0,25}{m}$.
\item Le verre correcteur est accolé à l'œil : on confond le centre optique du verre avec celui $O$ de l'œil réduit.
\item Ne confonds pas le défaut et sa correction : la myopie (œil trop convergent ou trop long) se corrige par un verre \emph{divergent} ; l'hypermétropie (œil pas assez convergent ou trop court) et la presbytie (perte d'accommodation liée à l'âge) se corrigent par un verre \emph{convergent}.
\item Pour la presbytie, le PR reste à l'infini : la vision de loin ne nécessite \emph{aucune} correction.
\end{itemize}
\end{attention}

\courssub{Démarche et corrigé commenté}

\begin{exoresolu}[Le cabinet d'optique de Bonanjo : corrigé complet]
\textbf{Étape 1 — Schématisation et positionnement de PR et PP.}

Pour chaque œil, on trace l'axe optique gradué en centimètres, le centre $O$ en 0, la rétine à \SI{15}{mm} derrière $O$ (non représentée à l'échelle sur l'axe des objets, qui couvre plusieurs dizaines de centimètres).

\begin{center}
\begin{tikzpicture}[scale=0.09]
\draw[->,thick] (-70,0) -- (10,0) node[below left]{\footnotesize axe optique (cm)};
\foreach \x/\lab in {-60/{-60},-50/{-50},-25/{-25},-12/{-12},0/{0}} {
  \draw (\x,0.5) -- (\x,-0.5) node[below]{\footnotesize \lab};
}
\node[above] at (0,0.7) {\footnotesize $O$};
% Emmétrope
\node[left,popblue] at (-68,4) {\footnotesize Emmétrope};
\fill[popblue] (-25,4) circle (0.8); \node[above,popblue] at (-25,4) {\footnotesize PP};
\draw[popblue,->] (-62,4) -- (-68,4); \node[above,popblue] at (-65,4) {\footnotesize PR $\to\infty$};
% Patient A
\node[left,poporange] at (-68,7.5) {\footnotesize Patient A};
\fill[poporange] (-60,7.5) circle (0.8); \node[above,poporange] at (-60,7.5) {\footnotesize PP};
\draw[poporange,->] (-62,7.5) -- (-68,7.5); \node[above,poporange] at (-65,7.5) {\footnotesize PR $\to\infty$};
% Patient B
\node[left,popgreen] at (-68,11) {\footnotesize Patient B};
\fill[popgreen] (-12,11) circle (0.8); \node[above,popgreen] at (-12,11) {\footnotesize PP};
\fill[popgreen] (-50,11) circle (0.8); \node[above,popgreen] at (-50,11) {\footnotesize PR};
% Patient C
\node[left,poppurple] at (-68,14.5) {\footnotesize Patient C};
\fill[poppurple] (-50,14.5) circle (0.8); \node[above,poppurple] at (-50,14.5) {\footnotesize PP};
\draw[poppurple,->] (-62,14.5) -- (-68,14.5); \node[above,poppurple] at (-65,14.5) {\footnotesize PR $\to\infty$};
\end{tikzpicture}
\end{center}

\textbf{Étape 2 — Diagnostic.}

\begin{itemize}
\item \textbf{Patient A (Étienne) :} PR à l'infini (vision de loin normale) mais PP reculé à \SI{60}{cm} au lieu de \SI{25}{cm}. À 52 ans, le cristallin a perdu son élasticité : c'est la \textbf{presbytie}.
\item \textbf{Patient B (M. Fotso) :} PR rapproché à \SI{50}{cm} (au lieu de l'infini) et PP à \SI{12}{cm} : toute la zone de vision distincte est rapprochée. L'œil est trop convergent (ou trop long) : c'est la \textbf{myopie}.
\item \textbf{Patient C (Maman Ngo Bell) :} PR à l'infini, PP reculé à \SI{50}{cm}, à 48 ans : c'est également la \textbf{presbytie}.
\end{itemize}

\textbf{Étape 3 — Calcul des vergences correctrices.}

\emph{Patient B — Myopie (vision de loin).} Le verre doit donner d'un objet à l'infini une image virtuelle au PR, situé à \SI{50}{cm} devant l'œil : $\overline{OA'} = \overline{OF_R} = -\SI{0,50}{m}$ et $\dfrac{1}{\overline{OA}} = 0$ (objet à l'infini). La relation de conjugaison donne :
\[ C_B = \frac{1}{\overline{OA'}} - \frac{1}{\overline{OA}} = \frac{1}{-0,50} - 0 = -2,0\ \delta. \]
Vérification dimensionnelle : l'inverse d'une longueur en mètres donne bien des $\mathrm{m}^{-1}$, c'est-à-dire des dioptries. Le verre est \textbf{divergent} ($C_B < 0$), de distance focale $f' = \dfrac{1}{C_B} = -\SI{0,50}{m}$.

\emph{Patient A — Presbytie (lecture à \SI{25}{cm}).} Le verre doit donner d'un objet $A$ à \SI{25}{cm} ($\overline{OA} = -\SI{0,25}{m}$) une image virtuelle au PP ($\overline{OA'} = -\SI{0,60}{m}$) :
\[ C_A = \frac{1}{-0,60} - \frac{1}{-0,25} = -1,67 + 4,00 = +2,33\ \delta. \]
Le verre est \textbf{convergent}, de focale $f' = \dfrac{1}{C_A} \approx \SI{0,43}{m}$.

\emph{Patient C — Presbytie (lecture à \SI{25}{cm}).} Même calcul avec $\overline{OA'} = -\SI{0,50}{m}$ :
\[ C_C = \frac{1}{-0,50} - \frac{1}{-0,25} = -2,00 + 4,00 = +2,0\ \delta. \]
Verre \textbf{convergent} de focale $f' = \SI{0,50}{m}$.

\textbf{Étape 4 — Constructions graphiques.}

\emph{Patient B :} les rayons parallèles issus de l'objet à l'infini émergent du verre divergent comme s'ils provenaient du foyer image du verre, placé au PR ; l'œil (sans accommoder) les conjugue alors exactement sur la rétine.

\begin{popfigure}
\begin{center}
\begin{tikzpicture}[scale=0.95]
\draw[->,thick] (-6.5,0) -- (4,0) node[below left]{\footnotesize axe optique};
% Verre divergent
\draw[poporange,thick,<->] (-2,-1.8) -- (-2,1.8);
\node[poporange,above] at (-2,1.9) {\footnotesize Verre divergent ($C_B=-2\,\delta$)};
% Œil réduit
\draw[popblue,thick,<->] (0,-1.8) -- (0,1.8);
\node[popblue,above] at (0,1.9) {\footnotesize Œil réduit};
\node[below] at (0,-0.1) {\footnotesize $O$};
% Rétine
\draw[popdark,thick] (1.5,-1.4) -- (1.5,1.4);
\node[popdark,below] at (1.5,-1.5) {\footnotesize Rétine};
% PR = F' du verre
\fill[poporange] (-4,0) circle (0.07); \node[below,poporange] at (-4,0) {\footnotesize PR $= F'_{\text{verre}}$};
% Rayons incidents parallèles
\draw[popgreen,thick,->] (-6.5,1.4) -- (-2,1.4);
\draw[popgreen,thick,->] (-6.5,0.7) -- (-2,0.7);
% Rayons émergents divergeant depuis PR (pointillés pour prolongement virtuel)
\draw[popgreen,dashed] (-4,0) -- (-2,1.4);
\draw[popgreen,dashed] (-4,0) -- (-2,0.7);
\draw[popgreen,thick,->] (-2,1.4) -- (0,0.6);
\draw[popgreen,thick,->] (-2,0.7) -- (0,0.3);
% Convergence par l'œil sur la rétine
\draw[popgreen,thick,->] (0,0.6) -- (1.5,0);
\draw[popgreen,thick,->] (0,0.3) -- (1.5,0);
\fill[popink] (1.5,0) circle (0.07); \node[above right,popink] at (1.5,0) {\footnotesize $B'$};
\end{tikzpicture}
\end{center}
\legende{Correction de la myopie : le verre divergent donne de l'objet à l'infini une image virtuelle au PR ; l'œil (au repos) la conjugue sur la rétine.}
\end{popfigure}

\emph{Patient C :} l'objet $A$ à \SI{25}{cm} est placé entre le foyer objet et le verre convergent ; le verre en donne une image virtuelle $A'$ au PP (\SI{50}{cm}), que l'œil voit nettement en accommodant au maximum.

\begin{popfigure}
\begin{center}
\begin{tikzpicture}[scale=0.6]
\draw[->,thick] (-12,0) -- (4.5,0) node[below left]{\footnotesize axe optique};
% Verre convergent
\draw[poporange,thick,<->] (0,-2.5) -- (0,2.5);
\node[poporange,above] at (0,2.6) {\footnotesize Verre convergent ($C_C=+2\,\delta$)};
% Œil réduit accolé
\draw[popblue,thick,<->] (0.4,-2.5) -- (0.4,2.5);
\node[popblue,below] at (0.4,-2.6) {\footnotesize Œil réduit};
% Rétine
\draw[popdark,thick] (2.5,-1.6) -- (2.5,1.6);
\node[popdark,below] at (2.5,-1.7) {\footnotesize Rétine};
% Objet A à -25 cm (échelle : -5), Image A' au PP à -50 cm (échelle : -10)
\draw[popink,thick,->] (-5,0) -- (-5,1.8) node[above]{\footnotesize $A$ (\SI{25}{cm})};
\draw[poppurple,thick,dashed,->] (-10,0) -- (-10,3.5) node[above]{\footnotesize $A'$ au PP (\SI{50}{cm})};
% Rayon 1 : passant par le centre O (non dévié)
\draw[popgreen,thick] (-5,1.8) -- (2.5,0);
\draw[popgreen,dashed] (-10,3.5) -- (-5,1.8);
% Rayon 2 : incident parallèle à l'axe, émerge par F' image
\draw[popgold,thick] (-5,1.8) -- (0,1.8);
\draw[popgold,thick,->] (0,1.8) -- (2.5,0);
\draw[popgold,dashed] (0,1.8) -- (-10,3.5);
% Image finale sur rétine
\fill[popink] (2.5,0) circle (0.09); \node[above right,popink] at (2.5,0.1) {\footnotesize Image nette sur la rétine};
\end{tikzpicture}
\end{center}
\legende{Correction de la presbytie : le verre convergent donne de l'objet proche une image virtuelle reculée au PP, vue nettement par l'œil en accommodation maximale.}
\end{popfigure}

\textbf{Étape 5 — Ordonnances.}

\begin{center}
\begin{tabularx}{\linewidth}{|l|X|X|X|}
\hline
\rowcolor{popblueL} \textbf{Patient} & \textbf{Diagnostic} & \textbf{Vergence} & \textbf{Prescription} \\
\hline
A — Étienne & Presbytie & $C_A = +2,33\ \delta$ & Verres convergents pour la lecture à \SI{25}{cm} \\
\hline
B — M. Fotso & Myopie & $C_B = -2,0\ \delta$ & Verres divergents portés en permanence (vision de loin) \\
\hline
C — Maman Ngo Bell & Presbytie & $C_C = +2,0\ \delta$ & Verres convergents pour la lecture à \SI{25}{cm} \\
\hline
\end{tabularx}
\end{center}

\textbf{Vérification croisée.} Tous les groupes doivent trouver les mêmes vergences (le calcul est unique). On remarque que le patient A, dont le PP est le plus reculé (\SI{60}{cm} contre \SI{50}{cm} pour C), nécessite la vergence la plus forte ($+2,33\,\delta$ contre $+2,00\,\delta$) : plus la presbytie est avancée, plus le verre correcteur est convergent.
\end{exoresolu}

\courssub{Bilan de l'activité}

Cette activité t'a permis de mettre en évidence les points essentiels de la leçon :

\begin{itemize}
\item Le \textbf{diagnostic} d'un défaut visuel repose entièrement sur la comparaison du couple (PR, PP) du patient avec celui de l'œil emmétrope :
  \begin{itemize}
  \item PR fini (rapproché) $\Rightarrow$ \textbf{myopie} (œil trop convergent ou trop long) ;
  \item PR à l'infini et PP reculé au-delà de \SI{25}{cm} chez un sujet de plus de 40--45 ans $\Rightarrow$ \textbf{presbytie} (perte d'élasticité du cristallin, amplitude d'accommodation réduite) ;
  \item PR virtuel (derrière l'œil) $\Rightarrow$ \textbf{hypermétropie} (œil pas assez convergent ou trop court) — non illustré ici mais traité dans le cours.
  \end{itemize}
\item La \textbf{correction} consiste toujours à faire en sorte que le verre donne de l'objet regardé une \emph{image virtuelle située dans la zone de vision distincte} du patient :
  \begin{itemize}
  \item Au PR pour la vision de loin (myopie, hypermétropie) ;
  \item Au PP pour la vision de près (presbytie, hypermétropie).
  \end{itemize}
\item Le calcul de la vergence n'est qu'une application directe de la \textbf{relation de conjugaison des lentilles minces}, à condition de respecter les signes algébriques (convention de Descartes : abscisses algébriques, origine au centre optique, sens de la lumière incident = positif) et de convertir les distances en mètres.
\item La \textbf{construction graphique} confirme le calcul et donne du sens physique : verre divergent pour écarter les rayons (myopie), verre convergent pour les rapprocher (presbytie, hypermétropie).
\item Enfin, tu as découvert le geste professionnel de l'opticien : traduire un diagnostic en une \textbf{ordonnance chiffrée en dioptries}, compétence directement mobilisable dans la vie courante — et au Probatoire.
\end{itemize}

\newpage

\coursec{Méthodes et exercices résolus}

\courssub{Méthodes et exercices résolus}

\begin{methode}[Tracer le schéma de l'œil réduit et identifier ses éléments]
\textbf{Objectif :} Représenter fidèlement le modèle de l'œil réduit de Gullstrand (ou Emsley) utilisé en Première.
\begin{enumerate}
    \item \textbf{Dessiner le système optique :} Un unique dioptre sphérique (ou une lentille mince équivalente) de centre $O$ (centre optique), séparant l'air ($n_1 = 1$) de l'humeur vitrée ($n_2 = 1,33$ ou $4/3$). Tracer l'axe optique horizontal.
    \item \textbf{Placer les points cardinaux :}
    \begin{itemize}
        \item Foyer objet $F$ (dans l'air) et foyer image $F'$ (dans l'œil). Rappel : $f' = OF' = -n_2 \times OF = -n_2 f$ (signes algébriques respectés : $f<0$, $f'>0$).
        \item Points principaux $H \equiv H' \equiv O$ (pour un dioptre unique ou lentille mince centrée en $O$).
    \end{itemize}
    \item \textbf{Repérer la rétine :} Plan image fixe, situé à une distance $OF' \approx 22,5\,\text{mm}$ (valeur type) ou donnée par l'énoncé. C'est l'écran de réception.
    \item \textbf{Indiquer le cristallin :} Représenté par une lentille biconvexe en $O$, symbole de la puissance variable $C_{\text{crist}}$.
    \item \textbf{Légender :} Distances algébriques ($p, p'$), vergence $C = \dfrac{n_2}{p'} - \dfrac{n_1}{p}$, puissance totale $C_{\text{œil}} = C_{\text{cornée}} + C_{\text{crist}}$.
\end{enumerate}
\cle{Réflexe :} Toujours orienter l'axe optique de gauche (objet) vers droite (image). L'objet réel est à gauche ($p < 0$), l'image réelle sur la rétine est à droite ($p' > 0$).
\end{methode}

\begin{exoresolu}[Schéma et caractéristiques de l'œil réduit]
\textbf{Énoncé :} On considère le modèle de l'œil réduit de Gullstrand : un unique dioptre sphérique de centre $O$, de rayon de courbure $R = 5,0\,\text{mm}$, séparant l'air ($n_1 = 1,00$) de l'humeur vitrée ($n_2 = 1,33$). La rétine est située au foyer image $F'$.
\begin{enumerate}
    \item Tracer le schéma annoté de cet œil réduit. Placer $O$, $F$, $F'$, la rétine, l'axe optique. Indiquer le sens de la lumière.
    \item Calculer les distances algébriques $OF$ et $OF'$ (en mm).
    \item Déduire la puissance $C$ de ce dioptre (en dioptries).
    \item Un objet $AB$ de hauteur $1,0\,\text{cm}$ est placé à $25\,\text{cm}$ de $O$. Déterminer par le calcul la position et la taille de son image $A'B'$ sur la rétine. L'image est-elle nette ?
\end{enumerate}

\tcblower
\textbf{Solution détaillée :}
\begin{enumerate}
    \item \textbf{Schéma :} (Voir figure ci-dessous). L'axe optique est horizontal. La lumière vient de la gauche. $O$ est le sommet du dioptre. $F$ est à gauche de $O$ (foyer objet), $F'$ à droite (foyer image). La rétine coïncide avec le plan focal image.
    \begin{popfigure}
    \begin{tikzpicture}[scale=1.2, >=Stealth]
        % Axes
        \draw[popline, thick, <->] (-3.5,0) -- (3.5,0) node[right] {Axe optique};
        % Dioptre
        \draw[popblue, very thick] (0,-1.2) -- (0,1.2);
        \draw[popblue, thick] (0,1.2) arc (90:-90:0.3 and 1.2); % Courbure suggérée
        \node[below left] at (0,0) {$O$};
        % Milieux
        \node[popmuted] at (-2,1) {Air $n_1=1$};
        \node[popmuted] at (2,1) {Humeur vitrée $n_2=1,33$};
        % Foyers
        \coordinate (F) at (-1.515,0);
        \coordinate (Fp) at (2.015,0);
        \draw[poporange, thick] (F) -- ++(0,0.3) node[above] {$F$};
        \draw[poporange, thick] (Fp) -- ++(0,0.3) node[above] {$F'$};
        % Rétine
        \draw[popred, very thick] (2.015,-1) -- (2.015,1) node[midway, right=2mm] {Rétine};
        % Flèches lumière
        \draw[popgreen, thick, ->] (-3, -0.5) -- (-0.5, -0.5);
        \draw[popgreen, thick, ->] (0.5, -0.5) -- (2, -0.5);
        % Objet AB (pour question 4)
        \draw[popdark, thick] (-2.5,0.5) -- (-2.5,0) node[below] {$A$} node[above] {$B$};
        \node[left] at (-2.5,0.6) {$AB$};
        % Construction image (simplifiée pour objet à 25 cm -> image derrière rétine)
        \draw[dashed, poppurple] (-2.5,0.5) -- (0,0.5) -- (2.15, -0.15); % Rayon parallèle -> F'
        \draw[dashed, poppurple] (-2.5,0.5) -- (0,0) -- (2.15, -0.15); % Rayon centre O
        % Image réelle (derrière la rétine)
        \draw[popdark, thick] (2.15, -0.15) -- (2.15, 0) node[below] {$A'$} node[above] {$B'$};
        \node[above right] at (2.15, 0.2) {Image floue};
    \end{tikzpicture}
    \legende{Modèle de l'œil réduit de Gullstrand (dioptre unique). Objet à 25 cm : image formée derrière la rétine (floue sans accommodation).}
    \end{popfigure}

    \item \textbf{Calcul des distances focales :}
    Formules de base pour un dioptre sphérique :
    \[
    OF = f = -\frac{n_1 R}{n_2 - n_1} \quad \text{et} \quad OF' = f' = \frac{n_2 R}{n_2 - n_1}
    \]
    Avec $R = +5,0\,\text{mm}$ (centre de courbure à droite du sommet, dioptre convergent), $n_1=1$, $n_2=1,33$.
    \[
    f = -\frac{1,00 \times 5,0}{1,33 - 1,00} = -\frac{5,0}{0,33} \approx -15,15\,\text{mm}
    \]
    \[
    f' = \frac{1,33 \times 5,0}{0,33} = \frac{6,65}{0,33} \approx 20,15\,\text{mm}
    \]
    \textbf{Résultat :} $OF \approx -15,2\,\text{mm}$ ; $OF' \approx 20,2\,\text{mm}$.
    \textit{Vérification dimensionnelle :} $[nR/(n_2-n_1)] = \text{mm}$. Cohérent. Relation $f' = -n_2 f$ vérifiée : $1,33 \times 15,15 \approx 20,15$.

    \item \textbf{Puissance du dioptre :}
    \[
    C = \frac{n_2 - n_1}{R} = \frac{0,33}{5,0 \times 10^{-3}} = 66,0\,\text{D}
    \]
    Ou $C = n_1/f = 1/(-0,01515) \approx -66\,\text{D}$ (puissance objet) et $C = n_2/f' = 1,33/0,02015 \approx 66\,\text{D}$ (puissance image).
    \textbf{Résultat :} $C = 66,0\,\text{D}$.

    \item \textbf{Imagerie d'un objet à 25 cm :}
    Données : $p = OA = -25\,\text{cm} = -0,250\,\text{m}$ (objet réel à gauche). $p' = OA'$ inconnu.
    Relation de conjugaison (formule de Descartes pour dioptre) :
    \[
    \frac{n_2}{p'} - \frac{n_1}{p} = C = \frac{n_2 - n_1}{R}
    \]
    \[
    \frac{1,33}{p'} - \frac{1,00}{-0,250} = 66,0
    \]
    \[
    \frac{1,33}{p'} + 4,00 = 66,0 \quad \Rightarrow \quad \frac{1,33}{p'} = 62,0
    \]
    \[
    p' = \frac{1,33}{62,0} \approx 0,02145\,\text{m} = 21,45\,\text{mm}
    \]
    La rétine est située en $F'$ soit $p'_{\text{rétine}} = 20,15\,\text{mm}$.
    L'image se forme à $21,45\,\text{mm}$, soit \textbf{derrière la rétine}.
    Grandissement transverse : $g = \dfrac{n_1 p'}{n_2 p} = \dfrac{1,00 \times 0,02145}{1,33 \times (-0,250)} \approx -0,0645$.
    Taille image : $A'B' = |g| \times AB = 0,0645 \times 10\,\text{mm} \approx 0,65\,\text{mm}$.
    \textbf{Conclusion :} L'image est réelle, inversée, réduite ($\approx 0,65\,\text{mm}$), formée \emph{derrière la rétine}. L'image n'est \textbf{pas nette} sur la rétine. Pour la voir nette, l'œil doit \textbf{accommoder} (augmenter sa puissance d'environ $4\,\text{D}$) pour ramener l'image sur la rétine. C'est la vision de près.
\end{enumerate}
\end{exoresolu}

\begin{methode}[Analyser l'accommodation et déterminer PP et PR]
\textbf{Objectif :} Relier la variation de puissance du cristallin aux points limites de la vision nette.
\begin{enumerate}
    \item \textbf{Rappeler le mécanisme :} L'accommodation est la variation de la puissance du cristallin ($C_{\text{crist}}$) par déformation (zonule de Zinn / muscle ciliaire). Puissance totale de l'œil : $C_{\text{œil}} = C_{\text{cornée}} + C_{\text{crist}}$.
    \item \textbf{Identifier les deux états extrêmes :}
    \begin{itemize}
        \item \textbf{Œil au repos (accommodation minimale) :} $C_{\text{œil}} = C_{\text{min}}$. L'image d'un objet à l'infini se forme sur la rétine. Le point image de l'objet à l'infini est le \textbf{Punctum Remotum (PR)}. \emph{Par définition, PR est le point conjugué de la rétine pour l'œil au repos.}
        \item \textbf{Œil accommodé au maximum :} $C_{\text{œil}} = C_{\text{max}} = C_{\text{min}} + A_{\text{max}}$ ($A_{\text{max}}$ : amplitude d'accommodation). L'image de l'objet le plus proche net se forme sur la rétine. Ce point objet est le \textbf{Punctum Proximum (PP)}. \emph{Par définition, PP est le point conjugué de la rétine pour l'œil accommodé au maximum.}
    \end{itemize}
    \item \textbf{Formules clés (vergences en dioptries, distances en mètres) :}
    On travaille en \textbf{vergences air} ($n=1$) pour la correction. La vergence de la rétine est $V_{\text{rétine}} = \dfrac{n'}{p'}$ ($n'=1,33$, $p'$ distance rétine en m).
    \begin{align*}
        C_{\text{min}} &= V_{\text{rétine}} - V_{PR} \\
        C_{\text{max}} &= V_{\text{rétine}} - V_{PP} \\
        A_{\text{max}} &= C_{\text{max}} - C_{\text{min}} = V_{PP} - V_{PR}
    \end{align*}
    \item \textbf{Convention de signe (algébrique) :} Objet réel à gauche $\Rightarrow$ distance $p < 0$, vergence $V = 1/p < 0$. Image réelle à droite $\Rightarrow$ $p' > 0$, $V' > 0$.
    \begin{itemize}
        \item Emmétrope : $PR = -\infty$ ($V_{PR}=0$), $PP$ fini (ex: $-25\,\text{cm}$).
        \item Myope : $PR$ fini et \textbf{négatif} (devant l'œil), $PP$ plus proche que l'emmétrope.
        \item Hypermétrope : $PR$ \textbf{positif} (virtuel, derrière l'œil), $PP$ plus loin que l'emmétrope.
    \end{itemize}
\end{enumerate}
\cle{Attention :} $PP$ et $PR$ sont des \textbf{points objets} (côté entrée de la lumière). Leurs coordonnées algébriques $x_{PP}, x_{PR}$ sont \textbf{négatives} s'ils sont réels (devant l'œil), \textbf{positives} s'ils sont virtuels (derrière l'œil). Ne pas confondre avec les distances focales image.
\end{methode}

\begin{exoresolu}[Détermination de PP, PR et amplitude d'accommodation]
\textbf{Énoncé :} Un œil réduit a une distance rétine-centre optique $d = 22,0\,\text{mm}$. Sa puissance minimale (œil au repos) est $C_{\text{min}} = 58,0\,\text{D}$. Son amplitude d'accommodation maximale est $A_{\text{max}} = 4,0\,\text{D}$.
\begin{enumerate}
    \item Calculer la vergence de la rétine $V_{\text{rétine}}$.
    \item Déterminer la position du Punctum Remotum $PR$ (en cm, signe algébrique). L'œil est-il emmétrope, myope ou hypermétrope ?
    \item Calculer la puissance maximale $C_{\text{max}}$.
    \item Déterminer la position du Punctum Proximum $PP$ (en cm, signe algébrique).
    \item Quelle est la distance de vision nette de cet œil (intervalle de distances d'objets vus nets sans correction) ?
\end{enumerate}

\tcblower
\textbf{Solution détaillée :}
\begin{enumerate}
    \item \textbf{Vergence rétine :} L'image se forme sur la rétine : $p' = +22,0\,\text{mm} = +0,0220\,\text{m}$. Milieu vitré $n' = 1,33$.
    $V_{\text{rétine}} = \dfrac{n'}{p'} = \dfrac{1,33}{0,0220} \approx 60,45\,\text{D}$.
    \item \textbf{Position du PR :} Œil au repos : $C_{\text{min}} = V_{\text{rétine}} - V_{PR}$.
    \[
    58,0 = 60,45 - \frac{1}{PR}
    \]
    \[
    \frac{1}{PR} = 60,45 - 58,0 = 2,45\,\text{D}
    \]
    \[
    PR = \frac{1}{2,45} \approx 0,408\,\text{m} = +40,8\,\text{cm}
    \]
    \textbf{Analyse du signe :} $1/PR = +2,45 \Rightarrow PR = +0,408\,\text{m}$.
    \textbf{Interprétation :} $PR > 0$ signifie que le point conjugué de la rétine (pour l'œil au repos) est un \textbf{objet virtuel situé derrière l'œil} (côté droit). L'œil au repos converge trop fort : il forme l'image d'un objet à l'infini \emph{devant} la rétine. Pour former l'image sur la rétine, il lui faut un objet virtuel (convergent) derrière lui.
    \textbf{Diagnostic :} \textbf{Hypermétropie}. L'œil manque de puissance (ou est trop court) par rapport à l'emmétropie. $PR$ est virtuel ($+40,8\,\text{cm}$).
    \item \textbf{Puissance maximale :} $C_{\text{max}} = C_{\text{min}} + A_{\text{max}} = 58,0 + 4,0 = 62,0\,\text{D}$.
    \item \textbf{Position du PP :} Œil accommodé max : $C_{\text{max}} = V_{\text{rétine}} - V_{PP}$.
    \[
    62,0 = 60,45 - \frac{1}{PP}
    \]
    \[
    \frac{1}{PP} = 60,45 - 62,0 = -1,55\,\text{D}
    \]
    \[
    PP = \frac{1}{-1,55} \approx -0,645\,\text{m} = -64,5\,\text{cm}
    \]
    \textbf{Interprétation :} $PP < 0$. Le point le plus proche vu net est un \textbf{objet réel} situé à $64,5\,\text{cm}$ devant l'œil. C'est loin (hypermétropie forte ou presbytie débutante).
    \item \textbf{Intervalle de vision nette :} De $PP = -64,5\,\text{cm}$ (objet réel le plus proche) à $PR = +40,8\,\text{cm}$ (objet virtuel le plus loin = infini non accessible sans accommodation).
    \textit{En pratique sans correction :} L'œil ne peut voir net que des objets situés entre $64,5\,\text{cm}$ et... l'infini n'est pas accessible (PR virtuel). Il doit accommoder en permanence. La vision nette \emph{sans effort} (au repos) correspond à un objet virtuel à $+40,8\,\text{cm}$ (impossible physiquement sans lentille). Donc \textbf{vision floue de loin et de près sans correction}.
\end{enumerate}
\end{exoresolu}

\begin{methode}[Corriger la myopie par une lentille divergente]
\textbf{Objectif :} Calculer la puissance de la lentille correctrice (lunettes) pour ramener le PR de l'œil myope à l'infini.
\begin{enumerate}
    \item \textbf{Diagnostiquer :} Myopie $\Leftrightarrow$ $PR$ \textbf{réel et négatif} (devant l'œil, ex $PR = -2,0\,\text{m}$). L'œil voit net de près, flou de loin. Puissance max trop forte ou axe trop long.
    \item \textbf{Principe de correction :} Placer une lentille mince (lunettes) devant l'œil (distance sommet-cornée $d \approx 0$ en 1ère C, sinon tenir compte du vertex). La lentille doit transformer un objet à l'infini ($V_{\text{obj}} = 0$) en une image \textbf{au PR de l'œil}.
    \item \textbf{Calcul de la puissance $C_L$ (vergences dans l'air, $n=1$) :}
    \[
    C_L = V'_{\text{lentille}} - V_{\text{obj}} = V_{PR} - 0 = \frac{1}{PR}
    \]
    Comme $PR < 0$, $C_L < 0$ : \textbf{lentille divergente}.
    \item \textbf{Vérification vision de près :} Avec la correction, le nouveau $PP'$ (point proche corrigé) est l'image par la lentille du $PP$ naturel.
    \[
    V_{PP'} = C_L + V_{PP} \quad \Rightarrow \quad PP' = \frac{1}{C_L + 1/PP}
    \]
    Il faut vérifier que $PP'$ reste accessible (généralement $< 25\,\text{cm}$ pour lecture).
\end{enumerate}
\cle{Formule raccourci (si $d=0$) :} $C_L = \dfrac{1}{PR}$ (en m, résultat en D). $PR$ en mètres, signe algébrique gardé.
\end{methode}

\begin{exoresolu}[Correction d'un œil myope]
\textbf{Énoncé :} Un œil myope a son Punctum Remotum en $PR = -1,25\,\text{m}$ et son Punctum Proximum en $PP = -10,0\,\text{cm}$. On place une lentille correctrice au contact de l'œil ($d=0$).
\begin{enumerate}
    \item Déterminer la nature et la puissance de la lentille correctrice pour que cet œil voie net à l'infini.
    \item Calculer la nouvelle position du Punctum Proximum $PP'$ avec cette correction.
    \item L'œil corrigé peut-il lire un journal placé à $25\,\text{cm}$ ?
\end{enumerate}

\tcblower
\textbf{Solution détaillée :}
\begin{enumerate}
    \item \textbf{Correction de loin :} On veut $V_{\text{obj}} = 0$ (objet à l'infini) $\to$ Image au $PR$ de l'œil.
    $V_{PR} = \dfrac{1}{PR} = \dfrac{1}{-1,25} = -0,80\,\text{D}$.
    Puissance lentille : $C_L = V_{PR} - V_{\text{obj}} = -0,80 - 0 = -0,80\,\text{D}$.
    \textbf{Nature :} Lentille \textbf{divergente} (puissance négative).
    \textbf{Résultat :} $C_L = -0,80\,\text{D}$.
    \item \textbf{Nouveau Punctum Proximum $PP'$ :} C'est l'image de l'ancien $PP$ par la lentille.
    $V_{PP} = \dfrac{1}{PP} = \dfrac{1}{-0,10} = -10,0\,\text{D}$.
    Relation lentille : $V_{PP'} = C_L + V_{PP} = -0,80 + (-10,0) = -10,80\,\text{D}$.
    $PP' = \dfrac{1}{V_{PP'}} = \dfrac{1}{-10,80} \approx -0,0926\,\text{m} = -9,26\,\text{cm}$.
    \textbf{Résultat :} $PP' \approx -9,3\,\text{cm}$.
    \item \textbf{Lecture à 25 cm :} L'objet est à $p = -0,25\,\text{m}$ ($V = -4,0\,\text{D}$).
    Vergence sortant de la lentille : $V' = C_L + V = -0,80 + (-4,0) = -4,80\,\text{D}$.
    Cette vergence correspond à un objet virtuel pour l'œil à $p' = 1/V' \approx -0,208\,\text{m} = -20,8\,\text{cm}$.
    L'œil nu voit net les objets entre $PP=-10\,\text{cm}$ et $PR=-1,25\,\text{m}$.
    $-20,8\,\text{cm}$ est bien dans cet intervalle ($-125 < -20,8 < -10$).
    \textbf{Conclusion :} Oui, l'œil corrigé voit net à $25\,\text{cm}$ (son $PP'$ est à $9,3\,\text{cm}$).
\end{enumerate}
\end{exoresolu}

\begin{methode}[Corriger l'hypermétropie par une lentille convergente]
\textbf{Objectif :} Calculer la puissance de la lentille correctrice pour un œil hypermétrope ($PR$ virtuel, positif).
\begin{enumerate}
    \item \textbf{Diagnostiquer :} Hypermétropie $\Leftrightarrow$ $PR > 0$ (virtuel, derrière l'œil). L'œil au repos forme l'image de l'infini \emph{derrière} la rétine. Il doit accommoder pour voir loin, et encore plus pour voir près.
    \item \textbf{Principe :} La lentille doit créer une image virtuelle de l'objet à l'infini \textbf{au $PR$ de l'œil} (côté droit, $p'_{L} = PR > 0$).
    \item \textbf{Calcul (vergences air, $d=0$) :}
    Objet à l'infini : $V_{\text{obj}} = 0$.
    Image par lentille : $V'_{\text{lentille}} = \dfrac{1}{p'_{L}} = \dfrac{1}{PR} > 0$.
    $C_L = V'_{\text{lentille}} - V_{\text{obj}} = \dfrac{1}{PR}$.
    Comme $PR > 0$, $C_L > 0$ : \textbf{lentille convergente}.
    \item \textbf{Vision de près corrigée :} L'objet à $25\,\text{cm}$ ($V=-4\,\text{D}$) donne une image par la lentille à $V' = C_L - 4$. Cette image devient l'objet pour l'œil. Il faut que cet objet soit entre $PP$ et $PR$ de l'œil nu.
    \item \textbf{Amplitude d'accommodation résiduelle :} $A_{\text{max}} = V_{PP} - V_{PR}$ (inchangée). La correction ne change pas $A_{\text{max}}$, elle décale l'intervalle.
\end{enumerate}
\cle{Piège :} Ne pas confondre $PR$ (point objet pour l'œil nu) et le foyer image de la lentille. Ici $F'_L = PR$.
\end{methode}

\begin{exoresolu}[Correction d'un œil hypermétrope]
\textbf{Énoncé :} Un œil hypermétrope présente un Punctum Remotum $PR = +50,0\,\text{cm}$ et un Punctum Proximum $PP = -80,0\,\text{cm}$. Distance sommet-cornée nulle.
\begin{enumerate}
    \item Déterminer la puissance de la lentille correctrice pour la vision de loin. Préciser sa nature.
    \item Calculer l'amplitude d'accommodation $A_{\text{max}}$ de cet œil.
    \item Déterminer la position du nouveau Punctum Proximum $PP'$ avec la correction.
    \item Cet œil corrigé peut-il lire à $25\,\text{cm}$ sans effort excessif ?
\end{enumerate}

\tcblower
\textbf{Solution détaillée :}
\begin{enumerate}
    \item \textbf{Correction de loin :} $PR = +0,50\,\text{m}$.
    $C_L = \dfrac{1}{PR} = \dfrac{1}{0,50} = +2,00\,\text{D}$.
    \textbf{Nature :} Lentille \textbf{convergente} ($+2,00\,\text{D}$).
    \item \textbf{Amplitude d'accommodation :}
    $V_{PR} = 1/0,50 = +2,00\,\text{D}$.
    $V_{PP} = 1/(-0,80) = -1,25\,\text{D}$.
    $A_{\text{max}} = V_{PP} - V_{PR} = -1,25 - (+2,00) = -3,25\,\text{D}$.
    L'amplitude est la valeur absolue : $A_{\text{max}} = 3,25\,\text{D}$.
    \item \textbf{Nouveau PP' (vision de près corrigée) :} C'est la position de l'objet réel qui, après la lentille, donne une image au $PP$ de l'œil nu ($V_{PP} = -1,25\,\text{D}$).
    $V_{PP'} = V_{PP} - C_L = -1,25 - 2,00 = -3,25\,\text{D}$.
    $PP' = 1/(-3,25) \approx -0,308\,\text{m} = -30,8\,\text{cm}$.
    \item \textbf{Lecture à 25 cm :} L'objet à $25\,\text{cm}$ ($V_{\text{obj}}=-4,00\,\text{D}$) donne une image par la lentille :
    $V_{\text{inter}} = C_L + V_{\text{obj}} = +2,00 - 4,00 = -2,00\,\text{D}$.
    Cet objet virtuel pour l'œil est à $50\,\text{cm}$ devant lui.
    L'œil nu voit net pour des vergences objets comprises entre $V_{PP}=-1,25\,\text{D}$ et $V_{PR}=+2,50\,\text{D}$.
    $V_{\text{inter}} = -2,00\,\text{D}$ est \textbf{inférieur à $V_{PP}$} (plus négatif), donc \textbf{en dehors de l'intervalle} (objet trop proche pour l'œil nu).
    Accommodation requise : $\Delta C = V_{PR} - V_{\text{inter}} = 2,50 - (-2,00) = 4,50\,\text{D}$.
    Comme $\Delta C = 4,50\,\text{D} > A_{\text{max}} = 3,25\,\text{D}$, \textbf{la lecture à 25 cm est impossible} sans effort excessif (dépassement de $1,25\,\text{D}$).
    \textit{Note :} Le $PP'$ corrigé est à $30,8\,\text{cm}$. L'objet à $25\,\text{cm}$ est \emph{plus proche} que $PP'$, l'œil ne peut pas accommoder suffisamment.
\end{enumerate}
\end{exoresolu}

\begin{methode}[Corriger la presbytie : addition pour la vision de près]
\textbf{Objectif :} Traiter la perte d'accommodation liée à l'âge ($A_{\text{max}}$ insuffisante pour la lecture à $25\,\text{cm}$).
\begin{enumerate}
    \item \textbf{Identifier le contexte :} Sujet âgé ($>45$ ans), $A_{\text{max}}$ faible (ex $< 3\,\text{D}$). Vision de loin souvent bonne (emmétropie ou myopie/hypermétropie corrigée). Problème : $PP$ trop loin (ex $PP = -50\,\text{cm}$), impossible de lire à $25\,\text{cm}$.
    \item \textbf{Principe :} Ajouter une lentille convergente (loupe) \textbf{uniquement pour la vision de près} (lunettes de lecture, verres progressifs, ou addition sur correction de loin).
    \item \textbf{Calcul de l'addition $Add$ (puissance de la loupe) :}
    On veut que l'objet à $25\,\text{cm}$ ($V_{\text{obj}} = -4,00\,\text{D}$) soit vu net par l'œil \emph{sans accommodation} (ou avec accommodation résiduelle confortable).
    Cas standard (vision de près sans effort, accommodation = 0) : La loupe doit former une image de l'objet à $25\,\text{cm}$ \textbf{au $PP$ naturel de l'œil} (point le plus proche vu net sans effort max).
    \[
    Add = V_{PP} - V_{\text{obj}} = \frac{1}{PP} - (-4,00) = \frac{1}{PP} + 4,00
    \]
    (Avec $PP$ en mètres, négatif).
    \item \textbf{Si l'œil a aussi un défaut de loin (myopie/hypermétropie) :}
    \begin{itemize}
        \item Corriger d'abord la vision de loin (lentille $C_L$).
        \item Calculer l'addition $Add$ par rapport au $PP$ de l'œil \emph{corrigé de loin} ($PP'$).
        \item Verre progressif : Partie haute $C_L$, partie basse $C_L + Add$.
    \end{itemize}
\end{enumerate}
\cle{Repère :} Emmétrope presbyte : $PP = -1/A_{\text{max}}$. Addition $Add \approx 4,00 - A_{\text{max}}$ (si lecture à 25 cm).
\end{methode}

\begin{exoresolu}[Presbytie et correction par addition]
\textbf{Énoncé :} Un sujet de 55 ans, emmétrope, a une amplitude d'accommodation résiduelle $A_{\text{max}} = 1,50\,\text{D}$. Il souhaite lire à $25\,\text{cm}$.
\begin{enumerate}
    \item Déterminer son Punctum Proximum $PP$ naturel.
    \item Calculer l'addition $Add$ nécessaire pour lire à $25\,\text{cm}$ sans effort d'accommodation (œil au repos).
    \item S'il accepte d'accommoder $0,50\,\text{D}$ pour la lecture, quelle addition choisir ?
    \item Ce sujet devient myope plus tard ($PR = -2,00\,\text{m}$). Quelle correction pour la vision de loin ? Quelle addition pour la lecture (verres progressifs) ?
\end{enumerate}

\tcblower
\textbf{Solution détaillée :}
\begin{enumerate}
    \item \textbf{PP naturel :} Emmétrope $\Rightarrow PR = -\infty$ ($V_{PR}=0$).
    $A_{\text{max}} = V_{PP} - V_{PR} = V_{PP}$.
    $V_{PP} = -1,50\,\text{D}$ (négatif car objet réel).
    $PP = \dfrac{1}{-1,50} \approx -0,667\,\text{m} = -66,7\,\text{cm}$.
    Il ne voit net de près qu'au-delà de $67\,\text{cm}$.
    \item \textbf{Addition pour lecture à 25 cm sans effort :} Objet $V_{\text{obj}} = -4,00\,\text{D}$. Image voulue au $PP$ naturel ($V_{PP} = -1,50\,\text{D}$).
    $Add = V_{PP} - V_{\text{obj}} = -1,50 - (-4,00) = +2,50\,\text{D}$.
    \textbf{Résultat :} $Add = +2,50\,\text{D}$ (lentille convergente).
    \item \textbf{Avec effort $0,50\,\text{D}$ :} L'œil peut fournir $0,50\,\text{D}$. L'image fournie par la loupe n'a pas besoin d'aller jusqu'au $PP$ max, elle peut s'arrêter à la limite de l'accommodation résiduelle.
    Vergence max acceptée par l'œil (au repos + effort) : $V_{\text{lim}} = V_{PR} - \Delta C = 0 - 0,50 = -0,50\,\text{D}$.
    $Add = V_{\text{lim}} - V_{\text{obj}} = -0,50 - (-4,00) = +3,50\,\text{D}$.
    \textit{Vérification :} $V' = Add + V_{\text{obj}} = 3,50 - 4,00 = -0,50\,\text{D} \Rightarrow p' = -2,00\,\text{m}$.
    L'œil voit cet objet virtuel à $2\,\text{m}$ en accommodant $0,50\,\text{D}$. C'est confortable.
    \textbf{Choix standard :} On vise souvent $Add = 2,50\,\text{D}$ (confort max) ou $3,00\,\text{D}$ (compromis). Ici $+3,50\,\text{D}$ est possible mais fort.
    \item \textbf{Myopie survenue ($PR = -2,00\,\text{m}$) :}
    \textbf{Correction loin :} $C_L = 1/PR = 1/(-2,00) = -0,50\,\text{D}$ (Divergente).
    \textbf{Nouvel œil corrigé de loin :} Devient "emmétrope artificiel". Son $PR'$ est à l'infini.
    Son $PP'$ corrigé de loin : Image de $PP$ par $C_L$.
    $V_{PP} = -1,50\,\text{D}$. $V_{PP'} = C_L + V_{PP} = -0,50 - 1,50 = -2,00\,\text{D}$.
    $PP' = -0,50\,\text{m} = -50\,\text{cm}$.
    \textbf{Addition pour lecture (sur correction de loin) :} On veut image au $PP'$ de l'œil corrigé ($V_{PP'} = -2,00\,\text{D}$) pour objet à $25\,\text{cm}$ ($-4,00\,\text{D}$).
    $Add = V_{PP'} - V_{\text{obj}} = -2,00 - (-4,00) = +2,00\,\text{D}$.
    \textbf{Verres progressifs :} Partie haute $-0,50\,\text{D}$, partie basse $-0,50 + 2,00 = +1,50\,\text{D}$.
\end{enumerate}
\end{exoresolu}

\begin{methode}[Méthode générale unifiée : "Vergence pas à pas" pour tout problème de correction]
\textbf{Objectif :} Résoudre systématiquement tout problème (myopie, hypermétropie, presbytie, astigmatisme non au prog) sans formules toutes faites.
\begin{enumerate}
    \item \textbf{Lister les données :} $n_{\text{air}}=1$, $n_{\text{œil}}=1,33$ (souvent implicite), $d_{\text{rétine}}$, $C_{\text{min}}$, $A_{\text{max}}$, $PR$, $PP$, distance lunettes-œil $e$ (souvent $0$).
    \item \textbf{Calculer les vergences caractéristiques de l'œil nu (dans l'air, au plan de la cornée/lentille) :}
    $V_{PR} = 1/PR$, $V_{PP} = 1/PP$. (Signes algébriques stricts : $<0$ réel devant, $>0$ virtuel derrière).
    \item \textbf{Identifier la demande :} "Voir net à l'infini" $\Rightarrow$ Objet $V_{\text{obj}}=0$. "Lire à 25 cm" $\Rightarrow$ Objet $V_{\text{obj}}=-4\,\text{D}$.
    \item \textbf{Écrire la chaîne de vergences :}
    \[
    V_{\text{obj}} \xrightarrow{\text{Lentille } C_L} V_{\text{inter}} \xrightarrow{\text{Œil (accommodation } \Delta C)} V_{\text{rétine}}
    \]
    Relation : $V_{\text{inter}} = C_L + V_{\text{obj}}$.
    Condition netteté : $V_{\text{inter}}$ doit être \textbf{comprise entre $V_{PP}$ et $V_{PR}$} (vergences objets que l'œil sait imager sur la rétine).
    \item \textbf{Résoudre pour $C_L$ (et $Add$) :}
    \begin{itemize}
        \item Vision de loin : On veut $V_{\text{inter}} = V_{PR}$ pour $V_{\text{obj}}=0$. $\Rightarrow C_L = V_{PR}$.
        \item Vision de près (avec correction de loin en place) : On veut $V_{\text{inter}} = V_{PP'}$ pour $V_{\text{obj}}=-4$. $\Rightarrow C_L + Add = V_{PP'} + 4$.
    \end{itemize}
    \item \textbf{Vérifier l'accommodation requise :} $\Delta C = V_{PR} - V_{\text{inter}}$. Doit être $\le A_{\text{max}}$ (en valeur absolue).
    \item \textbf{Conclure :} Nature de la lentille (Convergente $C>0$ / Divergente $C<0$), Puissance en Dioptries (arrondi au $0,25\,\text{D}$ standard), Position $PP'$ corrigé.
\end{enumerate}
\cle{Atout majeur :} Cette méthode évite d'apprendre des cas par cœur. Elle gère automatiquement les signes et les cas complexes (vertex $e \neq 0$, défauts mixtes).
\end{methode}

\begin{exoresolu}[Application méthode unifiée : Cas mixte myopie + presbytie]
\textbf{Énoncé :} Un sujet myope presbyte porte une correction de loin $C_L = -3,00\,\text{D}$ (lunettes au contact de l'œil). Son Punctum Proximum naturel (œil nu) est $PP = -12,0\,\text{cm}$. Son amplitude d'accommodation résiduelle est $A_{\text{max}} = 2,00\,\text{D}$. Il veut lire à $30\,\text{cm}$.
\begin{enumerate}
    \item Vérifier la cohérence des données de l'œil nu ($PR$, $PP$, $A_{\text{max}}$).
    \item Déterminer le Punctum Proximum $PP'$ de l'œil \textbf{corrigé de loin} (avec les lunettes $-3,00\,\text{D}$).
    \item Calculer l'addition $Add$ nécessaire pour lire à $30\,\text{cm}$ sans effort d'accommodation.
    \item Si on lui donne une addition $Add = +2,50\,\text{D}$, quel effort d'accommodation doit-il fournir pour lire à $30\,\text{cm}$ ? Est-ce possible ?
\end{enumerate}

\tcblower
\textbf{Solution détaillée (Méthode unifiée) :}
\begin{enumerate}
    \item \textbf{Cohérence œil nu :}
    $V_{PP} = 1/(-0,12) = -8,33\,\text{D}$.
    Myope : $C_L(\text{loin}) = -3,00\,\text{D} = V_{PR}$ (car $C_L = V_{PR}$ pour correction loin).
    Donc $V_{PR} = -3,00\,\text{D} \Rightarrow PR = -33,3\,\text{cm}$.
    $A_{\text{max}} = V_{PP} - V_{PR} = -8,33 - (-3,00) = -5,33\,\text{D} \Rightarrow |A| = 5,33\,\text{D}$.
    \textbf{Incohérence :} L'énoncé donne $A_{\text{max}} = 2,00\,\text{D}$.
    \textit{Hypothèse :} Le $PP = -12\,\text{cm}$ est celui de l'œil \emph{jeune} ou l'énoncé impose $A_{\text{max}}=2\,\text{D}$ pour l'âge actuel. On utilise les données "âge actuel" : $A_{\text{max}} = 2,00\,\text{D}$, $V_{PR} = -3,00\,\text{D}$.
    Alors $V_{PP}(\text{actuel}) = V_{PR} - A_{\text{max}} = -3,00 - 2,00 = -5,00\,\text{D}$.
    $PP_{\text{actuel}} = -20,0\,\text{cm}$.
    (On ignore le $PP=-12\,\text{cm}$ donné comme "naturel" s'il contredit $A_{\text{max}}$, ou on considère qu'il s'agit du $PP$ juvénile. \textbf{Au Probatoire, on utilise les données explicites du sujet âgé : $A_{\text{max}}$ et $PR$ déduit de $C_L$}).
    \item \textbf{PP' corrigé de loin :} Œil nu (actuel) : $V_{PR} = -3,00\,\text{D}$, $V_{PP} = -5,00\,\text{D}$.
    Lentille loin $C_L = -3,00\,\text{D}$.
    Correction de loin : La lentille envoie l'infini ($0$) sur $PR$ ($-3\,\text{D}$). $C_L = -3\,\text{D}$.
    Pour l'œil corrigé, le nouveau $PR'$ est l'infini ($V_{PR'}=0$).
    Le nouveau $PP'$ est l'image de $PP$ par la lentille.
    $V_{PP'} = C_L + V_{PP} = -3,00 + (-5,00) = -8,00\,\text{D}$.
    $PP' = 1/(-8,00) = -0,125\,\text{m} = -12,5\,\text{cm}$.
    \item \textbf{Addition pour lecture à 30 cm sans effort :}
    Objet : $30\,\text{cm} \Rightarrow V_{\text{obj}} = -1/0,30 = -3,33\,\text{D}$.
    On veut image au $PP'$ de l'œil corrigé ($V_{PP'} = -8,00\,\text{D}$) pour que l'œil soit au repos ($\Delta C = 0$).
    $C_{\text{totale}} = V_{PP'} - V_{\text{obj}} = -8,00 - (-3,33) = -4,67\,\text{D}$.
    $Add = C_{\text{totale}} - C_L = -4,67 - (-3,00) = -1,67\,\text{D}$.
    \textbf{Alerte :} Addition \textbf{négative} ? Impossible pour une loupe.
    \textit{Analyse :} $PP' = -12,5\,\text{cm}$. L'objet désiré est à $30\,\text{cm}$. L'objet est \textbf{plus loin} que le $PP'$.
    L'œil corrigé de loin voit net de $12,5\,\text{cm}$ à l'infini. $30\,\text{cm}$ est \textbf{dans cet intervalle}.
    \textbf{Conclusion :} \textbf{Aucune addition n'est nécessaire} ($Add = 0$). Il lit à $30\,\text{cm}$ avec sa seule correction de loin $-3,00\,\text{D}$.
    \item \textbf{Si $Add = +2,50\,\text{D}$ imposée (erreur de prescription) :}
    Puissance totale près : $C_{\text{près}} = C_L + Add = -3,00 + 2,50 = -0,50\,\text{D}$.
    Objet $30\,\text{cm}$ ($V=-3,33\,\text{D}$).
    Vergence entrant dans l'œil : $V_{\text{inter}} = C_{\text{près}} + V_{\text{obj}} = -0,50 - 3,33 = -3,83\,\text{D}$.
    L'œil corrigé a $V_{PR'}=0$ et $V_{PP'}=-8,00\,\text{D}$.
    $-3,83\,\text{D}$ est dans $[ -8,00 ; 0 ]$. Image nette possible.
    Accommodation requise : $\Delta C = V_{PR'} - V_{\text{inter}} = 0 - (-3,83) = 3,83\,\text{D}$.
    $|\Delta C| = 3,83\,\text{D} > A_{\text{max}} = 2,00\,\text{D}$.
    \textbf{Conclusion :} Effort $3,83\,\text{D}$ impossible (dépasse $A_{\text{max}}$). Vision floue. Mauvaise correction.
\end{enumerate}
\end{exoresolu}

\begin{methode}[Vérifier la cohérence d'une correction : Bilan "Vision de loin / Vision de près"]
\textbf{Objectif :} Valider une prescription complète (loin + addition) avant l'examen.
\begin{enumerate}
    \item \textbf{Étape 1 : Vision de loin corrigée.}
    Calculer $V_{PR'}$ (doit être $0$ pour emmétropie corrigée). Si $V_{PR'} \neq 0$, il reste un résidu (sous/sur-correction).
    \item \textbf{Étape 2 : Vision de près corrigée (avec addition).}
    Calculer $V_{PP'}$ (limite proche) et $V_{PR'}$ (limite loin) de l'œil \textbf{totalement corrigé} (loin + addition).
    $V_{PP'} = C_L + Add + V_{PP}$ (vergences air).
    $V_{PR'} = C_L + Add + V_{PR}$.
    \item \textbf{Étape 3 : Tester les distances usuelles.}
    \begin{itemize}
        \item Loin : $V_{\text{obj}}=0 \to V_{\text{inter}} = C_L$. Doit être $\in [V_{PP}, V_{PR}]$ de l'œil nu. (Vérifié par construction de $C_L$).
        \item Près (25 ou 30 ou 40 cm) : $V_{\text{obj}} = -1/d$. $V_{\text{inter}} = C_L + Add + V_{\text{obj}}$.
        Doit être $\in [V_{PP}, V_{PR}]$ de l'œil nu.
    \end{itemize}
    \item \textbf{Étape 4 : Calculer l'accommodation demandée.}
    $\Delta C_{\text{loin}} = C_L - V_{PR}$ (doit être $\approx 0$).
    $\Delta C_{\text{près}} = (C_L + Add + V_{\text{obj}}) - V_{PR}$.
    Vérifier $\Delta C_{\text{près}} \le A_{\text{max}}$.
    \item \textbf{Étape 5 : Champ de vision net corrigé.}
    Intervalle $[PP'_{\text{corr}} ; PR'_{\text{corr}}]$.
    $PP'_{\text{corr}} = 1 / (C_L + Add + V_{PP})$.
    $PR'_{\text{corr}} = 1 / (C_L + Add + V_{PR})$ (souvent $\infty$ si $Add$ calculée pour $PR'=\infty$).
\end{enumerate}
\cle{Rédiger au Probatoire :} "L'œil corrigé voit net de $X\,\text{cm}$ à $Y\,\text{cm}$ (ou l'infini). L'accommodation maximale demandée est $Z\,\text{D}$, compatible avec $A_{\text{max}}$."
\end{methode}

\begin{exoresolu}[Bilan complet d'une correction progressive]
\textbf{Énoncé :} Œil hypermétrope : $PR = +40\,\text{cm}$, $PP = -60\,\text{cm}$, $A_{\text{max}} = 3,00\,\text{D}$ (cohérence : $1/0,4 - 1/(-0,6) = 2,5 + 1,67 = 4,17 \neq 3$... On impose $A_{\text{max}}=3\,\text{D}$ et $PR=+40\,\text{cm}$ $\Rightarrow V_{PR}=+2,5\,\text{D}$, $V_{PP}=V_{PR}-A = +2,5 - 3 = -0,5\,\text{D} \Rightarrow PP = -2,00\,\text{m}$).
Correction loin : $C_L = +2,50\,\text{D}$. Addition pour lecture à $25\,\text{cm}$ : $Add = +2,00\,\text{D}$.
\begin{enumerate}
    \item Vérifier que la correction de loin ramène bien le $PR$ à l'infini.
    \item Déterminer l'intervalle de vision nette avec la correction de loin seule.
    \item Déterminer l'intervalle de vision nette avec la correction de près (loin + addition).
    \item Vérifier si la lecture à $25\,\text{cm}$ est possible sans effort excessif.
    \item Commenter la vision intermédiaire (écran d'ordinateur à $60\,\text{cm}$) avec les verres progressifs (partie haute $+2,50\,\text{D}$, partie basse $+4,50\,\text{D}$).
\end{enumerate}

\tcblower
\textbf{Solution détaillée :}
\textbf{Données cohérentes imposées :} $V_{PR} = +2,50\,\text{D}$, $A_{\text{max}} = 3,00\,\text{D} \Rightarrow V_{PP} = V_{PR} - A_{\text{max}} = +2,50 - 3,00 = -0,50\,\text{D}$ ($PP = -2,00\,\text{m}$).
\begin{enumerate}
    \item \textbf{Correction loin :} $C_L = V_{PR} = +2,50\,\text{D}$.
    Vergence entrant dans l'œil pour objet à l'infini ($V=0$) : $V_{\text{inter}} = C_L + 0 = +2,50\,\text{D} = V_{PR}$.
    L'œil reçoit sa vergence de repos. Image nette sur rétine sans accommodation.
    \textbf{Validé :} $PR' = \infty$.
    \item \textbf{Intervalle vision loin seule ($C_L = +2,50\,\text{D}$) :}
    Vergences entrant dans l'œil : $V_{\text{inter}} = 2,50 + V_{\text{obj}}$.
    Condition : $V_{PP} \le V_{\text{inter}} \le V_{PR}$ $\Rightarrow$ $-0,50 \le 2,50 + V_{\text{obj}} \le 2,50$.
    $-3,00 \le V_{\text{obj}} \le 0$.
    Objets réels : $V_{\text{obj}} \le 0$. Intervalle objets : $[ -3,00\,\text{D} ; 0 ]$.
    Distances : $PP' = 1/(-3,00) = -33,3\,\text{cm}$ à $PR' = \infty$.
    \textbf{Vision nette de loin seule :} De $33\,\text{cm}$ à l'infini.
    \item \textbf{Intervalle vision près (Progressif bas : $C_L + Add = +4,50\,\text{D}$) :}
    $V_{\text{inter}} = 4,50 + V_{\text{obj}}$.
    Condition : $-0,50 \le 4,50 + V_{\text{obj}} \le 2,50$.
    $-5,00 \le V_{\text{obj}} \le -2,00$.
    Distances : $PP'' = 1/(-5,00) = -20,0\,\text{cm}$ à $PR'' = 1/(-2,00) = -50,0\,\text{cm}$.
    \textbf{Vision nette partie basse :} De $20\,\text{cm}$ à $50\,\text{cm}$.
    \item \textbf{Lecture à 25 cm :} $V_{\text{obj}} = -4,00\,\text{D}$.
    Avec partie basse ($+4,50\,\text{D}$) : $V_{\text{inter}} = 4,50 - 4,00 = +0,50\,\text{D}$.
    $V_{PR} = +2,50\,\text{D}$. $V_{PP} = -0,50\,\text{D}$.
    $+0,50\,\text{D} \in [-0,50 ; +2,50]$. \textbf{Net.}
    Accommodation : $\Delta C = V_{PR} - V_{\text{inter}} = 2,50 - 0,50 = 2,00\,\text{D}$.
    $|\Delta C| = 2,00\,\text{D} \le A_{\text{max}} = 3,00\,\text{D}$. \textbf{Effort modéré, possible.}
    \item \textbf{Écran à 60 cm ($V = -1,67\,\text{D}$) :}
    \textit{Partie haute (loin $+2,50\,\text{D}$) :} $V_{\text{inter}} = 2,50 - 1,67 = +0,83\,\text{D} \in [-0,50; 2,50]$. Net. $\Delta C = 2,50 - 0,83 = 1,67\,\text{D}$. OK.
    \textit{Partie basse (près $+4,50\,\text{D}$) :} $V_{\text{inter}} = 4,50 - 1,67 = +2,83\,\text{D} > V_{PR}(+2,50)$.
    \textbf{Flou !} (Image formée derrière la rétine, œil pas assez puissant même au repos).
    \textbf{Conclusion :} Pour l'écran à $60\,\text{cm}$, il doit utiliser la \textbf{partie haute (vision de loin)} ou une zone intermédiaire du progressif (puissance $\approx +3,50\,\text{D}$ pour centrer l'intervalle sur $60\,\text{cm}$).
\end{enumerate}
\end{exoresolu}

\begin{attention}[Les 5 erreurs qui coûtent des points au Probatoire]
\begin{enumerate}
    \item \textbf{Signes algébriques inversés sur $PP$ et $PR$ :} Oublier que $PP$ et $PR$ sont des \textbf{points objets} (côté entrée lumière).
    \begin{itemize}
        \item Myope : $PR$ réel $\Rightarrow$ \textbf{Négatif} (ex $-2\,\text{m}$). $PP$ réel $\Rightarrow$ \textbf{Négatif}.
        \item Hypermétrope : $PR$ virtuel $\Rightarrow$ \textbf{Positif} (ex $+50\,\text{cm}$). $PP$ réel (souvent) $\Rightarrow$ \textbf{Négatif}.
        \item \textit{Conséquence :} Puissance lentille $C_L = 1/PR$ fausse (signe opposé $\Rightarrow$ nature de lentille inversée : convergente au lieu de divergente). \textbf{0 point sur tout l'exercice.}
    \end{itemize}
    \item \textbf{Confondre "Vergence air" et "Vergence dans l'œil" :} Utiliser $C = 1/f'$ (air) pour l'œil au lieu de $C = n'/f'$ (milieu vitré $n'=1,33$).
    \textit{Règle :} Pour les calculs de \textbf{correction (lunettes dans l'air)}, on travaille \textbf{uniquement en vergences air} ($V = 1/p$ en m$^{-1}$). La formule $C = V' - V$ s'applique avec $V, V'$ en air. On n'utilise $n'=1,33$ que pour calculer la puissance interne de l'œil ($C_{\text{œil}} = n'/p'$) si $p'$ (distance rétine) est donnée.
    \item \textbf{Oublier la distance sommet-cornée / vertex ($e$) :} Si l'énoncé donne $e = 12\,\text{mm}$ (ou $1,5\,\text{cm}$), la formule $C_L = 1/PR$ est fausse.
    \textit{Correction :} $V_{\text{sortie lentille}} = \dfrac{V_{PR}}{1 - e \cdot V_{PR}}$ (vergences air). Puis $C_L = V_{\text{sortie}} - V_{\text{entrée}}$. Ne pas l'appliquer coûte 1 à 2 points.
    \item \textbf{Calculer l'addition $Add$ par rapport au $PP$ de l'œil nu au lieu du $PP'$ corrigé de loin :} Pour un myope ou hypermétrope presbyte, l'addition s'ajoute à la correction de loin.
    \textit{Erreur :} $Add = 1/PP_{\text{nu}} + 4$.
    \textit{Juste :} $Add = 1/PP'_{\text{corr\_loin}} + 4$ (si lecture 25 cm sans effort). $PP'_{\text{corr\_loin}}$ est l'image de $PP_{\text{nu}}$ par $C_L$.
    \item \textbf{Rédaction insuffisante / Absence de conclusion physique :} Donner un nombre ($-2,50\,\text{D}$) sans écrire "Lentille divergente de puissance $-2,50\,\text{D}$". Ne pas vérifier si l'accommodation demandée dépasse $A_{\text{max}}$ (réponse "Oui/Non" non justifiée). Ne pas préciser "Vision nette de $X\,\text{cm}$ à l'infini" en conclusion.
\end{enumerate}
\end{attention}

\newpage

\coursec{Exercices}

\courssub{Je vérifie mes connaissances}

\begin{exercice}[QCM : l'œil réduit]
Pour chaque affirmation, choisir la (ou les) bonne(s) réponse(s) en justifiant brièvement.
\begin{enumerate}
\item Dans le modèle de l'œil réduit, le cristallin est assimilé à :
\begin{enumerate}
\item une lentille mince divergente ;
\item une lentille mince convergente de vergence variable ;
\item un miroir convergent ;
\item un diaphragme.
\end{enumerate}
\item Lorsque l'œil accommode au maximum :
\begin{enumerate}
\item la vergence du cristallin diminue ;
\item la vergence du cristallin augmente ;
\item la distance cristallin-rétine augmente ;
\item la rétine se rapproche du cristallin.
\end{enumerate}
\item Un œil emmétrope au repos voit net un objet situé :
\begin{enumerate}
\item à son punctum proximum ;
\item à 25 cm ;
\item à l'infini ;
\item sur sa rétine.
\end{enumerate}
\item La myopie se corrige avec :
\begin{enumerate}
\item une lentille convergente ;
\item une lentille divergente ;
\item un prisme ;
\item un miroir plan.
\end{enumerate}
\end{enumerate}
\end{exercice}

\begin{exercice}[Vrai ou faux]
Répondre par vrai ou faux en justifiant chaque réponse par une phrase ou un calcul.
\begin{enumerate}
\item L'image d'un objet réel donnée par l'œil se forme toujours sur la rétine pour être perçue nettement.
\item Le punctum proximum est le point le plus éloigné que l'œil peut voir nettement.
\item Un œil hypermétrope a une vergence trop faible : l'image d'un objet proche se formerait derrière la rétine.
\item La presbytie est due à un allongement du globe oculaire avec l'âge.
\item Pour un œil emmétrope de distance cristallin-rétine $d = \SI{15}{mm}$, la vergence au repos vaut environ $66,7\ \delta$.
\end{enumerate}
\end{exercice}

\begin{exercice}[Définitions et schéma]
\begin{enumerate}
\item Définir : œil réduit, punctum remotum (PR), punctum proximum (PP), accommodation, amplitude d'accommodation.
\item Schématiser l'œil réduit emmétrope au repos en faisant figurer : la lentille mince équivalente, son centre optique $O$, ses foyers $F$ et $F'$, la rétine, et la marche de deux rayons issus d'un point objet $B$ à l'infini.
\item Sur un second schéma, représenter le même œil accommodant au maximum pour observer un objet placé au PP.
\end{enumerate}
\end{exercice}

\begin{exercice}[Calculs directs]
On modélise un œil par une lentille mince convergente de centre $O$ ; la rétine est située à $d = \SI{15}{mm}$ du centre optique.
\begin{enumerate}
\item Calculer la vergence $C_{\min}$ de cet œil au repos, sachant qu'il voit alors net un objet à l'infini.
\item En accommodation maximale, sa vergence vaut $C_{\max} = 70\ \delta$. Déterminer la position algébrique $\overline{OA}$ du punctum proximum.
\item En déduire son amplitude d'accommodation $A = C_{\max} - C_{\min}$.
\end{enumerate}
\end{exercice}

\courssub{Je m'entraîne}

\begin{exercice}[Correction d'un œil myope]
Un œil myope a son punctum remotum à $80$ cm devant lui et son punctum proximum à $15$ cm.
\begin{enumerate}
\item Justifier qualitativement pourquoi cet œil ne peut pas voir net un objet à l'infini sans correction.
\item On accole à l'œil une lentille correctrice $L$ de vergence $C_c$. Déterminer la nature (convergente ou divergente) et la vergence de $L$ pour que l'œil voie net un objet à l'infini sans accommoder. On rappelle que la lentille doit donner de l'objet à l'infini une image virtuelle située au PR de l'œil.
\item Calculer la nouvelle position du punctum proximum de l'œil corrigé (objet réel donnant, à travers $L$, une image virtuelle à $15$ cm devant l'œil).
\item Représenter la marche de deux rayons issus d'un point à l'infini, avant et après correction.
\end{enumerate}
\end{exercice}

\begin{exercice}[Correction d'un œil hypermétrope]
Un œil hypermétrope a son punctum proximum situé à $75$ cm. Il souhaite lire un journal placé à $25$ cm de son œil.
\begin{enumerate}
\item Expliquer pourquoi la lecture lui est impossible sans correction, même en accommodant au maximum.
\item Déterminer la nature et la vergence $C_c$ de la lentille de contact (accolée à l'œil) qui ramène son PP apparent à $25$ cm. On écrira la relation de conjugaison en précisant les signes des grandeurs algébriques.
\item Construire la marche des rayons : objet réel à $25$ cm, lentille correctrice, image virtuelle servant d'objet pour l'œil.
\end{enumerate}
\end{exercice}

\begin{exercice}[Œil emmétrope : vergences extrêmes]
Un élève emmétrope a son PR à l'infini et son PP à $25$ cm. La distance cristallin-rétine vaut $d = \SI{15}{mm}$.
\begin{enumerate}
\item Calculer sa vergence minimale $C_{\min}$ (œil au repos).
\item Calculer sa vergence maximale $C_{\max}$ (accommodation maximale, image nette sur la rétine d'un objet placé au PP).
\item En déduire son amplitude d'accommodation.
\item Comparer la valeur de $C_{\min}$ à la vergence moyenne d'un œil emmétrope ($60\ \delta$). Commenter.
\end{enumerate}
\end{exercice}

\begin{exercice}[Presbytie et verres de lecture]
Une personne de 55 ans a son punctum proximum reculé à $1{,}00$ m ; son punctum remotum est resté à l'infini.
\begin{enumerate}
\item Quel est le défaut de cet œil ? Quelle en est la cause physiologique ?
\item Calculer la vergence $C_c$ des verres de lecture (accolés à l'œil) lui permettant de lire un texte placé à $25$ cm.
\item Expliquer pourquoi cette personne retire ses verres de lecture pour regarder un paysage au loin. On pourra calculer la position de l'image, à travers les verres, d'un objet à l'infini.
\end{enumerate}
\end{exercice}

\courssub{Je me prépare au Probatoire}

\begin{exercice}[Consultation chez l'ophtalmologue de Yaoundé]
Au cours d'une visite médicale scolaire, un élève de Première C est examiné. On modélise son œil par une lentille mince convergente de centre optique $O$, la rétine étant fixée à $d = \SI{15}{mm}$ de $O$. Les mesures donnent : vergence au repos $C_{\min} = 60\ \delta$ ; vergence en accommodation maximale $C_{\max} = 64\ \delta$.
\begin{enumerate}
\item Rappeler la définition de l'œil réduit et citer les trois éléments essentiels de ce modèle.
\item Déterminer la position du punctum remotum de cet élève. Justifier le calcul par la relation de conjugaison des lentilles minces appliquée à l'œil au repos.
\item Déterminer la position de son punctum proximum.
\item Calculer son amplitude d'accommodation $A$.
\item L'élève est-il emmétrope, myope ou hypermétrope ? Justifier la réponse à partir des positions du PR et du PP.
\item L'ophtalmologue annonce que, sans correction, cet œil donne d'un objet éloigné une image qui se formerait en avant de la rétine. Cette affirmation est-elle compatible avec les résultats précédents ? Conclure.
\end{enumerate}
On donne : $1/60 = \num{0,01667}$ ; $1/64 = \num{0,01563}$.
\end{exercice}

\begin{exercice}[Le chauffeur de mototaxi de Douala]
M. Etoundi, chauffeur de mototaxi à Douala, consulte car il ne distingue plus les panneaux routiers au loin. L'examen révèle un punctum remotum à $\overline{OR} = -\SI{2,0}{m}$ (l'origine $O$ étant au centre optique de l'œil, l'axe étant orienté dans le sens de la lumière) et un punctum proximum à $20$ cm.
\begin{enumerate}
\item Nommer le défaut de vision de M. Etoundi et expliquer son origine à l'aide d'un schéma de l'œil réduit non corrigé observant un objet à l'infini.
\item Déterminer la nature et la vergence $C_1$ de la lentille correctrice, accolée à l'œil, qui lui permet de voir net un objet à l'infini sans accommoder.
\item Quelques années plus tard, M. Etoundi devient presbyte : son PP passe de $20$ cm à $40$ cm, son PR corrigé restant à l'infini avec ses verres.
\begin{enumerate}
\item Calculer la vergence $C_2$ de la correction supplémentaire nécessaire pour lire un document placé à $25$ cm (on raisonnera sur l'œil déjà corrigé pour la vision de loin).
\item En déduire la vergence totale des verres de lecture.
\end{enumerate}
\item Expliquer l'intérêt des verres progressifs (double foyer) pour ce patient, en précisant le rôle de la partie supérieure et de la partie inférieure du verre.
\item Pour la sécurité routière, citer deux conséquences pratiques d'une myopie non corrigée chez un chauffeur.
\end{enumerate}
\end{exercice}

\begin{exercice}[Étude complète d'un œil hypermétrope]
Lors d'une campagne de dépistage visuel organisée dans un lycée de Bafoussam, on étudie l'œil d'une élève. Données : distance cristallin-rétine $d = \SI{15}{mm}$ ; punctum proximum à $75$ cm ; punctum remotum virtuel situé derrière l'œil à $\overline{OR} = +\SI{50}{cm}$ (l'œil doit accommoder pour voir net un objet à l'infini).
\begin{enumerate}
\item Faire un schéma de l'œil réduit de cette élève observant un objet rapproché, en montrant que l'image se formerait derrière la rétine. Nommer le défaut.
\item Calculer la vergence $C_{\max}$ de cet œil en accommodation maximale (image nette sur la rétine d'un objet au PP).
\item Calculer la vergence $C_0$ que devrait avoir cet œil au repos pour voir net un objet à l'infini. Comparer avec la vergence réelle au repos et commenter.
\item L'élève veut lire à $25$ cm sans effort excessif. Déterminer la nature et la vergence $C_c$ de la lentille de contact correctrice.
\item Vérifier par analyse dimensionnelle l'homogénéité de la relation de conjugaison utilisée à la question précédente.
\item Rédiger une courte conclusion (5 lignes maximum) destinée au carnet de santé scolaire, précisant le diagnostic et la correction prescrite.
\end{enumerate}
\end{exercice}

\newpage

\coursec{Corrigés des exercices}

\coursec{Exercices corrigés — L'œil réduit : fonctionnement, défauts et correction}

\begin{corrige}[Exercice 1 — QCM : l'œil réduit]
\begin{enumerate}
\item \textbf{Réponse b.} Le cristallin est assimilé à une \cle{lentille mince convergente} dont la vergence varie grâce à l'action des muscles ciliaires (accommodation). Il n'est ni divergent, ni un miroir ; le diaphragme correspond à l'iris, pas au cristallin.
\item \textbf{Réponse b.} Accommoder, c'est bomber le cristallin : sa distance focale image $f'$ diminue, donc sa vergence $C = 1/f'$ \textbf{augmente}. La distance cristallin-rétine $d$ est fixe (le globe oculaire ne se déforme pas), ce qui élimine les réponses c et d.
\item \textbf{Réponse c.} Un œil emmétrope au repos (sans accommodation) voit net un objet situé à son punctum remotum, qui est \textbf{à l'infini}. Le PP (25 cm pour l'œil normal) exige une accommodation maximale : l'œil n'est alors plus au repos.
\item \textbf{Réponse b.} La myopie (œil trop convergent) se corrige par une lentille \textbf{divergente} ($C_c < 0$) qui écarte les rayons et ramène l'image sur la rétine.
\end{enumerate}
\end{corrige}

\begin{corrige}[Exercice 2 — Vrai ou faux]
\begin{enumerate}
\item \textbf{Vrai.} La rétine joue le rôle d'écran : la vision n'est nette que si l'image $A'B'$ se forme exactement sur la rétine. Si elle se forme en avant ou en arrière, la perception est floue.
\item \textbf{Faux.} Le punctum proximum est le point le plus \emph{proche} vu nettement (accommodation maximale). Le point le plus éloigné est le punctum remotum.
\item \textbf{Vrai.} L'œil hypermétrope est trop peu convergent (vergence insuffisante ou globe trop court) : pour un objet proche, les rayons n'ont pas fini de converger à la rétine ; l'image se formerait derrière elle.
\item \textbf{Faux.} La presbytie est due au durcissement (perte d'élasticité) du cristallin et au relâchement des muscles ciliaires avec l'âge : l'amplitude d'accommodation diminue et le PP recule. Le globe oculaire ne s'allonge pas.
\item \textbf{Vrai.} Au repos, l'image d'un objet à l'infini se forme sur la rétine : $f' = d = \SI{15}{mm} = \SI{0,015}{m}$, d'où
\[ C_{\min} = \frac{1}{f'} = \frac{1}{\num{0,015}} \approx 66,7\ \delta. \]
\end{enumerate}
\end{corrige}

\begin{corrige}[Exercice 3 — Définitions et schéma]
\begin{enumerate}
\item \textbf{Définitions.}
\begin{itemize}
\item \cle{Œil réduit} : modèle simplifié de l'œil constitué d'une lentille mince convergente (le cristallin) de vergence variable, d'un diaphragme (l'iris) et d'un écran fixe (la rétine) placé à la distance $d \approx \SI{15}{mm}$ de la lentille.
\item \cle{Punctum remotum} (PR) : point de l'axe optique vu nettement \emph{sans accommodation} (œil au repos).
\item \cle{Punctum proximum} (PP) : point de l'axe optique vu nettement avec une \emph{accommodation maximale} ; c'est le point le plus proche visible nettement.
\item \cle{Accommodation} : modification de la courbure du cristallin (donc de sa vergence) par les muscles ciliaires, permettant de maintenir l'image nette sur la rétine quand l'objet se rapproche.
\item \cle{Amplitude d'accommodation} : $A = C_{\max} - C_{\min}$, écart entre la vergence maximale et la vergence au repos (en dioptries).
\end{itemize}
\item \textbf{Œil au repos, objet $B$ à l'infini.} Les rayons issus de $B$ arrivent parallèles entre eux ; le rayon passant par $O$ n'est pas dévié, le rayon parallèle à l'axe émerge en passant par $F'$. L'image se forme dans le plan focal image, qui coïncide avec la rétine.
\begin{popfigure}
\begin{center}
\begin{tikzpicture}[scale=0.9,>=Stealth]
\draw[->] (-5.5,0) -- (3.5,0) node[below left]{\small axe optique};
\draw[popblue,thick] (0,-1.8) -- (0,1.8);
\draw[popblue,->] (0,1.8) -- (-0.25,1.55);
\draw[popblue,->] (0,-1.8) -- (-0.25,-1.55);
\fill (0,0) circle (1.5pt) node[below right]{\small $O$};
\fill[popink] (-1.5,0) circle (1.5pt) node[below]{\small $F$};
\fill[popink] (1.5,0) circle (1.5pt) node[below]{\small $F'$};
\draw[poporange,very thick] (1.5,-1.8) -- (1.5,1.8) node[above]{\small rétine};
\draw[popgreen,thick,->] (-5,1.2) -- (0,1.2);
\draw[popgreen,thick,->] (0,1.2) -- (1.5,0);
\draw[popgreen,thick,->] (-5,0.6) -- (0,0.6);
\draw[popgreen,thick,->] (0,0.6) -- (1.5,0);
\node[popgreen] at (-4.2,1.6) {\small rayons parallèles (objet à l'infini)};
\end{tikzpicture}
\end{center}
\legende{Œil réduit emmétrope au repos : $F'$ est sur la rétine.}
\end{popfigure}
\item \textbf{Accommodation maximale (objet au PP).} Le cristallin se bombe : $f'$ diminue, $F'$ se rapproche de $O$ (en avant de la rétine). Les rayons issus du point proche $A$ (au PP) convergent alors exactement sur la rétine : le rayon passant par $O$ n'est pas dévié, le rayon parallèle à l'axe émerge par $F'$ ; leur intersection est sur la rétine. (Schéma identique au précédent, avec $A$ placé à $25$ cm environ devant $O$, $F'$ entre $O$ et la rétine, et les rayons émergents plus inclinés.)
\end{enumerate}
\end{corrige}

\begin{corrige}[Exercice 4 — Calculs directs]
\begin{enumerate}
\item Au repos, l'objet est à l'infini ($\overline{OA} \to -\infty$, donc $1/\overline{OA} \to 0$) et l'image est sur la rétine : $\overline{OA'} = d = \SI{0,015}{m}$. La relation de conjugaison donne :
\[ C_{\min} = \frac{1}{\overline{OA'}} - \frac{1}{\overline{OA}} = \frac{1}{\num{0,015}} \approx 66,7\ \delta. \]
\item En accommodation maximale, l'image reste sur la rétine ($\overline{OA'} = \SI{0,015}{m}$) :
\[ C_{\max} = \frac{1}{\overline{OA'}} - \frac{1}{\overline{OA}} \implies \frac{1}{\overline{OA}} = \frac{1}{\num{0,015}} - 70 = 66,67 - 70 = -3,33\ \text{m}^{-1}. \]
\[ \overline{OA} = \frac{1}{-3,33} \approx -\SI{0,30}{m} = -\SI{30}{cm}. \]
Le signe négatif confirme que le PP est un objet réel situé à $30$ cm devant l'œil.
\item Amplitude d'accommodation :
\[ A = C_{\max} - C_{\min} = 70 - 66,7 = 3,3\ \delta. \]
\end{enumerate}
\end{corrige}

\begin{corrige}[Exercice 5 — Correction d'un œil myope]
\begin{enumerate}
\item L'œil myope est trop convergent : son PR est à $80$ cm seulement. Un objet à l'infini envoie des rayons parallèles que l'œil converge trop fortement : l'image se forme \emph{en avant} de la rétine, la vision est floue. Aucune accommodation ne peut corriger cela (accommoder ne fait qu'augmenter encore la vergence).
\item La lentille $L$ doit donner d'un objet à l'infini une image virtuelle au PR : $\overline{OA'} = -\SI{0,80}{m}$. Relation de conjugaison appliquée à $L$ :
\[ C_c = \frac{1}{\overline{OA'}} - \frac{1}{\overline{OA}} = \frac{1}{-\num{0,80}} - 0 = -1,25\ \delta. \]
$C_c < 0$ : la lentille est \textbf{divergente}, de distance focale $f' = 1/C_c = -\SI{0,80}{m}$.
\item Pour le nouveau PP, l'objet réel (position $\overline{OA}$ inconnue) doit donner à travers $L$ une image virtuelle à l'ancien PP : $\overline{OA'} = -\SI{0,15}{m}$.
\[ \frac{1}{\overline{OA}} = \frac{1}{\overline{OA'}} - C_c = \frac{1}{-\num{0,15}} - (-1,25) = -6,67 + 1,25 = -5,42\ \text{m}^{-1}. \]
\[ \overline{OA} \approx -\SI{0,185}{m} \approx -\SI{18}{cm}. \]
L'œil corrigé peut donc voir net de $18$ cm à l'infini : le champ de vision est élargi.
\item \textbf{Marche des rayons.} Avant correction : les rayons parallèles issus du point à l'infini convergent en un point situé \emph{devant} la rétine, puis divergent jusqu'à elle (tache floue). Après correction : la lentille divergente rend les rayons légèrement divergents, comme s'ils provenaient du PR ($80$ cm) ; l'œil au repos les converge alors exactement sur la rétine (le rayon visant $O$ n'est pas dévié, le rayon parallèle à l'axe émerge de $L$ comme venant de son foyer image $F'_L$ confondu avec le PR).
\end{enumerate}
\end{corrige}

\begin{corrige}[Exercice 6 — Correction d'un œil hypermétrope]
\begin{enumerate}
\item L'œil hypermétrope est trop peu convergent : même en accommodation maximale, il ne voit net qu'à partir de $75$ cm (son PP). Un journal à $25$ cm enverrait des rayons trop divergents : l'image se formerait \emph{derrière} la rétine ; l'accommodation étant déjà maximale, la lecture est impossible.
\item La lentille accolée doit donner de l'objet réel à $25$ cm une image \emph{virtuelle} au PP ($75$ cm devant l'œil), qui servira d'objet pour l'œil. Grandeurs algébriques (axe orienté dans le sens de la lumière) :
\[ \overline{OA} = -\SI{0,25}{m} \quad (\text{objet réel}), \qquad \overline{OA'} = -\SI{0,75}{m} \quad (\text{image virtuelle}). \]
\[ C_c = \frac{1}{\overline{OA'}} - \frac{1}{\overline{OA}} = \frac{1}{-\num{0,75}} - \frac{1}{-\num{0,25}} = -1,33 + 4,00 = +2,67\ \delta. \]
$C_c > 0$ : lentille \textbf{convergente}, $f' \approx \SI{37,5}{cm}$.
\item \textbf{Construction.} De l'objet $B$ à $25$ cm : le rayon passant par le centre $O_L$ de la lentille n'est pas dévié ; le rayon parallèle à l'axe émerge en passant par $F'_L$. Ces deux rayons émergents divergent ; leurs prolongements (tracés en pointillés) se coupent en $B'$, image virtuelle située à $75$ cm devant l'œil. L'œil, accommodant au maximum, voit $B'$ nettement : le journal est lu sans effort excessif.
\end{enumerate}
\end{corrige}

\begin{corrige}[Exercice 7 — Œil emmétrope : vergences extrêmes]
\begin{enumerate}
\item Au repos, objet à l'infini, image sur la rétine ($\overline{OA'} = d = \SI{0,015}{m}$) :
\[ C_{\min} = \frac{1}{d} = \frac{1}{\num{0,015}} \approx 66,7\ \delta. \]
\item Accommodation maximale : objet au PP, $\overline{OA} = -\SI{0,25}{m}$, image toujours sur la rétine :
\[ C_{\max} = \frac{1}{\num{0,015}} - \frac{1}{-\num{0,25}} = 66,7 + 4,0 = 70,7\ \delta. \]
\item Amplitude d'accommodation :
\[ A = C_{\max} - C_{\min} = 70,7 - 66,7 = 4,0\ \delta. \]
On retrouve la valeur classique : pour un œil emmétrope, $A = 1/d_{\mathrm{PP}} = 1/\num{0,25} = 4\ \delta$.
\item La valeur $C_{\min} \approx 66,7\ \delta$ est supérieure à la vergence moyenne réelle ($\approx 60\ \delta$). L'écart vient du modèle : on a pris $d = \SI{15}{mm}$, alors que la distance cristallin-rétine réelle est plutôt de $17$ mm, et que la réfraction se répartit entre la cornée et le cristallin. Le modèle de l'œil réduit reste néanmoins excellent pour les calculs d'accommodation et de correction.
\end{enumerate}
\end{corrige}

\begin{corrige}[Exercice 8 — Presbytie et verres de lecture]
\begin{enumerate}
\item Il s'agit de la \cle{presbytie} : avec l'âge, le cristallin perd son élasticité (sclérose) et les muscles ciliaires s'affaiblissent ; l'amplitude d'accommodation diminue, donc le PP recule (ici à $1{,}00$ m). Le PR restant à l'infini, la vision de loin est conservée : ce n'est ni une myopie ni une hypermétropie.
\item Les verres doivent donner du texte placé à $25$ cm ($\overline{OA} = -\SI{0,25}{m}$) une image virtuelle au PP actuel ($\overline{OA'} = -\SI{1,00}{m}$) :
\[ C_c = \frac{1}{-\num{1,00}} - \frac{1}{-\num{0,25}} = -1,00 + 4,00 = +3,00\ \delta. \]
Verres \textbf{convergents} de vergence $+3\ \delta$.
\item Avec ces verres, un objet à l'infini donne une image au foyer image des verres :
\[ \overline{OA'} = f' = \frac{1}{C_c} = \frac{1}{3} \approx +\SI{0,33}{m}, \]
soit une image \emph{réelle} située $33$ cm \emph{derrière} l'œil : les rayons incidents sur l'œil sont convergents, alors que l'œil au repos n'accepte que des rayons provenant de l'infini (ou divergents depuis un point situé entre le PP et l'infini en accommodant). L'image sur la rétine serait floue. La personne retire donc ses verres de lecture pour regarder au loin.
\end{enumerate}
\end{corrige}

\begin{corrige}[Exercice 9 — Consultation chez l'ophtalmologue de Yaoundé]
\begin{enumerate}
\item L'\cle{œil réduit} est un modèle simplifié de l'œil réel. Trois éléments essentiels : une \textbf{lentille mince convergente} de vergence variable (le cristallin), un \textbf{diaphragme} (l'iris) et un \textbf{écran fixe} (la rétine) placé à la distance $d$ de la lentille.
\item \textbf{Punctum remotum.} Au repos, l'image du PR se forme sur la rétine : $\overline{OA'} = d = \SI{0,015}{m}$. Relation de conjugaison :
\[ C_{\min} = \frac{1}{\overline{OA'}} - \frac{1}{\overline{OR}} \implies \frac{1}{\overline{OR}} = \frac{1}{\num{0,015}} - 60 = 66,67 - 60 = 6,67\ \text{m}^{-1}. \]
\[ \overline{OR} = \frac{1}{6,67} \approx +\SI{0,15}{m} = +\SI{15}{cm}. \]
Le signe positif signifie que le PR est \emph{virtuel}, situé derrière l'œil : l'œil doit accommoder pour voir l'infini.
\item \textbf{Punctum proximum.} En accommodation maximale, image toujours sur la rétine :
\[ \frac{1}{\overline{OP}} = \frac{1}{\num{0,015}} - C_{\max} = 66,67 - 64 = 2,67\ \text{m}^{-1} \implies \overline{OP} \approx +\SI{0,375}{m}. \]
Résultat positif : le PP serait virtuel, derrière l'œil ! Concrètement, cela signifie que même en accommodant au maximum, l'œil ne peut voir net aucun objet réel proche : sa zone de vision nette est très réduite (objets réels vus nets seulement en accommodant, entre l'infini et une distance limite). On peut aussi exprimer : pour voir un objet réel à la distance $x$, il faut $C = 66,67 + 1/x \leq 64$, impossible car $66,67 > 64$.
\item Amplitude d'accommodation :
\[ A = C_{\max} - C_{\min} = 64 - 60 = 4,0\ \delta. \]
\item L'élève est \textbf{hypermétrope} : son PR est virtuel (derrière l'œil), signe d'un œil trop peu convergent. Un œil emmétrope aurait son PR à l'infini ; un myope l'aurait réel, devant l'œil.
\item Vérifions : pour un objet à l'infini, l'image se formerait à
\[ f'_{\text{repos}} = \frac{1}{C_{\min}} = \frac{1}{60} \approx \SI{0,01667}{m} = \SI{16,7}{mm} > d = \SI{15}{mm}. \]
L'image se formerait donc \emph{derrière} la rétine, et non en avant. L'affirmation de l'ophtalmologue est \textbf{incompatible} avec les mesures : « image en avant de la rétine » correspond à la myopie. Il y a probablement une erreur de formulation ; le diagnostic correct est l'hypermétropie.
\end{enumerate}
\textbf{Barème indicatif (10 pts) :} définition et éléments (1,5) ; PR avec justification (2) ; PP (2) ; amplitude (1) ; diagnostic justifié (1,5) ; vérification et conclusion critique (2).\par
\textbf{Points de vigilance :} citer explicitement la relation de conjugaison et préciser que l'image se forme sur la rétine ($\overline{OA'} = d$) ; interpréter le \emph{signe} des grandeurs algébriques (PR virtuel $\Rightarrow$ hypermétropie) ; conserver les unités SI (mètres) dans les calculs de vergence.
\end{corrige}

\begin{corrige}[Exercice 10 — Le chauffeur de mototaxi de Douala]
\begin{enumerate}
\item M. Etoundi est \textbf{myope} : son PR est réel mais rapproché ($2{,}0$ m). Son œil est trop convergent (ou son globe trop long) : les rayons parallèles venant d'un panneau lointain convergent \emph{en avant} de la rétine, puis divergent vers elle, donnant une tache floue. (Schéma : œil réduit, rayons parallèles, foyer image $F'$ situé entre le cristallin et la rétine.)
\item La lentille accolée doit donner de l'objet à l'infini une image virtuelle au PR : $\overline{OA'} = -\SI{2,0}{m}$.
\[ C_1 = \frac{1}{\overline{OA'}} - \frac{1}{\overline{OA}} = \frac{1}{-2,0} - 0 = -0,50\ \delta. \]
Lentille \textbf{divergente} de vergence $-0,50\ \delta$ (myopie légère).
\item \begin{enumerate}
\item L'œil corrigé pour la vision de loin a son PP à $40$ cm. La correction supplémentaire doit donner du document à $25$ cm ($\overline{OA} = -\SI{0,25}{m}$) une image virtuelle à $40$ cm ($\overline{OA'} = -\SI{0,40}{m}$) :
\[ C_2 = \frac{1}{-\num{0,40}} - \frac{1}{-\num{0,25}} = -2,50 + 4,00 = +1,50\ \delta. \]
\item Vergence totale des verres de lecture (verges accolées, les vergences s'ajoutent) :
\[ C_{\text{total}} = C_1 + C_2 = -0,50 + 1,50 = +1,00\ \delta. \]
\end{enumerate}
\item Les verres progressifs (double foyer) évitent de changer de lunettes : la \textbf{partie supérieure} du verre, de vergence $-0,50\ \delta$, corrige la myopie pour la vision de loin (conduite, panneaux) ; la \textbf{partie inférieure}, de vergence $+1,00\ \delta$, permet la lecture de près (compteur, téléphone, documents). Le regard bascule naturellement d'une zone à l'autre.
\item Conséquences pratiques d'une myopie non corrigée chez un chauffeur :
\begin{itemize}
\item lecture tardive des panneaux et des indications : distances de réaction réduites, freinages tardifs ;
\item mauvaise appréciation des obstacles lointains (piétons, nids-de-poule, véhicules à l'arrêt), surtout de nuit où la pupille dilatée accentue le flou : risque accru d'accident.
\end{itemize}
\end{enumerate}
\textbf{Barème indicatif (10 pts) :} diagnostic + schéma (2) ; $C_1$ avec justification (2) ; $C_2$ (2) ; vergence totale (1) ; verres progressifs (1,5) ; sécurité routière (1,5).\par
\textbf{Points de vigilance :} pour une correction, toujours écrire que la lentille donne une image \emph{virtuelle} au PR (vision de loin) ou au PP (vision de près) de l'œil ; les vergences de lentilles accolées s'additionnent ; ne pas oublier le signe négatif des distances virtuelles.
\end{corrige}

\begin{corrige}[Exercice 11 — Étude complète d'un œil hypermétrope]
\begin{enumerate}
\item \textbf{Schéma.} Œil réduit : lentille (cristallin), rétine à $15$ mm. Un objet proche envoie des rayons fortement divergents ; le cristallin, trop peu convergent, les fait converger en un point situé \emph{derrière} la rétine (prolongements en pointillés). Sur la rétine : tache floue. Le défaut est l'\cle{hypermétropie} (œil trop peu convergent ou globe trop court).
\item En accommodation maximale, l'objet est au PP ($\overline{OA} = -\SI{0,75}{m}$) et l'image est sur la rétine ($\overline{OA'} = \SI{0,015}{m}$) :
\[ C_{\max} = \frac{1}{\num{0,015}} - \frac{1}{-\num{0,75}} = 66,67 + 1,33 = 68,0\ \delta. \]
\item Pour voir net à l'infini sans accommoder, il faudrait :
\[ C_0 = \frac{1}{d} = \frac{1}{\num{0,015}} \approx 66,7\ \delta. \]
Vergence réelle au repos : le PR est virtuel à $\overline{OR} = +\SI{0,50}{m}$ :
\[ C_{\text{repos}} = \frac{1}{\num{0,015}} - \frac{1}{+\num{0,50}} = 66,67 - 2,00 = 64,7\ \delta. \]
La vergence au repos est donc inférieure d'environ $2\ \delta$ à la vergence idéale : l'œil est trop peu convergent, ce qui confirme l'hypermétropie. L'élève compense en accommodant ($64,7 + 2 = 66,7\ \delta$ pour l'infini), d'où une fatigue visuelle.
\item La lentille de contact doit donner du texte à $25$ cm ($\overline{OA} = -\SI{0,25}{m}$) une image virtuelle au PP ($\overline{OA'} = -\SI{0,75}{m}$) :
\[ C_c = \frac{1}{-\num{0,75}} - \frac{1}{-\num{0,25}} = -1,33 + 4,00 = +2,67\ \delta. \]
Lentille \textbf{convergente} de vergence $+2,67\ \delta$ (en pratique, $+2,75\ \delta$).
\item \textbf{Analyse dimensionnelle.} La relation $C_c = \dfrac{1}{\overline{OA'}} - \dfrac{1}{\overline{OA}}$ ne comporte que des inverses de longueurs :
\[ \left[\frac{1}{\overline{OA'}}\right] = \left[\frac{1}{\overline{OA}}\right] = \text{m}^{-1} = \text{dioptrie}. \]
La somme de deux termes de même dimension est homogène, et le résultat est bien une vergence (en $\delta$) : la relation est homogène.
\item \textbf{Conclusion pour le carnet de santé.} « L'élève présente une hypermétropie : son œil, trop peu convergent (vergence au repos $64,7\ \delta$ au lieu de $66,7\ \delta$), l'oblige à accommoder en permanence, même pour la vision de loin, et son punctum proximum est reculé à $75$ cm. Une correction par lentilles convergentes de vergence $+2,67\ \delta$ est prescrite pour la lecture à $25$ cm. Un suivi ophtalmologique annuel est recommandé. »
\end{enumerate}
\textbf{Barème indicatif (12 pts) :} schéma + diagnostic (2) ; $C_{\max}$ (2) ; $C_0$ et comparaison (2,5) ; $C_c$ avec relation de conjugaison et signes (2,5) ; analyse dimensionnelle (1,5) ; conclusion rédigée (1,5).\par
\textbf{Points de vigilance :} distinguer clairement les deux utilisations de la relation de conjugaison (appliquée à l'œil : image sur la rétine ; appliquée à la lentille correctrice : image virtuelle au PP ou au PR) ; justifier la nature de la lentille par le \emph{signe} de $C_c$ ; l'analyse dimensionnelle exige d'écrire les dimensions, pas seulement les unités du résultat.
\end{corrige}

\courssub{Synthèse et méthode de résolution}

\begin{methode}[Résolution type d'un problème de correction oculaire]
\begin{enumerate}
\item \textbf{Identifier le défaut} à partir des données (PR, PP, vergence au repos/max) et le type de lentille correctrice (divergente pour myopie, convergente pour hypermétropie/presbytie).
\item \textbf{Dessiner l'œil réduit} (lentille, rétine à $d$) et placer le PR ou le PP concerné.
\item \textbf{Écrire la relation de conjugaison pour l'œil} (image sur la rétine : $\overline{OA'} = +d$) pour trouver $C_{\min}$, $C_{\max}$, PR ou PP si besoin.
\item \textbf{Écrire la relation de conjugaison pour la lentille correctrice} : elle doit former une \emph{image virtuelle} au PR (vision de loin) ou au PP (vision de près) de l'œil nu. $\overline{OA'}_L$ est donc négatif (virtuel, même côté que l'objet).
\item \textbf{Calculer $C_c$} avec les signes algébriques rigoureux (axe orienté dans le sens de la lumière incidente).
\item \textbf{Conclure} sur la nature de la lentille (signe de $C_c$) et donner la vergence en dioptries ($\delta$).
\item \textbf{Vérifier} par analyse dimensionnelle et cohérence physique (signe, ordre de grandeur).
\end{enumerate}
\end{methode}

\begin{savaistu}[Histoire et contexte camerounais]
\emph{Histoire.} Johannes Kepler (1604) comprit le premier que la rétine est l'écran récepteur et que le cristallin fait office de lentille. Hermann von Helmholtz (1855) inventa l'ophtalmoscope et mesura l'accommodation, validant le modèle de l'œil réduit.
\emph{Au Cameroun.} Selon l'OMS, les erreurs de réfraction non corrigées sont la première cause de déficience visuelle. Le Programme National de Santé Oculaire (PNSO) et des ONG comme \emph{Sightsavers} ou \emph{CBM} organisent des campagnes de dépistage dans les écoles (Yaoundé, Douala, zones rurales) et distribuent des lunettes. La chirurgie de la cataracte (remplacement du cristallin par un implant intraoculaire de vergence calculée) est réalisée dans les hôpitaux de référence (Hôpital Central, Hôpital Général, Laquintinie). La téléophtalmologie via MTN/Orange commence à connecter les centres de santé périphériques aux spécialistes des grandes villes.
\end{savaistu}

\begin{aretenir}[Bilan des corrections]
\begin{center}
\begin{tabularx}{\linewidth}{|l|X|X|X|}\hline
\rowcolor{popblueL}
\textbf{Défaut} & \textbf{Origine (modèle réduit)} & \textbf{Lentille correctrice} & \textbf{Principe} \\ \hline
\cle{Myopie} & Œil trop convergent ($C_{\min}$ trop grande) ou globe trop long ($d$ trop grand). PR réel, fini. & \textbf{Divergente} ($C_c < 0$) & Ramène l'image de l'infini au PR (image virtuelle au PR). \\ \hline
\cle{Hypermétropie} & Œil pas assez convergent ($C_{\min}$ trop faible) ou globe trop court ($d$ trop petit). PR virtuel. & \textbf{Convergente} ($C_c > 0$) & Donne de l'objet proche une image virtuelle au PP (ou de l'infini une image virtuelle au PR). \\ \hline
\cle{Presbytie} & Perte d'élasticité du cristallin $\Rightarrow$ $A$ diminue, PP recule. PR à l'infini. & \textbf{Convergente} ($C_c > 0$) pour la vision de près & Donne de l'objet à $25$ cm une image virtuelle au PP actuel. \\ \hline
\end{tabularx}
\end{center}
\textbf{Formule unique (lentille correctrice accolée) :}
\[ C_c = \frac{1}{\overline{OA'}_L} - \frac{1}{\overline{OA}_L} \]
avec $\overline{OA}_L$ = distance objet (négatif si réel devant la lentille) et $\overline{OA'}_L$ = distance image virtuelle au PR ou PP de l'œil (négatif).
\end{aretenir}

\begin{attention}[Pièges fréquents au Probatoire]
\begin{itemize}
\item \textbf{Signe des distances :} $\overline{OA} < 0$ pour un objet réel ; $\overline{OA'} > 0$ pour une image réelle (sur la rétine) ; $\overline{OA'} < 0$ pour une image virtuelle (donnée par la lentille correctrice au PR/PP).
\item \textbf{Unités :} Les distances \textbf{doivent} être en \textbf{mètres} (m) dans la formule $C = 1/\overline{OA'} - 1/\overline{OA}$ pour obtenir des dioptries ($\delta = \text{m}^{-1}$).
\item \textbf{Deux relations de conjugaison distinctes :} Une pour l'œil (inconnue = vergence $C$ ou position objet/image), une pour la lentille correctrice (inconnue = $C_c$). Ne pas les mélanger.
\item \textbf{Amplitude d'accommodation :} $A = C_{\max} - C_{\min}$ (toujours positif). Ne pas confondre avec $1/\text{PP}$ (valable seulement pour l'emmétrope).
\item \textbf{Verres progressifs :} La vergence totale pour la vision de près est la \textbf{somme algébrique} des vergences (loin + près) car les verres sont accolés.
\end{itemize}
\end{attention}

\newpage

\coursec{Fiche bilan}

\begin{popfigure}
\centering
\begin{tikzpicture}[
    mindmap,
    grow cyclic,
    every node/.style={concept, execute at begin node=\hskip0pt},
    concept color=popblue,
    level 1/.style={level distance=4.5cm, sibling angle=360/7},
    level 2/.style={level distance=3cm, sibling angle=45},
    popblue/.style={concept color=popblue},
    poporange/.style={concept color=poporange},
    popgreen/.style={concept color=popgreen},
    poppurple/.style={concept color=poppurple},
    poppink/.style={concept color=poppink},
    popgold/.style={concept color=popgold},
    popdark/.style={concept color=popdark}
]
\node[concept, text width=2.8cm, align=center, font=\bfseries\large] {L'\textbackslash{} Œil Réduit}
    child[popblue] { node[concept, text width=2.2cm, align=center] {Modèle\\ Dioptre unique + Rétine + Cristallin} }
    child[poporange] { node[concept, text width=2.2cm, align=center] {Accommodation\\ Variation $C_{\text{crist}}$\\ PP, PR, Amplitude} }
    child[popgreen] { node[concept, text width=2.2cm, align=center] {Emmétropie\\ PR = $\infty$\\ PP $\approx$ 25 cm} }
    child[poppurple] { node[concept, text width=2.2cm, align=center] {Myopie\\ PR fini $<$ $\infty$\\ Verre Divergent} }
    child[poppink] { node[concept, text width=2.2cm, align=center] {Hypermétropie\\ PR virtuel $>$ $\infty$\\ Verre Convergent} }
    child[popgold] { node[concept, text width=2.2cm, align=center] {Presbytie\\ $\Delta C_{\text{acc}} \downarrow$\\ Addition convergente} }
    child[popdark] { node[concept, text width=2.2cm, align=center] {Correction\\ $C_L = -\frac{1}{PR}$\\ $C_L = \frac{1}{d}-\frac{1}{PP}$} };
\end{tikzpicture}
\legende{Carte mentale de la leçon : de l'œil réduit aux corrections des amétropies.}
\end{popfigure}

\begin{bilan}
\courssub{Définitions et notions clés}
\begin{itemize}
    \item \textbf{Œil réduit} : Modèle simplifié de l'œil : un unique dioptre sphérique (cornée, rayon $R \approx 7,8\,\text{mm}$, indice $n'=1,336$), un plan image fixe (rétine, distance $d \approx 22,6\,\text{mm}$) et un système lenticulaire variable (cristallin) de convergence $C_c$.
    \item \textbf{Accommodation} : Capacité du cristallin à modifier sa convergence $C_c$ (donc la convergence totale $C_{\text{œil}} = C_{\text{cornée}} + C_c$) pour netter une source à distance variable.
    \item \textbf{Punctum Remotum (PR)} : Point le plus lointain vu net \emph{sans} accommodation ($C_c = C_{c,\min}$). Emmétrope : $PR = \infty$.
    \item \textbf{Punctum Proximum (PP)} : Point le plus proche vu net \emph{avec} accommodation maximale ($C_c = C_{c,\max}$).
    \item \textbf{Amplitude d'accommodation} : $\mathcal{A} = \frac{1}{PP} - \frac{1}{PR}$ (en m$^{-1}$ ou dioptries $\delta$).
    \item \textbf{Emmétropie} : Œil sans défaut ; image nette sur la rétine pour un objet à l'infini ($PR=\infty$).
    \item \textbf{Amétropie} : Défaut de réfraction (myopie, hypermétropie, astigmatisme, presbytie).
\end{itemize}

\courssub{Formules et lois à connaître par cœur}
\begin{center}
\begin{tabularx}{\linewidth}{|>{\bfseries}l|X|c|}\hline
\textbf{Notion} & \textbf{Formule / Relation} & \textbf{Unités / Remarques} \\ \hline
Convergence totale & $C_{\text{œil}} = C_{\text{cornée}} + C_c$ & $C_{\text{cornée}} \approx 43\,\delta$ (fixe) \\ \hline
Relation de Descartes (Gauss) & $C_{\text{œil}} = V + V'$ & $V = 1/(\overline{OA})$, $V' = 1/(\overline{OA'})$\\ 
(signe algébrique : réel $>0$) & $V' = 1/d_{\text{rétine}} \approx +44,2\,\delta$ & $d_{\text{rétine}} \approx 22,6\,\text{mm}$ \\ \hline
Myopie (verre au plan cornéen) & $C_L = -\dfrac{1}{PR}$ & $PR$ en mètres, $C_L < 0$ (divergent) \\ \hline
Hypermétropie / Presbytie & $C_L = \dfrac{1}{d_{\text{voulu}}} - \dfrac{1}{PP_{\text{réel}}}$ & $d_{\text{voulu}} = 0,25\,\text{m}$ (lecture), $C_L > 0$ \\ \hline
Puissance addition (presbytie) & $\text{Add} = \dfrac{1}{d_{\text{lecture}}} - \dfrac{1}{PP_{\text{réel}}}$ & Souvent $+0,75\,\delta$ à $+3,00\,\delta$ \\ \hline
\end{tabularx}
\end{center}

\courssub{Méthodes express (résolution type Probatoire)}
\begin{enumerate}
    \item \textbf{Analyser l'énoncé} : Identifier le défaut (myope ? hypermétrope ? presbyte ?), relever $PR$ ou $PP$ donnés, distance de lecture souhaitée ($25\,\text{cm}$ par défaut).
    \item \textbf{Modéliser} : Schématiser l'œil réduit + verre correcteur au plan cornéen. Préciser le signe de la vergence demandée ($V$ objet).
    \item \textbf{Choisir la formule} :
    \begin{itemize}
        \item Myopie : $C_L = -1/PR$ (objet à l'infini $\Rightarrow V=0$).
        \item Hypermétropie : $C_L = 1/PP_{\text{normal}} - 1/PP_{\text{réel}}$ (souvent $PP_{\text{normal}}=25\,\text{cm}$).
        \item Presbytie : Addition $= 1/d_{\text{lecture}} - 1/PP_{\text{réel}}$.
    \end{itemize}
    \item \textbf{Calculer} : Convertir distances en \textbf{mètres}. Appliquer la formule. Résultat en dioptries ($\delta$ ou m$^{-1}$).
    \item \textbf{Conclure} : Nature du verre (convergent/divergent), puissance commerciale (arrondi au $0,25\,\delta$ près), position du verre.
\end{enumerate}
\end{bilan}

\begin{aretenir}[Checklist avant le Probatoire]
\begin{itemize}
    \item $\square$ Je sais schématiser l'œil réduit (dioptre cornéen, plan rétinien, cristallin) et nommer $C_{\text{cornée}}$, $C_c$, $d_{\text{rétine}}$.
    \item $\square$ Je définis rigoureusement $PR$ (accommodation nulle) et $PP$ (accommodation max) et je donne leurs valeurs pour un emmétrope standard ($PR=\infty$, $PP=25\,\text{cm}$).
    \item $\square$ J'explique le mécanisme d'accommodation : variation de $C_c$ par déformation du cristallin (zonule de Zinn / muscle ciliaire).
    \item $\square$ Je distingue les trois amétropies principales : \textbf{Myopie} ($PR$ fini, œil trop convergent/trop long), \textbf{Hypermétropie} ($PR$ virtuel, œil pas assez convergent/trop court), \textbf{Presbytie} (perte d'amplitude $\mathcal{A}$ avec l'âge, $PP$ s'éloigne).
    \item $\square$ J'associe le bon signe de verre correcteur : \textbf{Divergent} ($C_L<0$) pour myopie, \textbf{Convergent} ($C_L>0$) pour hypermétropie et presbytie.
    \item $\square$ J'applique sans erreur la formule de correction myopie : $C_L = -1/PR$ (avec $PR$ en \textbf{mètres}).
    \item $\square$ J'applique la formule de correction vision de près (hypermétropie/presbytie) : $C_L = 1/0,25 - 1/PP_{\text{réel}}$.
    \item $\square$ Je vérifie l'homogénéité de mes formules : $[C] = \text{m}^{-1}$, $[1/d] = \text{m}^{-1}$.
    \item $\square$ Je précise toujours que le verre correcteur est placé au \textbf{plan cornéen} (hypothèse du cours) sauf mention contraire.
    \item $\square$ Je sais lire une ordonnance ophtalmologique (sphère, cylindre, axe, addition) et lier "Sphère" à $C_L$ calculé.
    \item $\square$ Je traite un problème complet : données $\rightarrow$ modèle $\rightarrow$ calcul $\rightarrow$ réponse finale avec unité et phrase de conclusion.
\end{itemize}
\end{aretenir}

\findecours

\end{document}
